% !TeX root = ../../Book1.tex
% 纲要标题：## 第12章　向量丛、张量与整体构造
% 纲要位置：几何学计划纲要.md，第 568 行。
% 正文状态：初稿；尚未正式审核。

% 纲要内容开始
% 
% > 本章目标与证明范围：由切丛推广到一般向量丛，证明局部表示和整体对象之间的对应。多线性代数围绕几何运算重建，单位分解承担明确的构造任务。
% 
% 纲要内容结束

\chapter{向量丛、张量与整体构造}
\label{g1:ch:12}
\BookChapterNumber{g1:ch:12}

\begin{chapterabstract}
	切丛的各个纤维都是向量空间，但在不同坐标图中需要用随点变化的可逆矩阵互相识别。本章由这一结构定义一般向量丛，证明余循环粘合、常秩同态的核与像以及拉回、商丛的构造。随后以泛性质建立张量积，以行列式配对固定外幂和楔积的系数，区分体积形式与密度。在整体构造方面，完整证明光滑单位分解、闭集上的截面延拓、丛度量的存在与短正合丛列的分裂，最后把这些工具用于张量场、自然拉回和几类结构群约化。
\end{chapterabstract}

\begin{learninggoals}
	\begin{itemize}
		\item 从局部标架写出过渡矩阵及余循环条件，说明它们如何决定整体向量丛，并辨认局部平凡与整体平凡的区别。
		\item 通过泛性质与基展开解释张量、对称幂、外幂和收缩，计算楔积的系数、分次符号与换标架公式。
		\item 构造支集局部有限的光滑单位分解，明确闭集延拓及其相对版本所需的相容条件。
		\item 用单位分解构造丛度量，再用正交补分裂短正合丛列，同时辨认这些构造中依赖选择的部分。
		\item 把局部分量解释为光滑张量场，区分拉回丛与源流形上的张量，并用适配标架描述度量、定向及其他逐点几何结构。
	\end{itemize}
\end{learninggoals}

上一章的切空间在每张坐标图中都有一组基，但从一张图走进另一张图时，这组基会改变。过渡矩阵把一组基中的分量转换到另一组基中；这些局部识别的整体相容方式决定了能否选出一组全局基。Möbius 线丛将展示最小的例子：每个小圆弧上的纤维都是一条普通实直线，绕圆周一周后的符号变化，却阻止了处处非零的整体截面。局部模型相同，并不足以决定整体对象平凡。

\ProseParagraphBreak

另一些对象又确实总能从局部构造出来。例如，在每张局部标架中都能声明基向量正交归一；虽然这些局部内积未必一致，却可以用适当权重相加，得到整体丛度量。为什么同样的操作不能给 Möbius 线丛制造非零截面？关键不只在于权重是否光滑，还在于加权相加能否保持所需条件。正定性可以保持，非零性、非退化性和某些代数恒等式却未必可以。

\ProseParagraphBreak

本章沿用\chapref{g1:ch:11}的光滑流形、切丛和微分定义，默认底流形有限维、Hausdorff 且第二可数；正文讨论的实光滑向量丛均为有限秩，相关构造允许底流形带边界。读者只需熟悉有限维线性代数中的基、对偶、商空间和行列式。张量积的泛性质、对称与交替运算在这里由多线性映射重建。单位分解所需的紧穷竭与局部有限细化，也会恢复实际使用的构造。

\ProseParagraphBreak

阅读时可以把每一项构造分为两步：先检查同一纤维上的线性代数，再检查它随底点是否光滑，以及换标架后是否代表同一个对象。Lee 第二版的第10章可用于对照向量丛，第12--13章讨论张量与 Riemann 度量，第2章的 Partitions of Unity 一节对应光滑权重的构造；第14--15章及第16章的 Densities 一节可补充交替形式、定向与密度的阅读\cite[第40--47、249--271、304--399、427--433页]{Lee2012Smooth}。本章只建立形式的逐点代数及光滑性，外微分与积分留在\chapref{g1:ch:14}。

% !TeX root = ../../Book1.tex
% 纲要标题：### 12.1 向量丛及其粘合
% 纲要位置：几何学计划纲要.md，第 572 行。
% 正文状态：初稿；正文、证明、例子与配套习题解答已撰写，尚未正式审核。

% 纲要内容开始
% 
% * 局部平凡化、过渡函数和余循环
% * 丛同态、子丛及常秩条件
% * 拉回丛、直和与商丛
% * 标架丛与全局标架
% 
% 纲要内容结束

\section{向量丛及其粘合}
\label{g1:sec:1201}

切丛把每点的切空间组织成一个光滑对象。它的局部坐标同时记住基点与向量分量，而换坐标时，分量按随基点光滑变化的可逆矩阵变换。许多几何对象也有这种形式：曲面每点的法向量、线性方程在每点的解空间，以及随后要讨论的张量，都不是固定向量空间中的一组数，而是一族随基点变化的向量空间。向量丛保留了切丛构造中的这一共同部分。

\ProseParagraphBreak

本节讨论有限秩实光滑向量丛，底空间为 Hausdorff、第二可数的光滑流形，允许有边界；边界附近的光滑性沿用局部延拓到 Euclid 开集的定义。我们先弄清局部数据怎样粘合，再讨论如何逐纤维取子空间、商空间，以及如何更换底空间。关键问题始终是：逐点成立的线性代数，何时能随基点光滑地进行。

\subsection{局部平凡化、过渡函数和余循环}
\label{g1:subsec:1201:01}

\cref{g1:thm:tangent-bundle-atlas}已经证明，切丛在坐标域 \(U\) 上可写成 \(U\times\mathbb R^n\)，并验证了总空间的 Hausdorff 性与第二可数性。这种乘积表示只在基点的一个邻域上给定；换图时可以出现非恒定矩阵，因此它还没有给出整个切丛与 \(M\times\mathbb R^n\) 的同构。

\begin{definition}
	\label{g1:def:vector-bundle}
	设 \(M\) 为 Hausdorff、第二可数的光滑 \(n\) 维流形，可带边界，\(r\geq0\) 为整数。\(M\) 上的秩 \(r\) 的实\term[xiangliangcong]{向量丛}[vector bundle]由光滑流形 \(E\)、光滑满射 \(\pi:E\to M\) 及每个集合 \(E_p=\pi^{-1}(p)\) 上的 \(r\) 维实向量空间结构组成，并要求每个 \(p\in M\) 都有开邻域 \(U\) 和微分同胚
	\[
		\phi_U:\pi^{-1}(U)\longrightarrow U\times\mathbb R^r,
		\qquad \phi_U(v_p)=(p,[v_p]_U),
	\]
	使 \([\cdot]_U:E_q\to\mathbb R^r\) 对每个 \(q\in U\) 都是线性同构。这样的 \(\phi_U\) 称为\term[jubupingfanhua]{局部平凡化}[local trivialization]。\(E\)、\(M\)、\(E_p\) 分别称为总空间、底空间和 \(p\) 上的\term[xianwei]{纤维}[fiber]；\(E|_U=\pi^{-1}(U)\) 表示限制丛。
\end{definition}

这里总空间同样要求 Hausdorff、第二可数。局部模型说明 \(\dim E=n+r\)；若底空间有边界，模型为 \(\mathbb H^n\times\mathbb R^r\)，非负坐标来自底空间，所以 \(\partial E=\pi^{-1}(\partial M)\)。秩零的丛每个纤维只有零向量，其总空间可由 \(\pi\) 识别为 \(M\)。本文固定一个秩；若底空间不连通，也可在不同连通分支分别规定秩。

\begin{definition}
	\label{g1:def:bundle-section}
	设 \(\pi:E\to M\) 为有限秩实光滑向量丛，\(U\subset M\) 开。满足 \(\pi\circ s=\operatorname{id}_U\) 的光滑映射 \(s:U\to E\) 称为 \(E\) 在 \(U\) 上的\term[jiemian]{截面}[section]，全体光滑截面记为 \(\Gamma(U,E)\)，并简记 \(\Gamma(E)=\Gamma(M,E)\)。若秩为 \(r\)，一组截面 \(e_1,\ldots,e_r\) 在每个 \(p\in U\) 都组成 \(E_p\) 的有序基，就称为 \(U\) 上的\term[jububiaojia]{局部标架}[local frame]。
\end{definition}

局部平凡化把标准基向量送回纤维，便给出局部标架。反过来，给定局部标架 \(e=(e_1,\ldots,e_r)\)，记
\[
	e(p)v=\sum_{a=1}^r v^a e_a(p),\qquad v\in\mathbb R^r.
\]
在任一已有平凡化中，\(e_a\) 的分量组成一个光滑可逆矩阵 \(A(p)\)，因而 \((p,v)\mapsto e(p)v\) 及其逆都光滑。逆的光滑性来自 \(A^{-1}=\operatorname{adj}(A)/\det A\)，这里行列式处处非零。于是局部标架与局部平凡化确实等价，而不仅是逐点选基。

\ProseParagraphBreak

在开覆盖 \(\{U_i\}\) 上选标架 \(e_i\)，其中下标 \(i\) 标记开集，标架内部的向量仍用 \(e_{i,a}\) 表示。在交集上存在唯一光滑函数
\[
	g_{ij}:U_i\cap U_j\longrightarrow\operatorname{GL}(r,\mathbb R)
\]
使
\begin{equation}
	\label{g1:eq:bundle-frame-transition}
	e_j=e_i g_{ij},\qquad v_i=g_{ij}v_j.
\end{equation}
第二式中 \(v_i,v_j\) 是同一个纤维向量在两组标架下的列分量，因为 \(e_i v_i=e_jv_j\)。因此标架变换与分量变换在这个下标约定下都出现 \(g_{ij}\)，但前一式把第 \(j\) 组基写成第 \(i\) 组基的线性组合，后一式则把第 \(j\) 组分量转换成第 \(i\) 组分量。

\ProseParagraphBreak

这些\term[guoduhanshu]{过渡函数}[transition function]满足
\begin{equation}
	\label{g1:eq:bundle-cocycle}
	g_{ii}=I_r,\qquad g_{ji}=g_{ij}^{-1},\qquad
	g_{ij}g_{jk}=g_{ik}.
\end{equation}
最后一个等式在三重交集上成立，称为\term[yuxunhuantiaojian]{余循环条件}[cocycle condition]：先由第 \(k\) 组分量换到第 \(j\) 组，再换到第 \(i\) 组，必须与直接转换相同。它由 \(e_k=e_jg_{jk}=e_i g_{ij}g_{jk}\) 及基的线性无关性得到。对截面 \(s\)，局部表达 \(s=e_i s_i\) 中的列函数 \(s_i:U_i\to\mathbb R^r\) 光滑，且满足 \(s_i=g_{ij}s_j\)；反之，满足这些关系的光滑列函数唯一确定一个光滑截面。纤维中的加法与数乘由此定义了截面的逐点运算，并使 \(\Gamma(E)\) 成为 \(C^\infty(M)\) 上的模。

\begin{example}
	\label{g1:exm:tangent-bundle-transition-functions}
	对光滑流形 \(M\)，在坐标图 \(x_i:U_i\to x_i(U_i)\) 上取切标架 \(e_i=(\partial_{x_i^1},\ldots,\partial_{x_i^n})\)。由\eqref{g1:eq:tangent-bundle-transition}，
	\[
		g_{ij}(p)=\mathrm{D}(x_i\circ x_j^{-1})(x_j(p)).
	\]
	链式法则正好给出 \(g_{ij}g_{jk}=g_{ik}\)，所以切丛的过渡函数是坐标变换的 Jacobian 矩阵。一般向量丛的矩阵只须光滑、可逆并满足余循环条件，无须来自某个底空间坐标变换的导数。乘积丛 \(M\times\mathbb R^r\) 则有处处相同的标准标架，过渡函数都可取为 \(I_r\)。
\end{example}

余循环条件足以把一批局部乘积真正粘成向量丛。证明中既要保证同一个向量不会因走不同的换图路径而得到不同代表，也要保证粘合后的空间满足流形的整体拓扑条件。

\begin{theorem}[向量丛的粘合]
	\label{g1:thm:bundle-cocycle-gluing}
	设 \(M\) 为 Hausdorff、第二可数的光滑流形，可带边界；\(\{U_i\}_{i\in I}\) 为开覆盖，\(r\geq0\) 为固定整数。若光滑映射 \(g_{ij}:U_i\cap U_j\to\operatorname{GL}(r,\mathbb R)\) 在相应的二重、三重交集上满足 \(g_{ii}=I_r\)、\(g_{ji}=g_{ij}^{-1}\)、\(g_{ij}g_{jk}=g_{ik}\)，即\eqref{g1:eq:bundle-cocycle}，则它们确定一个秩 \(r\) 的实光滑向量丛，其局部平凡化按 \((p,v_j)\mapsto(p,g_{ij}(p)v_j)\) 转换。该丛在保持这些局部识别的丛同构意义下唯一。
\end{theorem}
\begin{proof}
	先把各局部乘积视为互不相交的副本，令
	\[
		\mathcal E=\coprod_{i\in I}(U_i\times\mathbb R^r),\qquad
		(i,p,v)\sim(j,q,w)
		\quad\Longleftrightarrow\quad
		p=q\ \text{且}\ v=g_{ij}(p)w.
	\]
	自反性来自 \(g_{ii}=I_r\)，对称性来自 \(g_{ji}=g_{ij}^{-1}\)。若 \(v=g_{ij}(p)w\)、\(w=g_{jk}(p)z\)，则余循环条件给出 \(v=g_{ik}(p)z\)，所以关系也传递。令 \(E=\mathcal E/\!\sim\)，赋予商拓扑，并记商映射为 \(Q\)。由于等价关系不改变基点，\(\pi([i,p,v])=p\) 良定义且连续。

	商映射 \(Q\) 是开映射。事实上，对 \(\mathcal E\) 中的开集 \(O\)，它的饱和集在第 \(j\) 个副本中的部分为
	\[
		\bigcup_{i\in I}\left\{(p,w):p\in U_i\cap U_j,\ (p,g_{ij}(p)w)\in O_i\right\},
	\]
	其中 \(O_i=O\cap(U_i\times\mathbb R^r)\)。每项都是光滑同胚 \((p,w)\mapsto(p,g_{ij}(p)w)\) 下的开原像，故此并为开。商拓扑的定义于是说明 \(Q(O)\) 开。

	对每个 \(i\)，映射
	\[
		\tau_i:U_i\times\mathbb R^r\longrightarrow\pi^{-1}(U_i),
		\qquad\tau_i(p,v)=[i,p,v]
	\]
	是双射：每个基点在 \(U_i\) 的等价类都可换到第 \(i\) 个副本，而在同一副本中相等的两个代表必有相同分量。它连续且开，故为同胚；其像 \(\pi^{-1}(U_i)\) 也是 \(E\) 中的开集。在交集上有
	\[
		\tau_i^{-1}\tau_j(p,w)=(p,g_{ij}(p)w).
	\]
	再将每个 \(U_i\) 限制到底空间坐标域，便得到 \(E\) 的光滑坐标图；上述公式以及底空间换图的光滑性保证它们相容。在边界附近，矩阵元有局部光滑延拓，乘法公式同样给出半空间模型之间的光滑映射。

	还需核验 Hausdorff 性。两个点若有不同基点，用底空间中不交开邻域的逆像将它们分离；若基点相同，就把两点放进同一个 \(\tau_i\) 的像内，利用 \(U_i\times\mathbb R^r\) 的 Hausdorff 性分离。该像是 \(E\) 的开集，所以所得邻域也是总空间的开邻域。

	第二可数性不能由任意多个局部图直接推出，不过底空间的第二可数性提供可数子覆盖。具体地，对底空间可数基中每个包含在某个 \(U_i\) 内的基元素，选择一个包含它的 \(U_i\)；这些被选开的集合构成可数覆盖。每个对应乘积 \(U_i\times\mathbb R^r\) 有可数基，把它们通过 \(\tau_i\) 送到 \(E\) 再取并，就得到总空间的可数基。

	在纤维 \(E_p\) 中，选任一个 \(p\in U_i\)，用 \(\tau_i(p,\cdot)\) 运输 \(\mathbb R^r\) 的向量空间运算。换到另一副本时使用的是线性同构 \(g_{ij}(p)\)，故运算与所选副本无关。投影在局部是 \((p,v)\mapsto p\)，光滑且满射，\(\tau_i^{-1}\) 则是所需局部平凡化。

	最后，若另一个丛 \(E'\) 有同样局部识别 \(\tau_i'\)，则 \([i,p,v]\mapsto\tau_i'(p,v)\) 由相同过渡关系而良定义。它在每个局部乘积上为恒等坐标映射，故是光滑丛同构。任何保持局部识别的映射都必须使用这个公式，因此唯一。秩零时取 \(\operatorname{GL}(0,\mathbb R)=\{I_0\}\)，所有论证仍成立。
\end{proof}

这里的唯一性保留了指定的局部识别；改变标架当然会改变矩阵。若在 \(U_i\) 上以 \(\widetilde v_i=h_i v_i\) 改换分量，其中 \(h_i:U_i\to\operatorname{GL}(r,\mathbb R)\) 光滑，则
\begin{equation}
	\label{g1:eq:bundle-cocycle-change}
	\widetilde e_i=e_i h_i^{-1},\qquad
	\widetilde g_{ij}=h_i g_{ij}h_j^{-1}.
\end{equation}
后一个等式由 \(\widetilde v_i=h_i g_{ij}h_j^{-1}\widetilde v_j\) 得出。反过来，两组余循环若满足此式，\([i,p,v]\mapsto[i,p,h_i(p)v]\) 就给出所粘丛之间的同构：把 \(v=g_{ij}w\) 代入，得到 \(h_i v=\widetilde g_{ij}h_jw\)，恰好保证映射与代表选择无关。不同覆盖上的表示可先限制到两覆盖的交集，再作同一比较。因此过渡矩阵是丛在选定标架中的数据，并不是无需改换的绝对记号。

\begin{example}[Möbius 线丛]
	\label{g1:exm:mobius-line-bundle}
	令 \(M=S^1=\{e^{\mathrm{i}t}:t\in\mathbb R\}\)，把 \(\mathbb R\times\mathbb R\) 按
	\[
		(t+2\pi,v)\sim(t,-v)
	\]
	及其迭代关系取商，记所得空间为 \(L\)，投影为 \([t,v]\mapsto e^{\mathrm{i}t}\)。每个足够短的圆弧 \(U_i\) 上可选光滑角度提升 \(t_i:U_i\to\mathbb R\)，从而有局部表示 \((p,v)\mapsto[t_i(p),v]\)。若在交叠上 \(t_j-t_i=2\pi k_{ij}\)，则 \(k_{ij}\) 为局部常值整数，且
	\[
		[t_j(p),v]=[t_i(p),(-1)^{k_{ij}(p)}v].
	\]
	所以过渡函数为 \(g_{ij}=(-1)^{k_{ij}}\)，满足余循环条件。\cref{g1:thm:bundle-cocycle-gluing}给出秩一光滑丛结构；这些局部表示也是所写商空间的图，因为任一长度小于 \(2\pi\) 的提升圆弧及其 \(2\pi\) 平移副本，经商映射的饱和集仍开。这个丛称为 Möbius 线丛。

	它的光滑截面恰对应满足
	\begin{equation}
		\label{g1:eq:mobius-section-antiperiodic}
		f(t+2\pi)=-f(t)
	\end{equation}
	的光滑函数 \(f:\mathbb R\to\mathbb R\)，截面公式为 \(s(e^{\mathrm{i}t})=[t,f(t)]\)。因为 \([t+2\pi,f(t+2\pi)]=[t,-f(t+2\pi)]\)，此条件正是代表选择无关性；光滑性可在各提升圆弧上核验。若截面处处非零，\(f(0)\) 与 \(f(2\pi)=-f(0)\) 异号，中间值定理迫使某处 \(f\) 为零，矛盾。故 \(L\) 没有处处非零的光滑截面，甚至没有这样的连续截面。乘积线丛却有常值分量为 \(1\) 的截面，任何丛同构都会保持非零性，因此 \(L\) 与 \(S^1\times\mathbb R\) 不同构。
\end{example}

\begin{figure}[htbp]
	\centering
	% !TeX root = ../Book1.tex
% Möbius线丛的端点粘合；纤维只截取有限区间用于示意。
\begin{tikzpicture}[x=0.95cm,y=1.15cm,>=stealth]
	\fill[blue!3] (0,-1.35) rectangle (6.283,1.35);
	\draw[gray!65,dashed] (0,1.35)--(6.283,1.35) (0,-1.35)--(6.283,-1.35);
	\draw[very thick,->] (0,-1.35)--(0,1.35);
	\draw[very thick,->] (6.283,1.35)--(6.283,-1.35);
	\draw[gray!65,->] (-0.3,0)--(6.8,0) node[right] {\(t\)};
	\node[above] at (0,1.4) {\(v\)};
	\node[below left] at (0,0) {\(0\)};
	\node[below] at (3.142,-0.04) {\(\pi\)};
	\node[below right] at (6.283,0) {\(2\pi\)};
	\draw[red!75!black,dashed,thick] (0,0.55)--(6.283,0.55);
	\draw[red!75!black,fill=white] (0,0.55) circle (2pt) (6.283,0.55) circle (2pt);
	\node[red!75!black,above] at (4.9,0.6) {恒正分量不相合};
	\draw[blue!75!black,very thick,domain=0:6.283,samples=80] plot (\x,{cos(deg(\x/2))});
	\fill[blue!75!black] (0,1) circle (2pt) (3.142,0) circle (2pt) (6.283,-1) circle (2pt);
	\node[blue!75!black,above right] at (0.5,1.02) {\(f(t)=\cos(t/2)\)};
	\node[left] at (0,1) {\(1\)};
	\node[right] at (6.283,-1) {\(-1\)};
	\node[below,align=center] at (3.142,-1.52) {\((2\pi,v)\sim(0,-v)\)：两侧沿相反方向粘合};
\end{tikzpicture}

	\caption[Möbius 线丛中的截面与端点粘合]{矩形只画出线丛纤维的有限一段。左右竖边按相反方向粘合，\((2\pi,v)\) 对应 \((0,-v)\)。蓝线为截面分量 \(f(t)=\cos(t/2)\)，它的两端经过粘合相合，而在中间穿过零。红色虚线的分量恒为正，左右端点却不能按粘合关系相合，因而不是圆周上的截面。}
	\label{g1:fig:mobius-section}
\end{figure}

这说明局部平凡并未排除全局扭转。扭转并不是某个纤维内部的缺陷：每条纤维都是一维实向量空间；障碍出现在沿圆周比较各条纤维的基时。

\subsection{丛同态、子丛及常秩条件}
\label{g1:subsec:1201:02}

对同一个底空间上的两族向量空间，可以逐点施加线性映射。不过逐点线性只保证每个纤维内的行为，并不保证映射随基点连续或光滑。局部标架把后一个要求化为矩阵元的光滑性。

\begin{definition}
	\label{g1:def:bundle-morphism}
	设 \(E\to M\)、\(F\to M\) 为有限秩实光滑向量丛。覆盖恒等映射的\term[congtongtai]{丛同态}[bundle morphism]是光滑映射 \(\Phi:E\to F\)，使 \(\pi_F\Phi=\pi_E\)，且每个限制 \(\Phi_p:E_p\to F_p\) 都线性。如果它有同样性质的光滑逆，称为\term[congtonggou]{丛同构}[bundle isomorphism]。若 \(E\to N\)、\(F\to M\) 的底空间不同，覆盖光滑映射 \(f:N\to M\) 的丛同态则要求 \(\pi_F\Phi=f\pi_E\)，并要求 \(E_q\to F_{f(q)}\) 线性。
\end{definition}

\begin{proposition}[丛同态的矩阵]
	\label{g1:prop:bundle-morphism-matrices}
	设 \(E,F\) 为 Hausdorff、第二可数光滑流形 \(M\) 上的秩 \(r,s\) 的实光滑向量丛，允许 \(M\) 有边界。在共同开覆盖上取局部标架，其过渡函数为 \(g_{ij}^E,g_{ij}^F\)。一个覆盖恒等映射的逐纤维线性映射 \(\Phi:E\to F\) 为光滑丛同态，当且仅当它的局部矩阵 \(A_i:U_i\to\operatorname{Mat}_{s\times r}(\mathbb R)\) 光滑；这些矩阵满足
	\begin{equation}
		\label{g1:eq:bundle-morphism-transition}
		A_i g_{ij}^E=g_{ij}^F A_j.
	\end{equation}
	反之，任何满足此关系的光滑矩阵族都确定唯一丛同态。特别地，逐纤维为线性同构的光滑丛同态，其逆自动光滑。
\end{proposition}
\begin{proof}
	局部矩阵由 \(\Phi(e_i v)=f_i A_i v\) 定义。若 \(\Phi\) 光滑，则其作用于每个光滑基截面 \(e_{i,a}\) 后的分量光滑，故 \(A_i\) 各列光滑。反过来，光滑矩阵给出的局部映射 \((p,v)\mapsto(p,A_i(p)v)\) 光滑。交叠上，先换源分量再映射得到 \(A_i g_{ij}^E v_j\)，先映射再换目标分量得到 \(g_{ij}^F A_jv_j\)，二者对所有 \(v_j\) 相同恰好等价于\eqref{g1:eq:bundle-morphism-transition}。因此相容矩阵族唯一粘成光滑映射。最后，若每个 \(A_i(p)\) 都可逆，其逆矩阵光滑且局部逆彼此相容，从而构成光滑丛逆。
\end{proof}

例如映射的微分 \(\mathrm{T}F:\mathrm{T}M\to \mathrm{T}N\) 覆盖 \(F:M\to N\)，局部矩阵为坐标表达的导数，见\cref{g1:prop:differential-chain-rule}。底空间不同的情形只须在目标变换矩阵中代入 \(F(p)\)，矩阵判据及其证明不变。另一个常用例子是同底丛同态作用于截面：\(s\mapsto\Phi\circ s\) 把 \(\Gamma(E)\) 送入 \(\Gamma(F)\)，并与光滑函数的逐点乘法相容。

\begin{definition}
	\label{g1:def:vector-subbundle}
	设 \(E\to M\) 为秩 \(r\) 的实光滑向量丛，\(0\leq k\leq r\)。若每个 \(p\) 给定一个 \(k\) 维线性子空间 \(F_p\subset E_p\)，且每点附近存在 \(k\) 个光滑截面 \(s_1,\ldots,s_k\)，在每个邻近点都线性无关并张成相应的 \(F_p\)，则称 \(F=\coprod_p F_p\subset E\) 为秩 \(k\) 的光滑\term[zicong]{子丛}[subbundle]。
\end{definition}

定义中的光滑局部张成条件不可删去。例如在 \(\mathbb R\times\mathbb R^2\) 中，令 \(F_t=\mathbb R(1,0)\) 当 \(t\leq0\)，令 \(F_t=\mathbb R(0,1)\) 当 \(t>0\)。各个子空间的维数都为一，但若零附近有一个处处非零的连续张成截面，它在零点的值必须既落在横轴又落在纵轴，因而只能为零，矛盾。恒定维数没有替代连续变化，更没有替代光滑变化。

\begin{proposition}[子丛的适配标架]
	\label{g1:prop:subbundle-adapted-frame}
	设 \(E\to M\) 是 Hausdorff、第二可数光滑流形上的秩 \(r\) 实光滑向量丛，允许有边界。若 \(s_1,\ldots,s_k\) 在点 \(p\) 处线性无关，则在 \(p\) 的某个邻域上可将这些光滑截面补成 \(E\) 的局部标架。特别地，秩 \(k\) 子丛 \(F\subset E\) 带有唯一的光滑向量丛结构，使其局部标架给出的坐标为光滑平凡化；其拓扑是 \(E\) 的子空间拓扑，并且存在局部标架使
	\[
		E|_U\simeq U\times\mathbb R^r,\qquad
		F|_U\simeq U\times(\mathbb R^k\times\{0\}).
	\]
\end{proposition}
\begin{proof}
	先在 \(p\) 附近固定一个 \(E\) 的平凡化，把 \(s_a\) 写成光滑列向量。在 \(p\) 处用常向量 \(b_{k+1},\ldots,b_r\) 将其补为一组基，并把这些常向量在该平凡化下视为光滑截面。所成 \(r\times r\) 矩阵的行列式在 \(p\) 非零，由连续性，在某邻域仍非零。因此这些截面在整个邻域组成标架。\(k=0\) 时直接选任意标架，\(k=r\) 时无须补充。

	对子丛应用这个结论，前 \(k\) 个坐标向量正好张成 \(F\)。因此环境平凡化限制到 \(F\) 后为到 \(U\times\mathbb R^k\) 的同胚，局部拓扑正是子空间拓扑。两组这种标架的变换矩阵把 \(\mathbb R^k\times\{0\}\) 映到自身，其左上 \(k\times k\) 块为光滑可逆矩阵，故这些限制图光滑相容。Hausdorff 性与第二可数性由子空间继承，局部图给出维数 \(\dim M+k\) 及向量丛结构。唯一性由规定的局部图立即得到。边界点处仅用了光滑矩阵及非零行列式，所得局部模型仍为 \(\mathbb H^n\times\mathbb R^k\)。
\end{proof}

现在考虑一个丛同态的核与像。每条纤维内它们都是线性子空间；要成为子丛，维数先要在附近不变，而常秩条件还会给出光滑的局部基。\cref{g1:thm:constant-rank}中用过“选非零子式，再消去其余行”的办法；此处只需要对一族矩阵作消元，不需要对底空间改换常秩标准坐标。

\begin{theorem}[常秩丛同态]
	\label{g1:thm:constant-rank-bundle-morphism}
	设 \(E,F\) 为 Hausdorff、第二可数光滑流形 \(M\) 上的秩 \(r,s\) 实光滑向量丛，允许 \(M\) 有边界；\(\Phi:E\to F\) 为覆盖恒等映射的光滑丛同态。若 \(\operatorname{rank}\Phi_p=k\) 对所有 \(p\in M\) 恒成立，则
	\[
		\ker\Phi=\coprod_{p\in M}\ker\Phi_p,\qquad
		\operatorname{im}\Phi=\coprod_{p\in M}\operatorname{im}\Phi_p
	\]
	分别是 \(E,F\) 的秩 \(r-k,k\) 的光滑子丛。若秩只要求局部恒定，则此结论按各秩恒定的开集分别成立。
\end{theorem}
\begin{proof}
	固定一点并取局部标架，把 \(\Phi\) 表示为光滑 \(s\times r\) 矩阵 \(A(p)\)。先设 \(k>0\)。在基点选择一个非零 \(k\) 阶子式，重排源与目标的基后，将对应子块放到左上角。缩小邻域使该块一直可逆，写成
	\[
		A=\begin{pmatrix}B&C\\D&H\end{pmatrix},\qquad B\in\operatorname{GL}(k,\mathbb R).
	\]
	左乘可逆矩阵作消元，得到
	\[
		\begin{pmatrix}I&0\\-DB^{-1}&I\end{pmatrix}A
		=\begin{pmatrix}B&C\\0&H-DB^{-1}C\end{pmatrix}.
	\]
	再把右上块通过列消元变为零，可见右边的秩为 \(k+\operatorname{rank}(H-DB^{-1}C)\)。常秩条件遂迫使 \(H=DB^{-1}C\) 在整个邻域上成立。

	于是 \(A(x,y)=0\) 等价于 \(Bx+Cy=0\)，亦即 \(x=-B^{-1}Cy\)。所以矩阵
	\[
		\begin{pmatrix}-B^{-1}C\\I_{r-k}\end{pmatrix}
	\]
	的 \(r-k\) 个列向量给出核的光滑局部基。从下方的单位矩阵可见，这些列向量线性无关。另一方面，\(A\) 的前 \(k\) 列组成 \(\begin{psmallmatrix}B\\D\end{psmallmatrix}\)，因 \(B\) 可逆而线性无关；其余列为这些列乘以 \(B^{-1}C\)，故前 \(k\) 列正好给出像的光滑局部基。两族列向量都在先前选定标架下表示光滑截面，因而核与像满足子丛定义。

	若 \(k=0\)，矩阵处处为零，核为 \(E\)，像为零子丛。若 \(k=r\) 或 \(k=s\)，上述相应的空块与空基直接略去，结论仍成立。边界点处，\(B^{-1}\) 的矩阵元是光滑函数除以非零的 \(\det B\)，所以在半空间图中仍可光滑延拓；所有等式只在实际底空间邻域内使用，不要求延拓后的矩阵在外侧仍保持秩 \(k\)。
\end{proof}

\begin{example}
	\label{g1:exm:variable-rank-bundle-map}
	在 \(\mathbb R\) 上的两个乘积线丛之间，令 \(\Phi(t,v)=(t,tv)\)。当 \(t\ne0\) 时，它的核为零、像为整条纤维；在 \(t=0\) 时，核为整条纤维、像为零。故核与像都不能在零附近成为向量丛。即使把底空间改成 \([0,\infty)\)，同样的问题仍在边界点零发生。失败原因是秩的变化，不是边界使矩阵消元失效。
\end{example}

特别地，一个逐纤维单射的光滑丛同态把源丛同构到目标的像子丛。像是子丛由常秩定理得到；把同态看成到这个像的映射后，它逐纤维可逆，\cref{g1:prop:bundle-morphism-matrices}保证逆光滑。这允许我们把一个抽象的单射丛同态与其实现的子丛明确对应起来。
\subsection{拉回丛、直和与商丛}
\label{g1:subsec:1201:03}

设 \(f:N\to M\) 光滑，\(E\to M\) 是向量丛。如果希望在 \(N\) 的每点使用 \(f(q)\) 处的纤维，就应把新纤维定义为 \(E_{f(q)}\)。不同的 \(q\) 即使具有相同的像，也必须保留为不同的基点，因此不能只取 \(E\) 中这些纤维的并，而要把 \(q\) 一同记下来。

\begin{proposition}[拉回丛]
	\label{g1:prop:pullback-bundle}
	设 \(M,N\) 为 Hausdorff、第二可数的光滑流形，均可带边界；\(f:N\to M\) 光滑，\(\pi:E\to M\) 为秩 \(r\) 的实光滑向量丛。集合
	\[
		f^*E=\{(q,v)\in N\times E:f(q)=\pi(v)\},\qquad
		\pi'(q,v)=q
	\]
	带有自然的秩 \(r\) 实光滑向量丛结构，其拓扑是 \(N\times E\) 的子空间拓扑，称为 \(E\) 沿 \(f\) 的\term[lahuicong]{拉回丛}[pullback bundle]。映射 \(\widetilde f:f^*E\to E\)，\(\widetilde f(q,v)=v\)，是覆盖 \(f\) 的光滑丛同态。

	它具有如下泛性质：对任意 \(N\) 上的有限秩实光滑向量丛 \(H\)，以及覆盖 \(f\) 的光滑丛同态 \(\Psi:H\to E\)，存在唯一覆盖 \(\operatorname{id}_N\) 的光滑丛同态 \(\widehat\Psi:H\to f^*E\)，使 \(\widetilde f\widehat\Psi=\Psi\)。
\end{proposition}
\begin{proof}
	若 \(e_i\) 是 \(E\) 在 \(U_i\) 上的局部标架，定义
	\[
		\tau_i':f^{-1}(U_i)\times\mathbb R^r\longrightarrow(\pi')^{-1}(f^{-1}(U_i)),\qquad
		(q,a)\longmapsto(q,e_i(f(q))a).
	\]
	它是双射，逆为 \((q,v)\mapsto(q,[v]_i)\)。两个公式对所给子空间拓扑都连续，且定义域与像分别在相应局部范围内开，所以 \(\tau_i'\) 是局部平凡化的拓扑候选。它们的过渡映射为
	\[
		(q,a_j)\longmapsto(q,g_{ij}(f(q))a_j).
	\]
	这些矩阵光滑、可逆，且满足余循环条件。由\cref{g1:thm:bundle-cocycle-gluing}，得到唯一相容的光滑向量丛结构；其拓扑与上述子空间拓扑相同，因为两者在覆盖总空间的每个局部乘积上都由同一同胚给定。Hausdorff 性与第二可数性也可直接由 \(N\times E\) 的子空间性质得到。边界来自 \(N\) 的半空间坐标，纤维仍为完整的 \(\mathbb R^r\)。

	映射 \(\widetilde f\) 在这些图中写成 \((q,a)\mapsto(f(q),a)\)，故光滑，并且逐纤维线性。对泛性质，唯一可能的公式是
	\[
		\widehat\Psi(u)=\bigl(\pi_H(u),\Psi(u)\bigr).
	\]
	因为 \(\pi\Psi=f\pi_H\)，这个有序对确实落在 \(f^*E\) 中。局部矩阵中，它与 \(\Psi\) 使用相同的纤维矩阵，只是底映射变为恒等，所以光滑且逐纤维线性。此公式同时给出 \(\widetilde f\widehat\Psi=\Psi\) 与 \(\pi'\widehat\Psi=\pi_H\)，而这两个等式决定有序对的两个分量，故唯一。
\end{proof}

这两个交换关系可集中画成
\[
	\begin{tikzcd}[column sep=large,row sep=large]
		H \arrow[rr,bend left=16,"\Psi"] \arrow[dr,"\pi_H"'] \arrow[r,dashed,"\widehat\Psi"] & f^*E \arrow[r,"\widetilde f"] \arrow[d,"\pi'"'] & E \arrow[d,"\pi"] \\
		& N \arrow[r,"f"'] & M.
	\end{tikzcd}
\]
虚线箭头不增加一种任意选择；它只是把 \(\Psi(u)\) 与其新基点记录在一起。特别地，对 \(E\) 的截面 \(s\)，
\[
	(f^*s)(q)=(q,s(f(q)))
\]
是 \(f^*E\) 的光滑截面。对流形映射的微分，这一构造把覆盖 \(f\) 的 \(\mathrm{T}f:\mathrm{T}N\to \mathrm{T}M\) 改写成同底丛同态 \(\mathrm{T}N\to f^*\mathrm{T}M\)，其在 \(q\) 处仍是 \(\mathrm{d}f_q\)。纤维 \((f^*\mathrm{T}M)_q\) 是 \(\mathrm{T}_{f(q)}M\)，通常不是 \(\mathrm{T}_qN\)。

\ProseParagraphBreak

如果再有光滑 \(h:P\to N\)，则
\[
	h^*(f^*E)\longrightarrow(f\circ h)^*E,
	\qquad(p,(h(p),v))\longmapsto(p,v)
\]
是自然丛同构，其逆把 \((p,v)\) 送到 \((p,(h(p),v))\)。在从 \(E\) 拉回的局部平凡化中，两个公式的纤维分量均为恒等，故光滑。类似地，\(\operatorname{id}_M^*E\simeq E\)。这些明确同构说明，连续两次拉回与沿复合映射拉回给出相同的几何对象。

\ProseParagraphBreak

线性代数的直和则保留同一基点上的两种向量。对 \(M\) 上的秩 \(r,s\) 向量丛 \(E,F\)，定义\term[zhihecong]{直和丛}[direct sum bundle]
\[
	E\oplus F=\{(u,v)\in E\times F:\pi_E(u)=\pi_F(v)\},
	\qquad(E\oplus F)_p=E_p\oplus F_p.
\]
在共同平凡化下，它为 \(U_i\times(\mathbb R^r\oplus\mathbb R^s)\)，过渡矩阵为
\[
	\begin{pmatrix}g_{ij}^E&0\\0&g_{ij}^F\end{pmatrix}.
\]
分块乘法验证余循环条件，粘合定理给出光滑丛结构。局部表示与在 \(E\times F\) 中的子空间拓扑相容，证明与拉回丛的局部表示相同：从局部乘积到这个集合的映射及其逆都由两组平凡化给出。逐纤维投影和包含在局部都是固定线性映射，所以是光滑丛同态。这里的总空间不是整个 \(E\times F\)：直和要求两个向量属于同一基点，因而其维数为 \(\dim M+r+s\)。

\ProseParagraphBreak

商丛的目标是在每条纤维中忽略一个指定子空间的分量。虽然计算时可补选基来表示剩余分量，但商本身不应依赖所补的基，也不应要求预先选出全局补空间。

\begin{proposition}[商丛]
	\label{g1:prop:quotient-vector-bundle}
	设 \(M\) 为 Hausdorff、第二可数的光滑流形，可带边界；\(E\to M\) 为秩 \(r\) 实光滑向量丛，\(F\subset E\) 为秩 \(k\) 光滑子丛。集合
	\[
		E/F=\coprod_{p\in M}E_p/F_p
	\]
	有自然的秩 \(r-k\) 实光滑向量丛结构，其总空间拓扑为映射 \(q:E\to E/F\)，\(v_p\mapsto[v_p]\)，确定的商拓扑。\(q\) 是光滑满射丛同态，其逐纤维核为 \(F\)。如果 \(H\to M\) 是实光滑向量丛，\(\Phi:E\to H\) 是同底光滑丛同态，且 \(\Phi|_F=0\)，则存在唯一光滑丛同态 \(\overline\Phi:E/F\to H\) 使 \(\Phi=\overline\Phi q\)。
\end{proposition}
\begin{proof}
	由\cref{g1:prop:subbundle-adapted-frame}，局部可取 \(E\) 的标架 \(e_i\)，其前 \(k\) 个向量组成 \(F\) 的标架。把分量写成 \((a,b)\in\mathbb R^k\oplus\mathbb R^{r-k}\)。在每条纤维中，两个向量代表同一个商类，当且仅当它们的 \(b\) 分量相同。因此
	\[
		[e_i(p)(a,b)]\longmapsto(p,b)
	\]
	是商集合在 \(U_i\) 上到 \(U_i\times\mathbb R^{r-k}\) 的双射。

	在交叠上，两组适配标架的转换保持前 \(k\) 维子空间，故其矩阵必为
	\[
		g_{ij}=\begin{pmatrix}B_{ij}&C_{ij}\\0&D_{ij}\end{pmatrix}.
	\]
	其限制到 \(F_p\) 可逆，所以 \(B_{ij}\) 可逆；再由 \(\det g_{ij}=\det B_{ij}\det D_{ij}\ne0\)，知 \(D_{ij}\) 也可逆。商分量转换恰为 \(b_i=D_{ij}b_j\)。原余循环等式的右下块给出 \(D_{ij}D_{jk}=D_{ik}\)，于是粘合定理赋予商集合光滑向量丛结构。换用另一组适配标架时，同样的分块计算给出光滑可逆过渡，所以所得结构与补基选择无关。

	按这些局部图，\(q\) 的表达是 \((p,a,b)\mapsto(p,b)\)，故连续、光滑，并且开。开性可在每个 \(U_i\) 上验证，因为 Euclid 纤维投影把开集送到开集。具体地，对总空间开集 \(O\)，\(q(O)\) 在商丛限制到 \(U_i\) 中的部分，正是 \(O\cap E|_{U_i}\) 的局部投影像，因而开。连续开满射总是商映射：如果 \(q^{-1}(A)\) 开，则 \(A=q(q^{-1}(A))\) 开。因此这里的拓扑确实等于原集合商的商拓扑。核、满射性与秩都由同一局部公式得到。

	最后，逐纤维定义 \(\overline\Phi([v])=\Phi(v)\)。若换代表，差向量落在 \(F\)，被 \(\Phi\) 送到零，所以定义无关选择。在适配标架与 \(H\) 的任意局部标架下，\(\Phi\) 的前 \(k\) 列为零，矩阵写成 \((0,A_i)\)。\(\overline\Phi\) 的局部矩阵就是光滑矩阵 \(A_i\)，故它光滑。由于 \(q\) 满射，因子映射也唯一。
\end{proof}

对常秩丛同态 \(\Phi:E\to F\)，现在得到纤维中的第一同构定理的光滑版本：
\begin{equation}
	\label{g1:eq:bundle-first-isomorphism}
	E/\ker\Phi\ \simeq\ \operatorname{im}\Phi,
	\qquad [v_p]\longmapsto\Phi_p(v_p).
\end{equation}
核与像由\cref{g1:thm:constant-rank-bundle-morphism}成为子丛；商丛泛性质给出这个光滑映射。线性代数说明它在每条纤维上双射，而\cref{g1:prop:bundle-morphism-matrices}使其逆也光滑。因此这不是只在每个点分别存在的一族同构，而是总空间上的丛同构。

\begin{definition}
	\label{g1:def:short-exact-bundle-sequence}
	设 \(F,E,Q\) 是同一光滑流形上的有限秩实光滑向量丛，\(\iota:F\to E\)、\(p:E\to Q\) 为同底光滑丛同态。如果对每个底点 \(x\)，\(\iota_x\) 单射、\(p_x\) 满射且 \(\operatorname{im}\iota_x=\ker p_x\)，则称
	\[
		0\longrightarrow F\xrightarrow{\ \iota\ }E\xrightarrow{\ p\ }Q\longrightarrow0
	\]
	为向量丛的\term[duanzhengheconglie]{短正合列}[short exact sequence]。
\end{definition}

在这样的列中，\(\iota\) 将 \(F\) 识别为 \(E\) 的像子丛，而\eqref{g1:eq:bundle-first-isomorphism}把 \(Q\) 识别为 \(E/\iota(F)\)。局部补全 \(\iota(F)\) 的标架后，商分量可用补出的基向量提升，因而总能局部选择光滑右逆 \(Q|_U\to E|_U\)。不同邻域补出的子空间却未必彼此相同；把这些局部选择组织成整体选择，需要随后单位分解所提供的工具。即使整体选择存在，它通常也不由商丛本身唯一决定。

\subsection{标架丛与全局标架}
\label{g1:subsec:1201:04}

局部标架使向量有了分量，但一个纤维里并没有优先选定的一组基。若把所有可能的基都保留下来，换基矩阵就从计算中的辅助量变成这个新空间上的作用。这样做也能准确表达“能否在整个底空间上光滑选基”。

\begin{definition}
	\label{g1:def:frame-bundle}
	设 \(E\to M\) 为秩 \(r\) 的实光滑向量丛。定义
	\[
		\operatorname{Fr}(E)_p=\{u:\mathbb R^r\to E_p\mid u\text{ 是线性同构}\},\qquad
		\operatorname{Fr}(E)=\coprod_{p\in M}\operatorname{Fr}(E)_p.
	\]
	每个 \(u\) 等价于有序基 \((u(\varepsilon_1),\ldots,u(\varepsilon_r))\)。对任一局部标架 \(e_i\)，将 \(u\in\operatorname{Fr}(E)_p\) 唯一写成 \(u=e_i(p)A_i\)，其中 \(A_i\in\operatorname{GL}(r,\mathbb R)\)，并规定 \(u\mapsto(p,A_i)\) 为局部平凡化。若 \(e_j=e_i g_{ij}\)，则 \(A_i=g_{ij}A_j\)；用这些平凡化指定拓扑与光滑结构。所得空间称为 \(E\) 的\term[biaojiacong]{标架丛}[frame bundle]，其投影为 \(u\mapsto p\)。
\end{definition}

在 \(U_i\) 上选定参考标架 \(e_i\) 后，任意 \(u\in\operatorname{Fr}(E)_p\) 都唯一写成
\[
	u=e_i(p)A_i,\qquad A_i\in\operatorname{GL}(r,\mathbb R).
\]
因此得到 \(\operatorname{Fr}(E)|_{U_i}\simeq U_i\times\operatorname{GL}(r,\mathbb R)\)。若参考标架满足 \(e_j=e_i g_{ij}\)，则同一 \(u\) 的矩阵满足
\begin{equation}
	\label{g1:eq:frame-bundle-transition}
	A_i=g_{ij}A_j.
\end{equation}
这些过渡映射光滑可逆，并由余循环条件相容。它们赋予 \(\operatorname{Fr}(E)\) 唯一的拓扑与光滑结构，理由可直接沿用粘合证明的局部步骤：局部乘积中的开集通过过渡同胚相容，故组成拓扑基；不同基点用底空间的开邻域分离，同一基点的不同标架在同一局部乘积中分离；从底空间取可数平凡化子覆盖，再取各乘积的可数基，得到第二可数性。这里 \(\operatorname{GL}(r,\mathbb R)\) 是矩阵空间中行列式非零的开集，因此维数为 \(r^2\)，总空间维数为 \(\dim M+r^2\)。若底空间有边界，其边界仍是底空间边界的逆像。标架纤维是可逆线性映射的集合，一般不是向量空间，所以这个构造没有把 \(\operatorname{Fr}(E)\) 再当作向量丛。

\ProseParagraphBreak

群 \(\operatorname{GL}(r,\mathbb R)\) 在标架丛上有自然右作用
\begin{equation}
	\label{g1:eq:frame-bundle-right-action}
	u\cdot B=u\circ B.
\end{equation}
确为右作用，因为 \((u\cdot B)\cdot C=u\circ(BC)=u\cdot(BC)\)。局部表示为 \((p,A)\cdot B=(p,AB)\)，故光滑。作用在每条纤维内自由且传递：\(uB=u\) 蕴含 \(B=I_r\)；给定同一纤维内的 \(u,v\)，唯一使 \(uB=v\) 的矩阵为 \(B=u^{-1}\circ v\)。因此任意两组基都可通过唯一的矩阵相互转换，却没有哪一组基天然是单位元；只有选定参考标架后，纤维才被识别为矩阵群。这是以后讨论主丛时要保留的结构。

\begin{proposition}[全局标架与平凡丛]
	\label{g1:prop:global-frame-triviality}
	设 \(M\) 为 Hausdorff、第二可数光滑流形，可带边界，\(E\to M\) 为秩 \(r\) 实光滑向量丛。下列条件等价：
	\begin{equivalences}
		\item \(E\) 与乘积向量丛 \(M\times\mathbb R^r\) 光滑同构。
		\item 存在 \(r\) 个全局光滑截面，在每个点组成纤维的有序基。
		\item 标架丛投影 \(\operatorname{Fr}(E)\to M\) 有光滑截面。
	\end{equivalences}
	满足这些条件时，称 \(E\) 为\term[pingfancong]{平凡丛}[trivial bundle]。
\end{proposition}
\begin{proof}
	若给定丛同构 \(\Theta:M\times\mathbb R^r\to E\)，令 \(s_a(p)=\Theta(p,\varepsilon_a)\)。这些截面光滑，且由 \(\Theta\) 的逐纤维线性同构性构成基，故第一项推出第二项。

	给定第二项中的全局标架，定义 \(\sigma(p):\mathbb R^r\to E_p\) 为 \(\sigma(p)v=\sum_a v^a s_a(p)\)。在任一参考局部标架下，\(s_a\) 的分量组成一个光滑可逆矩阵，因此 \(\sigma\) 是标架丛的光滑截面。反过来，给定这样的 \(\sigma\)，令 \(s_a(p)=\sigma(p)\varepsilon_a\)，同一局部矩阵计算给出光滑全局标架，故第二、三项等价。

	最后，给定 \(\sigma\)，令 \(\Theta(p,v)=\sigma(p)v\)。它逐纤维为线性同构，局部表达为 \((p,v)\mapsto(p,A(p)v)\)，其中 \(A\) 光滑且可逆，所以 \(\Theta\) 与其逆 \((p,w)\mapsto(p,A(p)^{-1}w)\) 都光滑。这证明第一项。秩零时取空标架，所有构造与结论仍成立。
\end{proof}

Möbius 线丛没有处处非零的截面，因而连秩一所需的全局标架也没有。另一方面，切丛有全局标架并不要求流形具有一张全局坐标图。例如圆周上的切向量 \((-\sin t,\cos t)\) 随 \(e^{\mathrm{i}t}\) 光滑且处处非零，给出 \(\mathrm{T}S^1\) 的全局标架；圆周仍不能由一张 Euclid 开区间坐标图覆盖。标架提供的是每个点的线性坐标，流形坐标还须把邻域中的不同基点同时参数化，这两种要求应当区分。
% 纲要条目排为编号小节标题，并进入目录与书签；说明段仍保留于纲要注释中。

% !TeX root = ../../Book1.tex
% 纲要标题：### 12.2 几何中的多线性代数
% 纲要位置：几何学计划纲要.md，第 579 行。
% 正文状态：初稿；尚未正式审核。

% 纲要内容开始
% 
% * 张量积的泛性质与收缩
% * 对偶、对称幂及外幂
% * 交替性、楔积与分次符号
% * 行列式线、密度与体积
% * 不同标架下的分量及不变量
% 
% 纲要内容结束

\section{几何中的多线性代数}
\label{g1:sec:1202}

在一点处，方向属于切空间，测量方向的线性函数属于余切空间。内积却要同时接收两个方向，有向面积也要同时接收两个方向，而且交换这两个方向时会变号。要把这些量沿流形逐点组织起来，先须知道它们在一个向量空间里是什么，以及换基时怎样保持同一个对象。本节暂时固定有限维实向量空间；随点变化及光滑性将在\secref{g1:sec:1204}处理。

\ProseParagraphBreak

所需的线性代数基础是基、对偶、线性映射和行列式。回顾\cref{g1:def:cotangent-space,g1:prop:covector-pullback}：\(V^*=\operatorname{Hom}_{\mathbb R}(V,\mathbb R)\) 中的元素把向量读成实数；线性映射 \(A:V\to W\) 的对偶为 \(A^*\alpha=\alpha\circ A\)，方向从 \(W^*\) 返回 \(V^*\)。下面把一次线性观测推广到多次线性观测，同时保留不依赖基的定义。

\subsection{张量积的泛性质与收缩}
\label{g1:subsec:1202:01}

设 \(V,W,U\) 为实向量空间。对映射 \(b:V\times W\to U\)，若固定任一输入后，它对另一输入线性，就称该映射为双线性的。内积就是这样的映射，但它一般不是乘积向量空间 \(V\times W\) 上的线性映射：例如同时将两个输入加倍，输出会乘以 \(4\)。我们希望构造一个新的向量空间，使每个双线性运算都能写成“先作同一个规定的运算，再作一个普通线性映射”。

\begin{theorem}[张量积的泛性质]
	\label{g1:thm:tensor-product-universal}
	对有限维实向量空间 \(V,W\)，存在实向量空间 \(V\otimes W\) 及双线性映射
	\[
		\tau:V\times W\longrightarrow V\otimes W,\qquad (v,w)\longmapsto v\otimes w,
	\]
	使得对任意实向量空间 \(U\) 及双线性映射 \(b:V\times W\to U\)，存在唯一线性映射 \(\widetilde b:V\otimes W\to U\) 满足 \(\widetilde b(v\otimes w)=b(v,w)\)。这样的二元组 \((V\otimes W,\tau)\) 在保持规定映射的唯一线性同构意义下唯一。若 \(e_1,\ldots,e_m\) 和 \(f_1,\ldots,f_n\) 分别为 \(V,W\) 的基，则
	\[
		\{e_i\otimes f_j:1\le i\le m,\ 1\le j\le n\}
	\]
	为 \(V\otimes W\) 的基，因而 \(\dim(V\otimes W)=mn\)。
\end{theorem}

\[
	\begin{tikzcd}[column sep=large,row sep=large]
		V\times W \arrow[r,"\tau"] \arrow[dr,"b"'] & V\otimes W \arrow[d,dashed,"\exists!\,\widetilde b"] \\
		& U
	\end{tikzcd}
\]
图中的 \(\tau,b\) 是双线性映射，虚线箭头是线性映射。交换性只说从 \((v,w)\) 出发的两条路径给出同一个值；唯一性则说任何双线性运算都恰好对应一个这样的线性运算。

\begin{proof}
	先用一组基制造所需空间，再证明结果不依赖这次选择。取一个以符号 \(E_{ij}\) 为基的 \(mn\) 维实向量空间 \(T\)，定义
	\[
		\tau\left(\sum_i v^ie_i,\sum_j w^jf_j\right)=\sum_{i,j}v^iw^jE_{ij}.
	\]
	每个系数对两组变量分别线性，故 \(\tau\) 双线性。给定 \(b\)，规定 \(\widetilde b(E_{ij})=b(e_i,f_j)\)，唯一地定义一个线性映射。由双线性，
	\[
		b(v,w)=\sum_{i,j}v^iw^j b(e_i,f_j)=\widetilde b(\tau(v,w)).
	\]
	反之，任何满足这个等式的线性映射都必须在 \(E_{ij}=\tau(e_i,f_j)\) 上取上述值，故唯一。把 \(T\) 记为 \(V\otimes W\)，把 \(\tau(v,w)\) 记为 \(v\otimes w\)，便得到构造及所述基。

	若另一个二元组 \((T',\tau')\) 具有同样的性质，分别取测试映射 \(\tau'\) 和 \(\tau\)，得到唯一线性映射 \(F:T\to T'\)、\(G:T'\to T\)，满足 \(F\tau=\tau'\)、\(G\tau'=\tau\)。于是 \(GF\tau=\tau\)，而恒等映射也具有这个性质；唯一性给出 \(GF=\operatorname{id}_T\)。同理 \(FG=\operatorname{id}_{T'}\)。因此换一组基构造的空间与原空间有一个由规定元素决定的唯一同构。这正是构造不依赖基的含义，而不是说坐标列表本身不改变。零维情形中 \(T=\{0\}\)，双线性映射也只能为零，上述结论仍成立。
\end{proof}

称 \(V\otimes W\) 为 \(V,W\) 的\term[zhangliangji]{张量积}[tensor product]。元素 \(v\otimes w\) 称为纯张量；一般元素是纯张量的有限和，未必本身就是一个纯张量。双线性意味着
\[
	\begin{aligned}
		(v+v')\otimes w&=v\otimes w+v'\otimes w,\\
		v\otimes(w+w')&=v\otimes w+v\otimes w',\\
		(cv)\otimes w&=v\otimes(cw)=c(v\otimes w).
	\end{aligned}
\]
这些关系不允许把一切和都合成一个纯张量。例如在两个二维空间中，\(e_1\otimes f_1+e_2\otimes f_2\) 若等于 \(v\otimes w\)，其系数矩阵将既是单位矩阵，又是列向量 \((v^i)\) 与行向量 \((w^j)\) 的乘积；后者的列都在同一条直线上，不可能等于单位矩阵。

\ProseParagraphBreak

泛性质也规定了张量积上的映射。若 \(A:V\to V'\)、\(B:W\to W'\) 线性，则 \((v,w)\mapsto Av\otimes Bw\) 双线性，因而唯一给出 \(A\otimes B\)，满足 \((A\otimes B)(v\otimes w)=Av\otimes Bw\)。复合和恒等式在纯张量上成立，便在整个空间成立，因为纯张量张成该空间。同样，对三个或更多输入可按基构造多线性映射的通用空间；\((V\otimes W)\otimes U\) 与 \(V\otimes(W\otimes U)\) 都有这一性质，规定 \((v\otimes w)\otimes u\mapsto v\otimes(w\otimes u)\) 的同构使我们能够省略括号。记 \(V^{\otimes k}\) 为 \(k\) 个因子的张量积，并约定 \(V^{\otimes0}=\mathbb R\)。

\begin{proposition}[线性映射的张量表示]
	\label{g1:prop:tensor-hom-duality}
	设 \(V,W\) 为有限维实向量空间。映射
	\[
		\Phi:W\otimes V^*\longrightarrow\operatorname{Hom}_{\mathbb R}(V,W),\qquad
		\Phi(w\otimes\alpha)(v)=\alpha(v)w
	\]
	是无需选基的线性同构。若 \(e_1,\ldots,e_n\) 为 \(V\) 的任意一组基，\(e^1,\ldots,e^n\) 为对偶基，则
	\[
		\Phi^{-1}(L)=\sum_i L(e_i)\otimes e^i.
	\]
	此外，求值映射 \(V\to V^{**}\)、\(v\mapsto[\alpha\mapsto\alpha(v)]\) 为自然线性同构。
\end{proposition}
\begin{proof}
	\((w,\alpha)\mapsto[v\mapsto\alpha(v)w]\) 是双线性的，故由泛性质得到 \(\Phi\)。将给定 \(L\) 对应的张量作用于 \(v=\sum_i e^i(v)e_i\)，得到
	\[
		\Phi\left(\sum_iL(e_i)\otimes e^i\right)(v)=\sum_i e^i(v)L(e_i)=L(v).
	\]
	反过来，对纯张量 \(w\otimes\alpha\)，先经 \(\Phi\) 再用所列逆公式，得到
	\[
		\sum_i\alpha(e_i)w\otimes e^i
		=w\otimes\left(\sum_i\alpha(e_i)e^i\right)=w\otimes\alpha.
	\]
	由线性可知两个映射互逆。虽然逆公式用了基，它是一个无基映射的逆，所以换基所得张量相同。最后，非零向量 \(v\) 至少有一个非零坐标，相应的对偶基元素在 \(v\) 上不为零，故求值映射 \(V\to V^{**}\) 单射；两空间同维，因而是同构。它由求值定义，也没有基的选择。
\end{proof}

一般地，定义
\[
	T^r_s(V)=V^{\otimes r}\otimes(V^*)^{\otimes s},\qquad r,s\ge0,
\]
称其元素为 \((r,s)\) 型\term[zhangliang]{张量}[tensor]。\symidx{\(T^r_s(V)\)：\((r,s)\) 型张量空间}它自然作用于 \(r\) 个协变向量和 \(s\) 个向量：纯张量 \(v_1\otimes\cdots\otimes v_r\otimes\alpha_1\otimes\cdots\otimes\alpha_s\) 对输入 \((\xi_1,\ldots,\xi_r;w_1,\ldots,w_s)\) 的值是
\[
	\prod_{a=1}^r\xi_a(v_a)\prod_{b=1}^s\alpha_b(w_b).
\]
这不只是举出一种作用方式，而是给出了张量空间与相应多线性函数空间的同构：选基后，任一多线性函数都由它在各输入基上的值唯一决定，逐项展开便得到张量基的唯一系数表示。这里把向量视为对协变向量的线性函数，使用的正是 \(V\simeq V^{**}\)。因此上标 \(r\) 数的是向量因子，实际输入却是 \(r\) 个协变向量；不要把两者混淆。

\begin{proposition}[收缩与迹]
	\label{g1:prop:tensor-contraction}
	设 \(V\) 为有限维实向量空间，\(r,s\ge1\)。指定一个向量因子的位置 \(a\) 和一个对偶因子的位置 \(b\)，在纯张量上将它们求值、删除这两个因子，并保留其余因子的先后次序，得到唯一线性映射
	\[
		C^a_b:T^r_s(V)\longrightarrow T^{r-1}_{s-1}(V).
	\]
	这一运算不依赖基，称为\term[shousuo]{收缩}[contraction]。特别地，经 \(V\otimes V^*\simeq\operatorname{End}(V)\) 识别，\(C^1_1\) 给出线性映射 \(L\) 的迹：
	\[
		\operatorname{tr}L=\sum_i e^i(L(e_i)).
	\]
	右端在任意基及其对偶基中均取同一值。
\end{proposition}
\begin{proof}
	具体规定为
	\[
		C^a_b(v_1\otimes\cdots\otimes v_r\otimes\alpha_1\otimes\cdots\otimes\alpha_s)
		=\alpha_b(v_a)\,v_1\otimes\cdots\widehat{v_a}\cdots\otimes v_r
		\otimes\alpha_1\otimes\cdots\widehat{\alpha_b}\cdots\otimes\alpha_s,
	\]
	帽号表示省略。右端对原来的每个输入分别线性，多因子张量积的泛性质便给出唯一线性映射。这个定义只使用内在的求值 \(\alpha_b(v_a)\)，没有使用基。对于 \(r=s=1\)，\cref{g1:prop:tensor-hom-duality}给出 \(L=\Phi(\sum_iL(e_i)\otimes e^i)\)，收缩后正是所列求和。因此矩阵对角元之和是一个无基运算的坐标表达，换基不改变它。
\end{proof}

收缩为什么只配对一个向量因子和一个对偶因子？因为 \(\alpha(v)\) 已有定义，而两个向量之间并没有预定的实数值。若另选内积，当然可以先用内积把一个向量变成协变向量再收缩，但这时结果包含了所选内积的信息。例如一个双线性形式属于 \(T^0_2(V)\)；仅凭这个张量本身没有规定的“迹”，要用内积的逆张量把它的两个指标配对，才能谈相对于该内积的迹。

\subsection{对偶、对称幂及外幂}
\label{g1:subsec:1202:02}

有些多线性观测不区分输入的排列，例如内积；另一些观测在交换两个输入时改变符号，例如二维有向面积。与其在每次计算时附加相同条件，不如把这些条件直接加入张量积的构造。这里的商空间只是在说：若所有允许的观测都必须把某种差别读成零，就把这种差别规定为零。

\begin{definition}
	\label{g1:def:symmetric-exterior-powers}
	设 \(V\) 为有限维实向量空间，\(k\ge2\)。其 \(k\) 次\term[duichengmi]{对称幂}[symmetric power]为
	\[
		\operatorname{Sym}^kV=V^{\otimes k}/R_{\mathrm{sym}},
	\]
	其中 \(R_{\mathrm{sym}}\) 由所有纯张量与其任意排列之差张成。纯张量的像记为 \(v_1\odot\cdots\odot v_k\)。其 \(k\) 次\term[waimi]{外幂}[exterior power]为
	\[
		\Lambda^kV=V^{\otimes k}/R_{\mathrm{alt}},
	\]
	其中 \(R_{\mathrm{alt}}\) 由任意两个位置输入相同的纯张量张成；纯张量的像记为 \(v_1\wedge\cdots\wedge v_k\)。约定 \(\operatorname{Sym}^0V=\Lambda^0V=\mathbb R\)，\(\operatorname{Sym}^1V=\Lambda^1V=V\)。
\end{definition}

外幂的定义立刻给出重复输入为零。将两个相同位置的输入都换成 \(u+v\) 并展开，两个重复项各自为零，余下两项相加为零。因此交换任意两个位置会变号。这也说明：一个多线性映射在重复输入时为零，当且仅当交换任意两个输入时变号；反向推导使用 \(a=-a\Rightarrow a=0\)，这里实数域的条件是有用的。这样的映射称为\term[jiaotixing]{交替的}[alternating]。

\ProseParagraphBreak

商的定义给出相应的泛性质：对称 \(k\) 线性映射唯一地通过 \(\operatorname{Sym}^kV\) 分解，交替 \(k\) 线性映射唯一地通过 \(\Lambda^kV\) 分解。理由是张量积已经把多线性映射变成线性映射；对称或交替条件恰好说这个线性映射消去相应关系子空间，因而能且只能下降到商。特别地，\((\Lambda^kV)^*\) 就是所有交替 \(k\) 线性实函数的空间。

\begin{proposition}[对称幂与外幂的基]
	\label{g1:prop:symmetric-exterior-bases}
	设 \(V\) 为 \(n\) 维实向量空间，基为 \(e_1,\ldots,e_n\)。对 \(k\ge1\)，
	\[
		e_1^{\odot a_1}\odot\cdots\odot e_n^{\odot a_n},
		\qquad a_i\ge0,\quad \sum_i a_i=k
	\]
	是 \(\operatorname{Sym}^kV\) 的一组基，而
	\[
		e_{i_1}\wedge\cdots\wedge e_{i_k},\qquad i_1<\cdots<i_k
	\]
	是 \(\Lambda^kV\) 的一组基。因此，当 \(n\ge1\) 时，
	\[
		\dim\operatorname{Sym}^kV=\binom{n+k-1}{k},\qquad
		\dim\Lambda^kV=\binom nk,
	\]
	其中 \(k>n\) 时外幂为零；零次幂均为一维，零维空间的正次幂均为零。
\end{proposition}
\begin{proof}
	展开各输入的基表示，并在对称幂中重新排列因子，得到第一组元素张成。为证独立，把每个 \(e_i\) 对应于形式变量 \(X_i\)，规定
	\[
		(v_1,\ldots,v_k)\longmapsto
		\prod_{\ell=1}^k\left(\sum_i e^i(v_\ell)X_i\right).
	\]
	它是取值于 \(k\) 次齐次多项式空间的对称多线性映射，故下降到 \(\operatorname{Sym}^kV\)，并把所列元素送到各不相同的单项式 \(X_1^{a_1}\cdots X_n^{a_n}\)。形式多项式的系数表示唯一，因此这些像线性无关，原元素也线性无关。非负整数解 \(a_1+\cdots+a_n=k\) 的数量是 \(\binom{n+k-1}{n-1}\)：把 \(k\) 个记号排成一行，再用 \(n-1\) 个隔板划分为各组，允许空组，即得到这个计数。

	在外幂中也先展开基表示；重复指标项为零，不重复项可按指标递增排列，代价只是排列的符号，所以第二组元素张成。对每个严格递增指标组 \(I=(i_1,\ldots,i_k)\)，定义
	\[
		D_I(v_1,\ldots,v_k)=\det\bigl[e^{i_a}(v_b)\bigr]_{a,b=1}^k.
	\]
	行列式对各列线性，重复列时为零，故它下降为 \(\Lambda^kV\) 上的线性函数。在所列的递增基元素 \(e_{j_1}\wedge\cdots\wedge e_{j_k}\) 上，它恰在 \(I=J\) 时取 \(1\)，其余取 \(0\)。逐个应用这些函数，就知所列元素独立，其数量为 \(\binom nk\)。若 \(k>n\)，每个展开项必有重复指标，所以全部为零。
\end{proof}

现在把外幂与交替形式准确地对应起来。对 \(\alpha_i\in V^*\) 和 \(v_j\in V\)，规定
\begin{equation}
	\label{g1:eq:exterior-determinant-pairing}
	\left\langle\alpha_1\wedge\cdots\wedge\alpha_k,\,
	v_1\wedge\cdots\wedge v_k\right\rangle
	=\det[\alpha_i(v_j)].
\end{equation}
行列式对两组输入分别多线性并交替，故由两次泛性质，这给出 \(\Lambda^k(V^*)\) 与 \(\Lambda^kV\) 之间的双线性配对。对偶基的递增楔积与向量基的递增楔积在此配对下互为对偶，因而得到自然同构
\[
	\Lambda^k(V^*)\simeq(\Lambda^kV)^*
	\simeq\{\text{\(V\) 上的交替 \(k\) 线性实函数}\}.
\]
这里“自然”具体指不需选基，而且与线性映射相容：若 \(A:V\to W\)，令 \(\Lambda^kA(v_1\wedge\cdots\wedge v_k)=Av_1\wedge\cdots\wedge Av_k\)，则配对的两边都给出同一个行列式 \(\det[\alpha_i(Av_j)]\)，所以 \((\Lambda^kA)^*\) 对应于 \(\Lambda^k(A^*)\)。以后简写 \(\Lambda^kV^*=\Lambda^k(V^*)\)，称其中元素为 \(V\) 上的 \(k\) 次形式。

\ProseParagraphBreak

对称幂有相似的对偶识别，但归一化不同。规定
\[
	\left\langle\alpha_1\odot\cdots\odot\alpha_k,\,
	v_1\odot\cdots\odot v_k\right\rangle
	=\frac1{k!}\sum_{\sigma\in S_k}\prod_{i=1}^k\alpha_i(v_{\sigma(i)}).
\]
这里 \(S_k\) 是 \(k\) 个位置的全排列集合。右端对两组输入都对称多线性，因而给出 \(\operatorname{Sym}^k(V^*)\to(\operatorname{Sym}^kV)^*\)。在上一命题的单项式基上，两个多重指标不同时配对为零；都为 \((a_1,\ldots,a_n)\) 时配对为 \(a_1!\cdots a_n!/k!\)，是非零实数。因此该映射也是无基的同构，并把 \(\operatorname{Sym}^k(V^*)\) 识别为对称 \(k\) 线性实函数。它与拉回相容，因为将各 \(v_j\) 换成 \(Av_j\)，等价于将各 \(\alpha_i\) 换成 \(A^*\alpha_i\)。例如
\[
	(\alpha\odot\beta)(u,v)=\frac{\alpha(u)\beta(v)+\alpha(v)\beta(u)}2,\qquad
	(\alpha\wedge\beta)(u,v)=\alpha(u)\beta(v)-\alpha(v)\beta(u).
\]
这两个公式中的系数差别应记住；本书的有向面积形式采用第二个公式，没有额外的 \(1/2\)。

\subsection{交替性、楔积与分次符号}
\label{g1:subsec:1202:03}

将两个外幂中的代表张量直接首尾相接，便得到乘法
\[
	\Lambda^pV\times\Lambda^qV\longrightarrow\Lambda^{p+q}V.
\]
它确实不依赖代表：若一个因子改变了一个重复输入生成的关系项，相接后的张量仍有重复输入，在目标商空间中为零。于是乘法由纯元素上的
\[
	(v_1\wedge\cdots\wedge v_p)\wedge(w_1\wedge\cdots\wedge w_q)
	=v_1\wedge\cdots\wedge v_p\wedge w_1\wedge\cdots\wedge w_q
\]
决定，称为\term[xieji]{楔积}[wedge product]。在 \(V^*\) 上也有同一构造。把它视为交替形式的乘法时，必须说明它对一般输入怎样求值。

\ProseParagraphBreak

先对任意 \(k\) 线性实函数 \(T\) 定义\term[jiaotihua]{交替化}[alternation]
\[
	\operatorname{Alt}_kT(v_1,\ldots,v_k)
	=\frac1{k!}\sum_{\sigma\in S_k}\operatorname{sgn}(\sigma)
	T(v_{\sigma(1)},\ldots,v_{\sigma(k)}).
\]
交换两个输入会将每项对应到一个符号相反的排列，所以结果交替；若 \(T\) 已交替，每一项连同符号都等于 \(T\)，故 \(\operatorname{Alt}_kT=T\)。于是 \(\operatorname{Alt}_k\) 是投影到交替形式空间的线性算子。由行列式展开，\(\alpha_1\wedge\cdots\wedge\alpha_k\) 作为形式等于 \(k!\operatorname{Alt}_k(\alpha_1\otimes\cdots\otimes\alpha_k)\)，不能把这里的阶乘删去。

\begin{proposition}[楔积的代数性质]
	\label{g1:prop:wedge-algebra}
	设 \(V\) 为有限维实向量空间，\(\alpha\in\Lambda^pV^*\)、\(\beta\in\Lambda^qV^*\)、\(\gamma\in\Lambda^rV^*\)。在行列式配对\eqref{g1:eq:exterior-determinant-pairing}的约定下，
	\begin{equation}
		\label{g1:eq:wedge-shuffle}
		\begin{split}
			(\alpha\wedge\beta)(v_1,\ldots,v_{p+q})
			&=\sum_{\sigma\in\operatorname{Sh}(p,q)}\operatorname{sgn}(\sigma)
			\alpha(v_{\sigma(1)},\ldots,v_{\sigma(p)})
			\beta(v_{\sigma(p+1)},\ldots,v_{\sigma(p+q)})\\
			&=\frac1{p!\,q!}\sum_{\sigma\in S_{p+q}}\operatorname{sgn}(\sigma)
			\alpha(v_{\sigma(1)},\ldots,v_{\sigma(p)})
			\beta(v_{\sigma(p+1)},\ldots,v_{\sigma(p+q)}).
		\end{split}
	\end{equation}
	其中 \(\operatorname{Sh}(p,q)\) 由满足 \(\sigma(1)<\cdots<\sigma(p)\)、\(\sigma(p+1)<\cdots<\sigma(p+q)\) 的排列组成。楔积双线性，并满足
	\[
		(\alpha\wedge\beta)\wedge\gamma=\alpha\wedge(\beta\wedge\gamma),\qquad
		\alpha\wedge\beta=(-1)^{pq}\beta\wedge\alpha.
	\]
	零次形式是实数，\(1\) 是乘法单位。对任意有限维实向量空间 \(W\) 及线性映射 \(A:V\to W\)，按
	\[
		(A^*\eta)(v_1,\ldots,v_k)=\eta(Av_1,\ldots,Av_k)
	\]
	定义形式的拉回，则 \(A^*(\eta\wedge\zeta)=A^*\eta\wedge A^*\zeta\)，且复合的拉回按相反顺序复合。
\end{proposition}
\begin{proof}
	先说明公式与已经定义的乘法一致。由于一形式的楔积张成各次外幂，只须取 \(\alpha=\alpha_1\wedge\cdots\wedge\alpha_p\)、\(\beta=\beta_1\wedge\cdots\wedge\beta_q\)。相接所得形式在 \(v_1,\ldots,v_{p+q}\) 上的值，是以 \(\alpha_i,\beta_j\) 为行的 \((p+q)\) 阶行列式。按哪些 \(p\) 列被分配给前 \(p\) 行将行列式展开归组，每组选列对应唯一的 \((p,q)\) 洗牌排列；组内的两个行列式正是 \(\alpha\) 和 \(\beta\) 的值，分组选列的符号就是该排列的符号。这得到第一式。任意全排列又唯一分为一个洗牌排列和两个块内排列。块内排列的符号会分别被 \(\alpha,\beta\) 的交替性抵消，所以一个洗牌项在全排列和中重复 \(p!q!\) 次，得到第二式。

	结合性也能直接从相接代表看出：两种加括号方式都给出同一个首尾相接的张量。用形式公式核对，展开任一边时，都是将所有输入分为内部递增、大小分别为 \(p,q,r\) 的三个有序块；每种分法出现一次，系数都是将原输入排列成三个块的符号。因此公式中的归一化与结合性一致。交换前后两个块时，前块的每一个位置都须越过后块的每一个位置，共 \(pq\) 次交换，故产生 \((-1)^{pq}\)，这就给出所列交换律。零次时没有输入，所有公式退化为实数乘法。

	拉回仍为交替多线性函数。将洗牌公式的每个 \(v_i\) 换成 \(Av_i\)，两边分别就是 \(A^*(\eta\wedge\zeta)\) 和 \(A^*\eta\wedge A^*\zeta\)。若再有线性映射 \(B:W\to Z\)，则两种拉回都把输入送成 \(BAv_i\)，所以 \((B\circ A)^*=A^*\circ B^*\)。
\end{proof}

这里的交换律称为分次交换律：符号由两个次数的乘积决定。特别地，奇数次形式满足 \(\alpha\wedge\alpha=0\)，因为它等于自身的相反数；偶数次形式没有这一结论。例如在四维空间中，
\[
	(e^1\wedge e^2+e^3\wedge e^4)^2
	=2e^1\wedge e^2\wedge e^3\wedge e^4\ne0.
\]
右边的两项相加，是因为交换两个二次块经过 \(2\cdot2\) 次换位，符号为正。由此也可见，“交替”指的是一个形式内部交换向量输入的规律，不能理解为任意两个形式相乘都要变号。

\begin{definition}
	\label{g1:def:interior-product}
	设 \(V\) 为实向量空间，\(k\ge1\) 为整数，\(v\in V\)，\(\alpha\in\Lambda^kV^*\)。定义\term[charu]{插入}[interior product]
	\[
		(\iota_v\alpha)(v_1,\ldots,v_{k-1})=\alpha(v,v_1,\ldots,v_{k-1}).
	\]
	它属于 \(\Lambda^{k-1}V^*\)，并约定 \(\iota_v c=0\) 对零次形式 \(c\in\mathbb R\) 成立。
\end{definition}

插入固定了一个输入，所以次数降低 \(1\)。输入固定在最左端是本书的约定；把它放在末端会产生与次数有关的符号，必须相应修改下面的公式。

\begin{proposition}[插入的乘法公式]
	\label{g1:prop:interior-product-rule}
	设 \(V\) 为有限维实向量空间，\(v\in V\)，\(\alpha\in\Lambda^pV^*\)、\(\beta\in\Lambda^qV^*\)。则
	\[
		\iota_v(\alpha\wedge\beta)
		=(\iota_v\alpha)\wedge\beta+(-1)^p\alpha\wedge(\iota_v\beta).
	\]
	因此 \(\iota_v\) 是外代数上次数为 \(-1\) 的分次导子。另有 \(\iota_v^2=0\)；若 \(A:V\to W\) 线性，则 \(\iota_v(A^*\eta)=A^*(\iota_{Av}\eta)\)。
\end{proposition}
\begin{proof}
	先取一形式 \(\eta_1,\ldots,\eta_k\)。把 \(\eta_1\wedge\cdots\wedge\eta_k\) 在 \((v,v_1,\ldots,v_{k-1})\) 上的行列式沿第一列展开，得到
	\[
		\iota_v(\eta_1\wedge\cdots\wedge\eta_k)
		=\sum_{j=1}^k(-1)^{j-1}\eta_j(v)
		\,\eta_1\wedge\cdots\widehat{\eta_j}\cdots\wedge\eta_k.
	\]
	若前 \(p\) 个因子来自 \(\alpha\)，后 \(q\) 个来自 \(\beta\)，前 \(p\) 项合成 \((\iota_v\alpha)\wedge\beta\)；第 \(p+j\) 项的符号为 \((-1)^{p+j-1}=(-1)^p(-1)^{j-1}\)，后 \(q\) 项因而合成第二项。纯楔积张成所有形式，按双线性便得到一般公式。\(p=0\) 或 \(q=0\) 时公式就是插入对常数乘法的线性性质。两次插入同一向量得到 \(\alpha(v,v,\ldots)=0\)，故 \(\iota_v^2=0\)；最后一式的两边在任意输入上都等于 \(\eta(Av,Av_1,\ldots)\)。
\end{proof}

\subsection{行列式线、密度与体积}
\label{g1:subsec:1202:04}

对 \(n\) 维空间，最高次外幂只有一维，却保存着换基的行列式。若把一组有序基写作行向量 \(e=(e_1,\ldots,e_n)\)，换基约定为
\[
	e'=eA,\qquad e'_j=\sum_i e_i A^i{}_j,\qquad A\in\operatorname{GL}(n,\mathbb R),
\]
那么按多线性展开并消去重复项，有
\[
	e'_1\wedge\cdots\wedge e'_n=(\det A)e_1\wedge\cdots\wedge e_n.
\]
行列式的符号记录定向，绝对值记录平行多面体的体积缩放。要把这两个作用区分开，就会得到不同的一维空间。

\begin{definition}
	\label{g1:def:determinant-density-line}
	设 \(V\) 为 \(n\) 维实向量空间。其\term[hanglieshixian]{行列式线}[determinant line]为 \(\det V=\Lambda^nV\)。最高次形式组成的是另一条线 \(\det(V^*)=\Lambda^n(V^*)\)，经行列式配对自然同构于 \((\det V)^*\)。

	\(V\) 上的一个\term[midu]{密度}[density]是有序基集合上的实值函数 \(\mu\)，满足对每组有序基 \(e\) 和每个可逆实矩阵 \(A\)，
	\[
		\mu(eA)=|\det A|\,\mu(e).
	\]
	全体密度按逐点加法与数乘构成实向量空间，记为 \(\mathcal D(V)\)，也记为 \(|\Lambda^nV^*|\)。若 \(\mu(e)>0\) 对一组、从而对每组有序基成立，称 \(\mu\) 为正密度。
\end{definition}

密度线的记号 \(|\Lambda^nV^*|\) 是一种惯用记号，其中的元素仍可取负值，不能把它理解为只有非负元素的集合。固定一组基 \(e\) 后，任意指定实数 \(c\) 都唯一地定义密度：若 \(f=eA\)，令 \(\mu(f)=|\det A|c\)。两个连续换基的行列式绝对值相乘，故这个定义满足规定关系。因此 \(\mu\mapsto\mu(e)\) 是 \(\mathcal D(V)\to\mathbb R\) 的线性同构，密度线确为一维；正密度只构成其中一个正半轴，不是线性子空间。

\ProseParagraphBreak

最高次形式 \(\omega\in\Lambda^nV^*\) 给出密度 \(|\omega|\)，定义为 \(|\omega|(e)=|\omega(e_1,\ldots,e_n)|\)。这是因为 \(\omega(eA)=\det A\,\omega(e)\)。但是 \(\omega\mapsto|\omega|\) 不是线性同构：\(\omega\) 和 \(-\omega\) 给出同一密度，而两个相反形式之和为零。对 \(n\ge1\)，定向可由有序基按正行列式换基划成的两类中选择一类来指定。选定定向以后，令 \(s(e)=1\) 对正向基成立，\(s(e)=-1\) 对负向基成立，则
\[
	\omega\longmapsto\bigl[e\longmapsto s(e)\omega(e)\bigr]
\]
才是最高次形式线到密度线的线性同构。正密度与所选定向一起唯一确定一个在正向基上取正值的最高次形式；仅有正密度时，两个可能的定向仍未选定。

\ProseParagraphBreak

这些对象的换基行为可以集中写成下表。设 \(e^1,\ldots,e^n\) 为 \(e\) 的对偶基，\(\varepsilon=e^1\wedge\cdots\wedge e^n\)，\(\varepsilon'=e'^1\wedge\cdots\wedge e'^n\)。
\[
	\begin{array}{c|c|c}
		\text{空间} & \text{由基产生的生成元} & \text{换到 }e'=eA\\ \hline
		\det V & e_1\wedge\cdots\wedge e_n & (\det A)\,e_1\wedge\cdots\wedge e_n\\
		\det(V^*) & \varepsilon & \varepsilon'=(\det A)^{-1}\varepsilon\\
		\mathcal D(V) & |\varepsilon| & |\varepsilon'|=|\det A|^{-1}|\varepsilon|
	\end{array}
\]
生成元的变换与同一对象的系数变换互为倒数。例如 \(\mu=m|\varepsilon|=m'|\varepsilon'|\)，便有 \(m'=|\det A|m\)，且 \(m=\mu(e)\)。这一点在把密度逐点粘成丛时会再次使用。

\ProseParagraphBreak

若 \(L:V\to W\) 是同维实向量空间之间的线性同构，则 \((L^*\mu)(e)=\mu(Le)\) 定义了密度的线性拉回，并与复合相容，因为 \(L(eA)=(Le)A\)。这里要求 \(L\) 可逆，才能直接把有序基送成有序基。对不可逆的同维线性映射，可以另用绝对行列式约定得到零密度，但不能据此把密度对光滑映射的拉回与普通微分形式的拉回不加区别地处理；行列式绝对值在变号处有不同的光滑性问题。

\begin{proposition}[内积的体积密度]
	\label{g1:prop:gram-volume-density}
	设 \(V\) 为有限维实向量空间，\(g\) 为正定内积。对有序基 \(v=(v_1,\ldots,v_n)\)，公式
	\[
		\mu_g(v_1,\ldots,v_n)=\sqrt{\det[g(v_i,v_j)]_{i,j=1}^n}
	\]
	定义一个正密度。若另给定向，则存在唯一在每组正向正交单位基上取值 \(1\) 的最高次形式 \(\operatorname{vol}_g\)，并有 \(|\operatorname{vol}_g|=\mu_g\)。对任意 \(n\) 个向量，同一平方根公式在它们线性相关时取零，在线性无关时就是相应平行多面体的体积。
\end{proposition}
\begin{proof}
	先说明根号内为正。逐个从基向量中减去沿先前向量的正交投影，再除以自身的正长度，得到正交单位基 \(u_1,\ldots,u_n\)；每一步的余向量不为零，因为原向量组线性无关。将 \(v\) 写成 \(uB\)，其内积矩阵就是 \(B^{\mathsf T}B\)，行列式为 \((\det B)^2>0\)。换基 \(v'=vA\) 后，矩阵变成 \(A^{\mathsf T}[g(v_i,v_j)]A\)，所以平方根乘以 \(|\det A|\)。这正是密度的变换规则，并证明了正性。

	选一组正向正交单位基 \(u\)，定义 \(\operatorname{vol}_g\) 在它上取值 \(1\)。其他正向正交单位基的变换矩阵 \(Q\) 满足 \(Q^{\mathsf T}Q=I\) 及 \(\det Q>0\)，因而 \(\det Q=1\)；所以定义不依赖所选正向正交单位基。最高次形式由在一组基上的值唯一决定，故唯一。对 \(v=uB\)，有 \(\operatorname{vol}_g(v)=\det B\)，而 \(\mu_g(v)=|\det B|\)，得到绝对值关系。若给定的 \(n\) 个向量线性相关，其坐标矩阵 \(B\) 奇异，Gram 行列式为零；否则 \(|\det B|\) 是正交坐标中平行多面体的体积。这也解释了该公式的几何意义。零维时采用空行列式等于 \(1\) 的约定。
\end{proof}

例如反射一组正交单位基中的一个向量，最高次体积形式的值由 \(1\) 变为 \(-1\)，体积密度的值仍为 \(1\)。在没有整体定向的几何问题里，长度和体积仍应是可测量的正量；密度保留的正是这一部分信息。

\subsection{不同标架下的分量及不变量}
\label{g1:subsec:1202:05}

前面的对象都先于基而存在，实际计算却需要系数。本小节把这两件事接起来。固定 \(n\) 维实向量空间的一组基 \(e\)，也称一个标架，并始终用 \(e'=eA\) 表示新基；矩阵的第 \(j\) 列是第 \(j\) 个新基向量在旧基中的分量。若 \(v\) 的两组分量列向量为 \([v]_e,[v]_{e'}\)，则 \(e[v]_e=eA[v]_{e'}\)，所以
\[
	[v]_{e'}=A^{-1}[v]_e.
\]
对偶基必须满足 \(e'^i(e'_j)=\delta^i_j\)，故 \(e'^i=\sum_a(A^{-1})^i{}_a e^a\)。但同一协变向量的系数是它在各基向量上的值，因而满足
\[
	[\alpha]_{e'}=A^{\mathsf T}[\alpha]_e.
\]
对偶基的变化和协变向量系数的变化是两件互逆的事。

\begin{proposition}[张量的分量变换]
	\label{g1:prop:tensor-component-change}
	设 \(V\) 为 \(n\) 维实向量空间，\(T\in T^r_s(V)\)，有序基满足 \(e'=eA\)。在各自的张量基中写
	\[
		T=\sum T^{i_1\cdots i_r}{}_{j_1\cdots j_s}
		\,e_{i_1}\otimes\cdots\otimes e_{i_r}\otimes e^{j_1}\otimes\cdots\otimes e^{j_s},
	\]
	以及对应的带撇表达，则
	\begin{equation}
		\label{g1:eq:tensor-change-components}
		T'^{i_1\cdots i_r}{}_{j_1\cdots j_s}
		=\sum_{\substack{a_1,\ldots,a_r\\ b_1,\ldots,b_s}}
		(A^{-1})^{i_1}{}_{a_1}\cdots(A^{-1})^{i_r}{}_{a_r}
		A^{b_1}{}_{j_1}\cdots A^{b_s}{}_{j_s}
		T^{a_1\cdots a_r}{}_{b_1\cdots b_s}.
	\end{equation}
	收缩的分量由一对上下指标取相同值后求和得到；这与换基相容。尤其所有输入都与张量配对之后得到的实数，与标架无关。
\end{proposition}
\begin{proof}
	以多线性函数的观点看，原系数恰为
	\[
		T^{i_1\cdots i_r}{}_{j_1\cdots j_s}
		=T(e^{i_1},\ldots,e^{i_r};e_{j_1},\ldots,e_{j_s}).
	\]
	新系数把这些输入全部换成新基及新对偶基。分别代入 \(e'_j=\sum_b e_bA^b{}_j\) 和 \(e'^i=\sum_a(A^{-1})^i{}_a e^a\)，再逐项使用多线性，即得到\eqref{g1:eq:tensor-change-components}。

	对纯张量收缩时，\(e^j(e_i)=\delta_i^j\)，故分量规则就是指定的上下指标相等后求和。若先换基再收缩，被收缩的一对指标带来的因子满足
	\[
		\sum_k(A^{-1})^k{}_a A^b{}_k=\delta^b_a,
	\]
	两矩阵相消，其余因子恰是低两阶张量的换基公式。因此先收缩或先换基给出同一张量。所有输入的分量也按相应规则变换，完全配对时每个因子都有逆因子与之抵消；所得数就是原来无基定义的多线性函数值。
\end{proof}

例如线性映射 \(L:V\to V\) 的矩阵满足 \(L'=A^{-1}LA\)，而双线性形式 \(g\) 的矩阵满足 \(G'=A^{\mathsf T}GA\)。它们的两种变换公式不同，是因为 \(L\) 有一个向量因子和一个对偶因子，\(g\) 有两个对偶因子。对任意 \(v,w\)，
\[
	[v]_{e'}^{\mathsf T}G'[w]_{e'}
	=[v]_e^{\mathsf T}G[w]_e;
\]
这是同一个数 \(g(v,w)\)。又如
\[
	L=\begin{pmatrix}2&0\\0&3\end{pmatrix},\qquad
	A=\begin{pmatrix}1&1\\0&1\end{pmatrix}
	\quad\Longrightarrow\quad
	L'=\begin{pmatrix}2&-1\\0&3\end{pmatrix}.
\]
新矩阵不再对角，但迹仍为 \(5\)，其线性映射也仍是原来那个映射。矩阵形状的改变与几何对象的改变须分开判断。

\begin{example}[极坐标中的度量与面积]
	\label{g1:exm:tensor-polar-components}
	在平面原点之外的一个角扇区上，取 \(x=r\cos\theta\)、\(y=r\sin\theta\)，其中 \(r>0\)。\cref{g1:exm:tangent-polar-coordinates}已经给出切基的变换；固定其中一点，写成矩阵就是
	\[
		(\partial_r,\partial_\theta)=(\partial_x,\partial_y)A,\qquad
		A=\begin{pmatrix}\cos\theta&-r\sin\theta\\ \sin\theta&r\cos\theta\end{pmatrix}.
	\]
	此处每个向量都在同一个切空间中。Euclid 内积在直角坐标基中的矩阵为 \(I\)，所以在极坐标基中为
	\[
		A^{\mathsf T}A=\begin{pmatrix}1&0\\0&r^2\end{pmatrix},\qquad
		g=\mathrm{d}r\otimes\mathrm{d}r+r^2\,\mathrm{d}\theta\otimes\mathrm{d}\theta.
	\]
	因此 \(a\partial_r+b\partial_\theta\) 的长度平方为 \(a^2+r^2b^2\)。出现 \(r^2\) 是因为 \(\partial_\theta\) 表示每单位角度的速度，其长度为 \(r\)，并不表示 Euclid 内积换成了另一种内积。

	对偶变换也可直接求导：
	\[
		\mathrm{d}x=\cos\theta\,\mathrm{d}r-r\sin\theta\,\mathrm{d}\theta,\qquad
		\mathrm{d}y=\sin\theta\,\mathrm{d}r+r\cos\theta\,\mathrm{d}\theta.
	\]
	使用楔积的双线性和一次形式交换变号，得到
	\[
		\mathrm{d}x\wedge\mathrm{d}y
		=r\,\mathrm{d}r\wedge\mathrm{d}\theta,\qquad
		\mu_g=|\mathrm{d}x\wedge\mathrm{d}y|
		=r\,|\mathrm{d}r\wedge\mathrm{d}\theta|.
	\]
	度量的 Gram 行列式为 \(r^2\)，平方根也给出密度系数 \(r\)，与行列式换基结果一致。原点处 \(A\) 不可逆，所以这些公式是在极坐标有效的区域内比较两组标架；不能用退化的极坐标基描述原点处的张量。
\end{example}
% 纲要条目排为编号小节标题，并进入目录与书签；说明段仍保留于纲要注释中。

% !TeX root = ../../Book1.tex
% 纲要标题：### 12.3 光滑单位分解与延拓
% 纲要位置：几何学计划纲要.md，第 587 行。
% 正文状态：初稿；尚未正式审核。

% 纲要内容开始
% 
% * 光滑截断函数和局部有限细化
% * 单位分解的完整构造
% * 截面延拓与相对版本
% * 丛度量、正交补和短正合丛列的分裂
% 
% 纲要内容结束

\section{光滑单位分解与延拓}
\label{g1:sec:1203}

在每张局部标架中，把纤维坐标的通常内积搬回来，就得到一个局部丛度量。可是两张标架之间的过渡矩阵通常不是正交矩阵，两个局部内积因而未必相同，不能直接粘合。截面延拓也有相似的困难：局部公式可以很好地描述一块区域，却可能在接近定义域边界时发散。我们需要一族光滑权重，让每个局部对象只在自己的定义域内起作用，并使这些权重在各点的总和为 \(1\)。

\ProseParagraphBreak

本节中的流形均为有限维、Hausdorff 且第二可数的光滑流形，允许有边界；向量丛均为有限秩实光滑向量丛。边界处的光滑性沿用\cref{g1:def:manifold-with-boundary}的局部延拓定义。构造首先在 Euclid 空间及半空间中完成，再借助局部有限的坐标覆盖搬到流形上；局部有限性保证以下无穷和在每一点附近其实都是有限和。

\subsection{光滑截断函数和局部有限细化}
\label{g1:subsec:1203:01}

一个函数在开集 \(U\) 外等于零，并不表示它的支集包含于 \(U\)。按照\cref{g1:def:support-cutoff}，支集还包括非零点的所有极限点：
\[
	\operatorname{supp}f=\overline{\{p:f(p)\ne0\}}.
\]
例如 \(f(t)=\max\{1-|t|,0\}\) 在 \((-1,1)\) 外为零，支集却是 \([-1,1]\)。我们要把支集也留在允许的开集内部，并把折角换成任意阶可微的过渡。

\begin{lemma}[光滑截断函数]\label{g1:lem:smooth-bump-function}
	函数
	\[
		\eta(t)=
		\begin{cases}
			e^{-1/t},&t>0,\\
			0,&t\le0
		\end{cases}
	\]
	在 \(\mathbb R\) 上光滑，且 \(\eta^{(k)}(0)=0\) 对每个非负整数 \(k\) 成立。由它可构造光滑函数 \(\vartheta:\mathbb R\to[0,1]\)，满足 \(\vartheta=0\) 于 \((-\infty,0]\)、\(\vartheta=1\) 于 \([1,\infty)\)。因此，若 \(0<r<s<R\)，在 Euclid 空间或半空间 \(H^m=\{x^m\ge0\}\) 内，存在光滑函数，等于 \(1\) 于半径 \(r\) 的闭球或相对闭半球，支集包含于半径 \(s\) 的相对闭球，因而紧含于半径 \(R\) 的相对开球。
\end{lemma}
\begin{proof}
	首先要检查拼接点 \(0\)，仅知道正半轴上的公式光滑是不够的。对 \(t>0\) 反复求导得到
	\[
		\eta^{(k)}(t)=e^{-1/t}P_k(1/t),
	\]
	其中 \(P_k\) 是多项式。确实 \(P_0=1\)，写 \(u=1/t\)，下一次求导给出 \(P_{k+1}(u)=u^2(P_k(u)-P_k'(u))\)。对任意非负整数 \(N\)，有 \(u^N e^{-u}\to0\)（\(u\to+\infty\)）：例如从指数级数取第 \(N+1\) 项，便有 \(u^N e^{-u}\le (N+1)!/u\)。

	把每个 \(e^{-1/t}P_k(1/t)\) 在 \(t\le0\) 补成零，所得函数在 \(0\) 连续。而它在 \(0\) 的右侧差商又是 \(e^{-1/t}P_k(1/t)/t\)，仍趋于零；左侧差商恒为零。因此这个延拓的导数在 \(0\) 为零，在两侧则是下一阶公式。逐阶归纳即得 \(\eta\in C^\infty(\mathbb R)\)，且所有阶导数在 \(0\) 消失。

	令
	\[
		\vartheta(t)=\frac{\eta(t)}{\eta(t)+\eta(1-t)}.
	\]
	分母处处为正：若第一项为零，则 \(t\le0\)，此时 \(1-t>0\)。其余性质由非负性及 \(\eta\) 的零点集直接得到。对中心 \(a\) 定义
	\[
		b_{a,r,s}(x)=
		\vartheta\!\left(\frac{s^2-|x-a|^2}{s^2-r^2}\right).
	\]
	它在 \(|x-a|\le r\) 时为 \(1\)，在 \(|x-a|\ge s\) 时为零，支集包含于闭球 \(\overline B_s(a)\)。限制到 \(H^m\) 后仍是光滑函数，因为原来的公式就是其 Euclid 光滑延拓；支集按半空间的相对拓扑计算，包含于 \(\overline B_s(a)\cap H^m\)。
\end{proof}

若坐标图中半径 \(R\) 的相对闭球仍留在坐标域内，就可把上述函数搬到流形，并在坐标域外补零。补零为何光滑？支集已留在半径 \(s<R\) 的紧集内，每个坐标域外的点都有一个与该支集不交的邻域，所得函数在这个邻域恒为零。边界坐标图使用相对球，论证完全相同。零维流形的每个单点都是开集，其示性函数就是所需的光滑截断函数。

\ProseParagraphBreak

许多局部函数相加时，另一个要求是控制它们在同一邻域内出现的数目。回忆\cref{g1:def:locally-finite-refinement}：集合族 \(\{C_i\}\) 局部有限，是指每一点都有一个开邻域，只与有限个 \(C_i\) 相交；逐点有限则只要求每一点属于有限个 \(C_i\)。后一个条件不能控制邻近点发生的变化。

\begin{example}\label{g1:exm:pointwise-finite-smooth-sum}
	在 \(\mathbb R\) 中令 \(c_n=2^{-n}\)、\(r_n=2^{-n-2}\)，并用前述构造选光滑函数 \(b_n\)，使 \(0\le b_n\le1\)、\(b_n(c_n)=1\)，且
	\[
		\operatorname{supp}b_n\subseteq[c_n-r_n,c_n+r_n].
	\]
	这些闭区间两两不交，所以每一点最多落在一个支集中，\(\sum_n b_n(x)\) 在每一点都只有至多一个非零项。然而这个和在 \(0\) 等于零，在每个 \(c_n\) 处等于 \(1\)，因而在 \(0\) 不连续。问题正是任何原点邻域都会遇到无穷多个支集。\cref{g1:fig:locally-finite-bumps}把这种聚集与沿直线向远处移动的支集作了对照。
\end{example}

% !TeX root = ../Book1.tex
% 支集与峰值的示意图；曲线形状不表示截断函数的精确公式。
\begin{figure}[htbp]
	\centering
	\begin{tikzpicture}[font=\small]
		\begin{scope}[x=0.95cm,y=1.25cm]
			\node[align=center] at (2.3,1.65) {局部有限};
			\fill[orange!15] (1.92,-0.08) rectangle (2.68,1.12);
			\draw[orange!75!black,dashed] (1.92,-0.08)--(1.92,1.12);
			\draw[orange!75!black,dashed] (2.68,-0.08)--(2.68,1.12);
			\draw[->,GeometryGray] (-0.1,0)--(4.85,0);
			\foreach \c/\lab in {0.7/-1,2.3/0,3.9/1}{
				\draw[GeometryBlue,thick] ({\c-0.65},0) .. controls ({\c-0.40},0) and ({\c-0.28},1) .. (\c,1) .. controls ({\c+0.28},1) and ({\c+0.40},0) .. ({\c+0.65},0);
				\draw[GeometryBlue,very thick] ({\c-0.65},-0.09)--({\c+0.65},-0.09);
				\node[below] at (\c,-0.1) {\(\lab\)};
			}
			\node at (0.0,0.52) {\(\cdots\)};
			\node at (4.65,0.52) {\(\cdots\)};
			\node[orange!75!black] at (2.3,1.25) {\(N\)};
			\node[align=center,text width=4.8cm] at (2.3,-0.96) {固定邻域 \(N\) 只遇有限个支集；\\更远的支集不会返回这里。};
		\end{scope}
		\begin{scope}[xshift=5.7cm,x=6.2cm,y=1.25cm]
			\node[align=center] at (0.35,1.65) {逐点有限，局部并非有限};
			\draw[->,GeometryGray] (-0.025,0)--(0.75,0);
			\draw[GeometryGray,densely dashed] (0,1)--(0.65,1);
			\node[left] at (0,1) {\(1\)};
			\foreach \n in {1,...,6}{
				\pgfmathsetmacro{\c}{pow(2,-\n)}
				\pgfmathsetmacro{\r}{\c/4}
				\draw[GeometryBlue,thick] ({\c-\r},0) .. controls ({\c-0.65*\r},0) and ({\c-0.45*\r},1) .. (\c,1) .. controls ({\c+0.45*\r},1) and ({\c+0.65*\r},0) .. ({\c+\r},0);
				\fill[GeometryBlue] (\c,1) circle[radius=1.2pt];
			}
			\fill (0,0) circle[radius=1.5pt];
			\node[below] at (0,0) {\(0\)};
			\node[below] at (0.25,0) {\(\tfrac14\)};
			\node[below] at (0.5,0) {\(\tfrac12\)};
			\node[orange!75!black,align=center] at (0.24,1.25) {\(c_n\longrightarrow0\)};
			\node[align=center,text width=4.9cm] at (0.35,-0.96) {各支集两两不交，却聚集到 \(0\)；\\峰值为 \(1\)，原点处的和为 \(0\)。};
		\end{scope}
	\end{tikzpicture}
	\caption{左侧支集沿整数方向排列；右侧支集为以 \(c_n=2^{-n}\) 为中心、半径 \(2^{-n-2}\) 的区间。右侧每点至多属于一个支集，却无法在原点附近把求和固定为有限和。图中曲线仅示意支集与峰值。}
	\label{g1:fig:locally-finite-bumps}
\end{figure}


下面的细化既给出覆盖，又为每个截断函数预留一圈支集边缘。记号 \(V\Subset U\) 表示 \(\overline V\) 紧且 \(\overline V\subseteq U\)，闭包一律在 \(M\) 中计算。

\begin{lemma}[相对紧坐标细化]\label{g1:lem:precompact-coordinate-refinement}
	设 \(M\) 为有限维、Hausdorff 且第二可数的光滑流形，允许有边界，\(\{O_\alpha\}_{\alpha\in I}\) 为任意开覆盖。存在至多可数族坐标邻域对 \(V_i\Subset U_i\) 及指标选择 \(\alpha(i)\)，使
	\[
		M=\bigcup_iV_i,\qquad
		\overline{U_i}\text{ 紧},\qquad
		\overline{U_i}\subseteq O_{\alpha(i)},
	\]
	且 \(\{U_i\}\) 局部有限。每对 \(V_i,U_i\) 可在同一坐标图中选为同心的相对开球；存在光滑函数 \(b_i:M\to[0,1]\)，满足 \(b_i=1\) 于 \(\overline{V_i}\)、\(\operatorname{supp}b_i\Subset U_i\)。
\end{lemma}
\begin{proof}
	\cref{g1:thm:compact-exhaustion,g1:thm:lch-paracompact}分别给出了紧穷竭与壳层细化。这里恢复构造，并在每一层加入所需的两层坐标球。这样从第十一章开始阅读的读者也能看到第二可数性与局部紧性各用在何处。

	先构造紧集列 \(K_n\subseteq\operatorname{int}K_{n+1}\)，使其并为 \(M\)。流形的坐标球或半球可以缩小为紧闭球或紧闭半球，所以每一点的每个邻域都包含一个相对紧开邻域。从一个可数拓扑基中仅保留相对紧的成员，仍得到基：在上述较小邻域中再选基元素，其闭包就是紧闭包的闭子集。把非空成员列为 \(B_1,B_2,\ldots\)，有限时允许重复。令 \(K_1=\overline B_1\)。若 \(K_n\) 已选，由紧性用有限个 \(B_j\) 覆盖 \(K_n\cup\overline B_{n+1}\)，令 \(K_{n+1}\) 为这些 \(B_j\) 的闭包之并。于是
	\[
		K_n\cup\overline B_{n+1}\subseteq\operatorname{int}K_{n+1},
		\qquad
		M=\bigcup_{n\ge1}K_n=\bigcup_{n\ge1}\operatorname{int}K_n.
	\]
	各 \(K_n\) 紧且在 Hausdorff 空间中闭；后一事实可由外点与紧集各点的分离邻域、再取有限子覆盖得到。空流形无需作任何选择。

	约定 \(K_{-1}=K_0=\varnothing\)，定义需要覆盖的紧壳层和允许选择邻域的较大开区域
	\[
		S_n=K_n\setminus\operatorname{int}K_{n-1},
		\qquad
		G_n=\operatorname{int}K_{n+1}\setminus K_{n-2}.
	\]
	有 \(S_n\subseteq G_n\)：外侧由内部包含关系保证，内侧则因为 \(K_{n-2}\subseteq\operatorname{int}K_{n-1}\)。每一点首次进入某个 \(K_n\) 时属于 \(S_n\)，所以各壳层覆盖 \(M\)。

	对每个 \(p\in S_n\)，选择含它的 \(O_{\alpha(p)}\)，在开邻域 \(G_n\cap O_{\alpha(p)}\) 中取一张坐标图。选以 \(p\) 的坐标为中心的三层相对球，使内层 \(V_p\) 的闭包在中层 \(U_p\) 内，中层的闭包又紧含于该坐标图；特别地，\(\overline{U_p}\subseteq G_n\cap O_{\alpha(p)}\)。由壳层的紧性，保留有限个内层球 \(V_{n,1},\ldots,V_{n,k_n}\) 覆盖 \(S_n\)，并同时保留其外层球 \(U_{n,j}\)。所有层合起来至多可数，内层球覆盖 \(M\)。

	还须检查局部有限性；仅仅每层有限并不够。给定 \(p\)，取 \(m\) 使 \(p\in\operatorname{int}K_m\)。当 \(n\ge m+2\) 时，这一层的全部 \(\overline{U_{n,j}}\) 位于 \(G_n\subseteq M\setminus K_{n-2}\)，而 \(K_m\subseteq K_{n-2}\)，故它们与 \(\operatorname{int}K_m\) 不交。这个固定邻域只能遇到前 \(m+1\) 层的有限个外层球，局部有限性得证。最后，在每对内外球之间选一个中间半径，由\cref{g1:lem:smooth-bump-function}构造 \(b_i\)，再补零到 \(M\)。它的支集位于中间闭球内，因而紧含于 \(U_i\)。
\end{proof}

\subsection{单位分解的完整构造}
\label{g1:subsec:1203:02}

现在已有非负函数 \(b_i\)，每一点至少落在一个 \(b_i=1\) 的区域内，各点附近又只出现有限项。只需除以这些函数的总和，就能得到总和为 \(1\) 的权重。这里不能省略支集控制：局部对象只在 \(O_\alpha\) 上定义时，权重的非零点及其极限点都应留在 \(O_\alpha\) 内。

\begin{definition}[光滑单位分解]\label{g1:def:smooth-partition-unity}
	设 \(\{O_\alpha\}_{\alpha\in I}\) 是光滑流形 \(M\) 的开覆盖。一族光滑函数 \(\varphi_\alpha:M\to[0,1]\) 若满足
	\[
		\operatorname{supp}\varphi_\alpha\subseteq O_\alpha,\qquad
		\{\operatorname{supp}\varphi_\alpha\}_{\alpha\in I}\text{ 局部有限},\qquad
		\sum_{\alpha\in I}\varphi_\alpha(p)=1,
	\]
	则称为从属于这个覆盖的\term[guanghuadanweifenjie]{光滑单位分解}[smooth partition of unity]。允许某些 \(\varphi_\alpha\) 恒为零。这是\cref{g1:def:partition-of-unity}中连续单位分解的光滑版本。
\end{definition}

\begin{theorem}[光滑单位分解的存在]\label{g1:thm:smooth-partition-unity}
	有限维、Hausdorff 且第二可数的光滑流形 \(M\)，包括带边界流形，其每个开覆盖都有从属于该覆盖的光滑单位分解。构造可以先得到一个从属于可数细化、每项紧支且支集族局部有限的光滑单位分解，再按原覆盖指标归组；归组后的各项不一定紧支。若 \(M\) 紧，可以只使用有限个非零项。
\end{theorem}
\begin{proof}
	给定覆盖 \(\{O_\alpha\}\)，用\cref{g1:lem:precompact-coordinate-refinement}选取 \(V_i,U_i,b_i\)。由于 \(\operatorname{supp}b_i\subseteq U_i\)，支集族局部有限。因此
	\[
		b=\sum_i b_i
	\]
	在每一点的某个邻域内都是有限个光滑函数之和，从而光滑。内层 \(V_i\) 覆盖 \(M\)，故每个点处至少有一个 \(b_i=1\)，得到 \(b\ge1\)。定义
	\[
		\psi_i=\frac{b_i}{b}.
	\]
	这些函数光滑、非负、总和为 \(1\)，且 \(\psi_i\) 与 \(b_i\) 有相同的非零点集，所以
	\[
		\operatorname{supp}\psi_i=\operatorname{supp}b_i\Subset U_i\subseteq O_{\alpha(i)}.
	\]
	这给出了细化上的单位分解。

	若需要与原覆盖同指标，令
	\[
		\varphi_\alpha=\sum_{\alpha(i)=\alpha}\psi_i.
	\]
	每个和局部只有有限项，因此光滑。说明支集包含关系时，要防止无限并增加边界点。局部有限闭集族的任意子族之并仍闭：在并集外的一个点附近，只遇到有限个闭集，去掉这有限个闭集后仍留下一个开邻域。将这点用于 \(\{\operatorname{supp}\psi_i\}\)，得到
	\[
		\operatorname{supp}\varphi_\alpha
		\subseteq\bigcup_{\alpha(i)=\alpha}\operatorname{supp}\psi_i
		\subseteq O_\alpha.
	\]
	前一个包含成立，是因为右侧是包含所有非零点的闭集。一个只与有限个 \(\operatorname{supp}\psi_i\) 相交的邻域，也只可能遇到相应有限个 \(\operatorname{supp}\varphi_\alpha\)，所以归组后支集族仍局部有限。各点处重排的只是有限和，故 \(\sum_\alpha\varphi_\alpha=1\)。非负性又保证各项不超过 \(1\)。若某个原指标没有分配到成员，就取零函数。

	当 \(M\) 紧时，先从覆盖 \(M\) 的内层球 \(V_i\) 中选有限子覆盖，只保留相应的 \(b_i\)，再作同样的归一化即可。一般情形下，按原指标归组会把无穷多个紧支集并起来，这个并虽然闭，却未必紧。
\end{proof}

局部有限还保证可以逐项求任意阶导数，因为在计算点的一个固定邻域内，本来就只有有限个函数。这里既不需要估计导数级数，也不涉及交换无限求和与极限；\cref{g1:exm:pointwise-finite-smooth-sum}说明仅有逐点有限时，连连续性都可能失败。

\begin{corollary}[闭集的光滑截断]\label{g1:cor:smooth-cutoff-closed}
	设 \(M\) 为有限维、Hausdorff 且第二可数的光滑流形，允许有边界，\(A\subseteq M\) 闭，\(U\subseteq M\) 开且 \(A\subseteq U\)。存在光滑函数 \(\chi:M\to[0,1]\)，使
	\[
		\chi=1\text{ 于 }A\text{ 的某个开邻域},\qquad
		\operatorname{supp}\chi\subseteq U.
	\]
	若 \(A\) 紧，还可要求 \(\operatorname{supp}\chi\) 紧。
\end{corollary}
\begin{proof}
	对开覆盖 \(\{U,M\setminus A\}\) 取单位分解 \(\chi+\rho=1\)，其中 \(\operatorname{supp}\chi\subseteq U\)、\(\operatorname{supp}\rho\subseteq M\setminus A\)。开集 \(M\setminus\operatorname{supp}\rho\) 包含 \(A\)，在这里 \(\rho=0\)，所以 \(\chi=1\)。

	若 \(A\) 紧，在 \(U\) 中给每个 \(p\in A\) 选相对紧坐标邻域，再取有限个覆盖 \(A\)。它们的并 \(W\) 满足 \(A\subseteq W\) 且 \(\overline W\) 紧、\(\overline W\subseteq U\)。把前面的构造用于 \(A\subseteq W\)，所得支集闭且包含于紧集 \(\overline W\)，因而紧。若 \(A=\varnothing\)，直接取 \(\chi=0\)。
\end{proof}

“支集包含于 \(U\)”与“支集紧”是两件事。即使 \(U\) 非紧，只要 \(A\) 紧，仍可选紧支截断；反过来，若闭集 \(A\) 非紧，任何在 \(A\) 上等于 \(1\) 的函数都不可能紧支，因为其支集包含 \(A\)。例如 \(A=\mathbb Z\subset U=\mathbb R\)，截断函数存在，却不能紧支。

\subsection{截面延拓与相对版本}
\label{g1:subsec:1203:03}

回忆\cref{g1:def:bundle-section}，丛 \(E\to M\) 的截面 \(s\) 在每一点 \(p\) 取值于对应纤维 \(E_p\)，光滑性可在局部标架中逐分量判断。截面的支集同样定义为
\[
	\operatorname{supp}s=\overline{\{p:s(p)\ne0_p\}}.
\]
对开集 \(U\) 上的光滑截面 \(s\)，若光滑函数 \(\chi\) 满足 \(\operatorname{supp}\chi\subseteq U\)，就可以令 \(\chi s\) 在 \(U\) 外等于零。它在 \(U\) 内逐分量光滑，在 \(U\) 外的每个点附近恒为零，因此是整体光滑截面。这里不要求 \(s\) 在接近 \(U\) 的边界时有界。

\begin{definition}[闭子集上的光滑截面]\label{g1:def:smooth-section-closed-set}
	设 \(E\to M\) 为光滑向量丛，\(A\subseteq M\) 为闭集。给定一个赋值 \(a(p)\in E_p\)（\(p\in A\)）。局部光滑延拓的条件是：每个 \(p\in A\) 都有 \(M\) 中的开邻域 \(O_p\) 和光滑截面 \(s_p\in\Gamma(E|_{O_p})\)，满足
	\[
		s_p(q)=a(q)\qquad(q\in A\cap O_p).
	\]
	若满足这一条件，就称这一赋值为 \(A\) 上的光滑截面。
\end{definition}

这个定义允许 \(A\) 不是子流形，只要求所给纤维值在各点附近可以光滑延拓。不同 \(s_p\) 在 \(A\) 外可以不同；当前数据没有指定法向导数或更高阶导数。

\begin{theorem}[截面延拓]\label{g1:thm:section-extension}
	设 \(M\) 为有限维、Hausdorff 且第二可数的光滑流形，允许有边界，\(E\to M\) 为有限秩实光滑向量丛，\(A\subseteq M\) 闭，\(U\supseteq A\) 开。
	\begin{enumerate}
		\item 若 \(s\in\Gamma(E|_U)\)，则存在整体光滑截面 \(\widetilde s\)，在 \(A\) 的某开邻域与 \(s\) 相同，且 \(\operatorname{supp}\widetilde s\subseteq U\)。
		\item 若 \(a\) 为\cref{g1:def:smooth-section-closed-set}意义下 \(A\) 上的光滑截面，则存在整体光滑截面 \(\widetilde a\)，满足 \(\widetilde a|_A=a\) 且 \(\operatorname{supp}\widetilde a\subseteq U\)。
	\end{enumerate}
	若 \(A\) 紧，两项都可要求所得截面紧支。
\end{theorem}
\begin{proof}
	第一项中，由\cref{g1:cor:smooth-cutoff-closed}选取在 \(A\) 附近等于 \(1\)、支集包含于 \(U\) 的 \(\chi\)。令
	\[
		\widetilde s(p)=
		\begin{cases}
			\chi(p)s(p),&p\in U,\\
			0_p,&p\notin U.
		\end{cases}
	\]
	前面的零延拓论证说明其光滑，支集包含于 \(\operatorname{supp}\chi\)，并且在 \(\chi=1\) 的那个邻域与 \(s\) 相同。紧性结论也由 \(\chi\) 的紧支性得到。

	第二项先粘合各个局部延拓，再控制支集。为每个 \(p\in A\) 取定义中的 \(O_p,s_p\)，并加入开集 \(O_*=M\setminus A\) 及其零截面 \(s_*=0\)，得到 \(M\) 的开覆盖。取从属单位分解 \(\{\varphi_p\}\cup\{\varphi_*\}\)。每个乘积 \(\varphi_p s_p\) 都可在 \(O_p\) 外光滑补零，支集仍包含于 \(\operatorname{supp}\varphi_p\)。因此
	\[
		t=\sum_{p\in A}\varphi_p s_p+\varphi_*s_*
	\]
	是整体光滑截面：每一点附近只需求有限项相加。若 \(q\in A\)，则 \(\varphi_*(q)=0\)，而每个在 \(q\) 处权重非零的 \(s_p\) 都满足 \(s_p(q)=a(q)\)，从而
	\[
		t(q)=\left(\sum_{p\in A}\varphi_p(q)\right)a(q)=a(q).
	\]
	最后再用第一项中的截断函数，取 \(\widetilde a=\chi t\)，即可控制支集并保持 \(A\) 上的值；\(A\) 紧时选紧支的 \(\chi\)。
\end{proof}

第一项保留的是一个邻域上的整条截面，第二项保留的是闭集上指定的值；不能从第二项自动推断它同时保留各个局部延拓的全部导数。闭性也有实际作用。取平凡线丛 \(\mathbb R\times\mathbb R\)，在非闭集 \(A=(0,1)\) 上给定 \(a(x)=1/x\)。它甚至已在开集 \(A\) 上光滑，却不可能延拓为 \(\mathbb R\) 上的连续截面，因为 \(x\downarrow0\) 时发散。

\ProseParagraphBreak

有时已经有一个整体截面 \(s_0\)，只希望在某闭集附近把它换成指定截面，并在另一闭集附近保持原样。两块区域相遇时需要相容；仅在相遇点取值相同，未必足以要求双方都在邻域内保持不变。

\begin{proposition}[相对截面延拓]\label{g1:prop:relative-section-extension}
	设 \(M\) 为 Hausdorff、第二可数的有限维光滑流形，允许带边界，\(E\to M\) 为有限秩实光滑向量丛，\(A,B\subseteq M\) 为闭集，\(s_0\in\Gamma(E)\)。这些流形与丛条件保证可以应用\cref{g1:thm:section-extension}。
	\begin{enumerate}
		\item 设 \(U\supseteq A\) 开，\(s\in\Gamma(E|_U)\)。若有开集 \(W\supseteq A\cap B\)，使 \(s=s_0\) 于 \(U\cap W\)，则存在 \(\widetilde s\in\Gamma(E)\)，在 \(A\) 的某开邻域等于 \(s\)，在 \(B\) 的某开邻域等于 \(s_0\)，且 \(\operatorname{supp}(\widetilde s-s_0)\subseteq U\)。
		\item 设 \(a\) 是闭集 \(A\) 上的光滑截面。若有开集 \(W\supseteq B\)，使 \(a=s_0|_A\) 于 \(A\cap W\)，则存在整体延拓 \(\widetilde a|_A=a\)，并且 \(\widetilde a=s_0\) 于 \(B\) 的某开邻域。
	\end{enumerate}
	若 \(A\) 紧，还可在两项中要求对 \(s_0\) 的改变量紧支。
\end{proposition}
\begin{proof}
	第一项中，闭集 \(B\setminus W\) 与 \(A\) 不交，所以
	\[
		O=U\setminus(B\setminus W)
	\]
	是包含 \(A\) 的开集。选在 \(A\) 附近为 \(1\)、支集包含于 \(O\) 的截断函数 \(\chi\)，把 \(\chi(s-s_0)\) 从 \(U\) 补零，令
	\[
		\widetilde s=s_0+\chi(s-s_0).
	\]
	在 \(A\) 附近有 \(\chi=1\)，故 \(\widetilde s=s\)。在 \(W\) 上改变量为零，因为 \(U\cap W\) 内两截面相等，\(U\) 外则已补零。在 \(M\setminus\operatorname{supp}\chi\) 上改变量也为零。这两个开集的并包含 \(B\)：属于 \(B\cap W\) 的点在前者中，属于 \(B\setminus W\) 的点不在支集内，因而在后者中。支集控制由构造得到。

	第二项中，先用\cref{g1:thm:section-extension}取 \(a\) 的一个整体延拓 \(t\)。闭集 \(C=A\setminus W\) 与 \(B\) 不交，故可选在 \(C\) 附近为 \(1\)、支集包含于 \(M\setminus B\) 的光滑截断函数 \(\chi\)。令
	\[
		\widetilde a=s_0+\chi(t-s_0).
	\]
	在 \(A\setminus W\) 上 \(\chi=1\)，在 \(A\cap W\) 上则有 \(t=a=s_0\)，所以两部分都给出 \(\widetilde a=a\)。开集 \(M\setminus\operatorname{supp}\chi\) 包含 \(B\)，其上 \(\widetilde a=s_0\)。若 \(A\) 紧，第一项中的 \(A\) 与第二项中的 \(C\) 都紧，可选紧支的 \(\chi\)，于是改变量紧支。
\end{proof}

例如在平凡线丛上取 \(A=[0,1]\)、\(B=\{0\}\)、\(a(x)=x\)、\(s_0=0\)。虽然 \(a\) 与 \(s_0\) 在 \(A\cap B\) 上相等，却不能找一个同时延拓 \(a\) 并在 \(0\) 附近恒为零的截面。因此相对版本中的邻域相容条件不能只改成相交点上的相等。当 \(A,B\) 不交时，第一项的相容条件自动满足，可取 \(W=\varnothing\)。

\subsection{丛度量、正交补和短正合丛列的分裂}
\label{g1:subsec:1203:04}

截面的局部公式相加前，必须先确认它们在同一点属于同一个向量空间。丛的局部标架本身不提供这种相加规则：不同标架中的同一列数字，通常表示不同的向量。单位分解粘合的是已解释为纤维对象的局部截面、双线性型或丛同态，再通过标架检验其光滑性。

\begin{definition}[丛度量]\label{g1:def:bundle-metric}
	设 \(E\to M\) 为有限秩实光滑向量丛。给每个纤维 \(E_p\) 指定正定对称双线性型 \(h_p\)，若在任一光滑局部标架 \(e_1,\ldots,e_r\) 中，所有函数 \(h(e_a,e_b)\) 都光滑，则称 \(h\) 为 \(E\) 上的\term[congduliang]{丛度量}[bundle metric]。
\end{definition}

只需在一组覆盖 \(M\) 的局部标架中检查光滑性。事实上，按照\cref{g1:def:bundle-section}的行标架约定 \(e_j=e_i g_{ij}\)，度量矩阵满足
\[
	H_j=g_{ij}^{\mathsf T}H_i g_{ij}.
\]
右侧是光滑矩阵的乘积，因而光滑。这里说的是各纤维内的内积；在尚未另作定义之前，它不是底空间 \(M\) 上两点间的距离。

\begin{theorem}[丛度量的存在]\label{g1:thm:bundle-metric-existence}
	设 \(M\) 为有限维、Hausdorff 且第二可数的光滑流形，允许有边界，\(E\to M\) 为有限秩实光滑向量丛，则 \(E\) 存在光滑丛度量。
\end{theorem}
\begin{proof}
	取局部标架覆盖 \(\{O_\alpha\}\)。在每个 \(O_\alpha\) 上，把该标架声明为正交归一，得到局部度量
	\[
		h_\alpha\!\left(\sum_a v^a e_{\alpha,a},\sum_a w^a e_{\alpha,a}\right)
		=\sum_a v^a w^a.
	\]
	它在别的标架中仍有光滑矩阵，因为过渡矩阵光滑。取从属于 \(\{O_\alpha\}\) 的光滑单位分解，定义
	\[
		h=\sum_\alpha\varphi_\alpha h_\alpha,
	\]
	每个加权局部双线性型在定义域外补零。其支集控制保证补零光滑；局部有限性保证每一矩阵分量在各点附近都是有限个光滑函数之和。

	所得 \(h_p\) 对称且双线性。若 \(v\in E_p\) 非零，则每个有定义的 \(h_{\alpha,p}(v,v)>0\)，而权重非负且至少一个为正，故
	\[
		h_p(v,v)=\sum_\alpha\varphi_\alpha(p)h_{\alpha,p}(v,v)>0.
	\]
	这给出正定性，完成构造。
\end{proof}

正定性能够保留，是因为正定对称双线性型构成凸集。给定两个丛度量 \(h_0,h_1\)，对 \(0\le t\le1\)，\((1-t)h_0+th_1\) 仍是丛度量。存在性并未选出一个唯一度量，局部标架与单位分解的选择通常都会改变结果。

\begin{proposition}[正交补子丛]\label{g1:prop:orthogonal-complement-bundle}
	设 \(E\to M\) 为有限秩实光滑向量丛，\(h\) 为光滑丛度量，\(F\subseteq E\) 为秩 \(k\) 的光滑子丛。则逐纤维的正交补
	\[
		F^\perp_p=\{v\in E_p:h_p(v,w)=0\text{ 对所有 }w\in F_p\}
	\]
	构成光滑子丛，并有光滑直和分解 \(E=F\oplus F^\perp\)。到 \(F\) 的正交投影是光滑丛同态，\(h\) 在两个子丛上的限制仍为光滑丛度量。
\end{proposition}
\begin{proof}
	每个纤维中的线性代数已给出向量空间直和；要补足的是正交补和投影如何随底点光滑变化。由\cref{g1:prop:subbundle-adapted-frame}，在一点附近取环境丛标架 \(e_1,\ldots,e_r\)，使前 \(k\) 项张成 \(F\)。令
	\[
		A_{ab}=h(e_a,e_b)\quad(1\le a,b\le k).
	\]
	矩阵 \(A\) 光滑且正定，故可逆，\(A^{-1}=\operatorname{adj}(A)/\det A\) 也是光滑矩阵。对纤维向量 \(v\) 定义
	\[
		P(v)=\sum_{a,b=1}^k e_a(A^{-1})_{ab}\,h(e_b,v).
	\]
	对每个 \(c\le k\)，有
	\[
		h(e_c,P(v))=\sum_{a,b=1}^k A_{ca}(A^{-1})_{ab}h(e_b,v)
		=h(e_c,v).
	\]
	所以 \(v-P(v)\) 与 \(F\) 正交；若 \(v\in F\)，则公式又给出 \(P(v)=v\)。这正是纤维上的唯一正交投影，因而不同局部标架得到的公式自动相同，可以粘合为整体光滑丛同态。

	对 \(j=k+1,\ldots,r\) 令 \(f_j=e_j-P(e_j)\)。它们光滑且落在 \(F^\perp\) 中；若其线性组合为零，在商空间 \(E_p/F_p\) 中取像，得到 \(e_{k+1},\ldots,e_r\) 的对应线性组合为零，从而全部系数为零。它们共有 \(r-k\) 个，正好构成 \(F^\perp\) 的局部标架，证明了子丛的光滑性。映射 \(F\oplus F^\perp\to E\) 的逆为 \(v\mapsto(Pv,v-Pv)\)，两者都光滑。度量限制的矩阵光滑，正定性逐纤维保持；\(k=0\) 或 \(k=r\) 时按零子丛解释同样成立。
\end{proof}

有了正交补，就能把逐纤维的短正合列拆成整体直和。回忆\cref{g1:def:short-exact-bundle-sequence}，短正合要求第一张箭头在每个纤维单射、第二张箭头满射，且前者的像等于后者的核；这些条件本身还没有指定一个与核互补的子丛。

\begin{theorem}[短正合丛列的光滑分裂]\label{g1:thm:exact-bundle-splitting}
	设 \(M\) 为有限维、Hausdorff 且第二可数的光滑流形，允许有边界，且
	\[
		0\longrightarrow E'\xrightarrow{\ i\ }E\xrightarrow{\ q\ }E''\longrightarrow0
	\]
	是 \(M\) 上有限秩实光滑向量丛的短正合列。则存在光滑丛同态 \(j:E''\to E\)，使 \(qj=\operatorname{id}_{E''}\)，并由
	\[
		E'\oplus E''\longrightarrow E,\qquad (u,w)\longmapsto i(u)+j(w)
	\]
	得到光滑丛同构。这样的分裂一般不唯一。
\end{theorem}
\begin{proof}
	由\cref{g1:thm:constant-rank-bundle-morphism}，\(F=i(E')=\ker q\) 是光滑子丛；这里 \(i,q\) 的纤维秩由逐纤维单射、满射保证恒定。为 \(E\) 选丛度量，令 \(C=F^\perp\)。限制 \(q|_C:C\to E''\) 逐纤维为线性同构：其核是 \(C_p\cap F_p=0\)，任意 \(q_p(v)\) 又等于分解 \(v=v_F+v_C\) 后的 \(q_p(v_C)\)，故它满射。

	\cref{g1:prop:bundle-morphism-matrices}说明，逐纤维可逆的光滑丛同态具有光滑逆；其理由是在局部标架中取可逆光滑矩阵的逆，仍由伴随矩阵除以不为零的行列式给出。于是取 \(j=(q|_C)^{-1}\)，再包含到 \(E\) 中，就有 \(qj=\operatorname{id}\)。对任意 \(v\in E\)，
	\[
		v=(v-jqv)+jqv,\qquad v-jqv\in F,
	\]
	所以所述直和同构的逆为
	\[
		v\longmapsto\bigl(i^{-1}(v-jqv),qv\bigr).
	\]
	这里 \(i:E'\to F\) 也是光滑丛同构，因此该逆光滑。

	最后说明选择的自由度。若 \(j_0,j_1\) 都是分裂，则 \(q(j_1-j_0)=0\)，故其差唯一写成 \(i\circ a\)，其中 \(a:E''\to E'\) 为光滑丛同态；反过来，任意这样的 \(a\) 都使 \(j_0+i\circ a\) 成为分裂。例如平凡丛列 \(\mathbb R\to\mathbb R^2\to\mathbb R\) 中，第一箭头包含为第一坐标、第二箭头投影到第二坐标，则每个光滑函数 \(a\) 都给出分裂 \(t\mapsto(a(p)t,t)\)。因此短正合列本身通常没有指定一个分裂。
\end{proof}

这些构造也说明单位分解可以保持什么。若 \(L:E\to F\) 是固定的丛同态，局部截面都满足同一个线性条件 \(L(s_\alpha)=t\)，则
\[
	L\!\left(\sum_\alpha\varphi_\alpha s_\alpha\right)
	=\sum_\alpha\varphi_\alpha L(s_\alpha)
	=\left(\sum_\alpha\varphi_\alpha\right)t=t.
\]
正定性也由非负权重保持。但两个处处非零的局部截面可以相互抵消，例如平凡线丛中的 \(1\) 与 \(-1\) 以相同权重平均就得到零。逐向量平均两组正交归一标架，甚至未必得到一组标架，更不能自动保持正交归一；非退化交替形式与几乎复结构的代数条件也不能直接由平均保持，\secref{g1:sec:1204}将给出具体例子；含导数的条件还须另作检验。因此“局部存在”能否推出“整体存在”，仍须检验所需条件是否被这个加权构造保留。

% 纲要条目排为编号小节标题，并进入目录与书签；说明段仍保留于纲要注释中。

% !TeX root = ../../Book1.tex
% 纲要标题：### 12.4 张量场与几何结构
% 纲要位置：几何学计划纲要.md，第 594 行。
% 正文状态：初稿；尚未正式审核。

% 纲要内容开始
% 
% * 度量、微分形式和场的光滑性
% * 自然丛运算与拉回
% * 结构群约化的模型及其存在问题
% 
% ---
% 
% 纲要内容结束

\section{张量场与几何结构}
\label{g1:sec:1204}

前面已经分别建立了两种构造：过渡函数把一族向量空间粘成向量丛，单位分解把相容的局部要求组织成整体截面。现在把它们用于几何。度量给每个切空间一个内积，微分形式给每个切空间一个交替的多线性函数；局部矩阵按换标架规律相互转换，并随基点光滑变化，从而给出整体几何对象。哪些变化规律来自代数运算，哪些几何条件可以通过粘合保持，是本节要分开的两个问题。

\subsection{度量、微分形式和场的光滑性}
\label{g1:subsec:1204:01}

先说明怎样把\secref{g1:sec:1202}的逐点运算变成一个光滑空间。设 \(E,F\) 是同一光滑流形 \(M\) 上的有限秩实光滑向量丛。直和已经构造；对偶、张量积和外幂也可以在每个纤维中进行，但要确认它们的换标架仍满足余循环条件。

\begin{proposition}[由纤维运算构造向量丛]
	\label{g1:prop:tensor-bundle-constructions}
	设 \(M\) 为 Hausdorff、第二可数光滑流形，可带边界，\(E,F\to M\) 为有限秩实光滑向量丛。纤维分别为 \(E_p^*\)、\(E_p\otimes F_p\)、\(\operatorname{Hom}(E_p,F_p)\)、\(\operatorname{Sym}^kE_p\)、\(\Lambda^kE_p\) 的不交并，都有与局部标架相容的自然光滑向量丛结构。逐点的求值、收缩、对称化、楔积以及线性映射的复合，均给出相应光滑丛映射。
\end{proposition}

\begin{proof}
	取共同细化的平凡化覆盖，设 \(E,F\) 的向量分量过渡分别为 \(v_i=g_{ij}v_j\)、\(w_i=h_{ij}w_j\)。对偶、张量积和同态的分量过渡依次为
	\[
		\alpha_i=g_{ij}^{-\mathsf T}\alpha_j,\qquad
		z_i=(g_{ij}\otimes h_{ij})z_j,\qquad
		A_i=h_{ij}A_jg_{ij}^{-1}.
	\]
	对称幂和外幂使用由 \(g_{ij}\) 诱导的 \(\operatorname{Sym}^kg_{ij}\)、\(\Lambda^kg_{ij}\)。这些矩阵元由原矩阵元的多项式与可逆矩阵的逆组成，故光滑。每一种代数构造保持恒等映射与复合，因此新的过渡函数也满足余循环条件。例如 \((g_{ij}\otimes h_{ij})(g_{jk}\otimes h_{jk})=g_{ik}\otimes h_{ik}\)。\cref{g1:thm:bundle-cocycle-gluing}便给出相应向量丛。

	如果改变原平凡化，新、旧标架之间的同构也诱导上述各空间之间的同构，所以纤维上的恒等识别给出同一个光滑结构。在这些图中，求值与收缩是坐标分量的有限次乘法和求和，对称化是有限次线性组合，楔积是带符号的有限次乘法和求和，复合是矩阵乘法，故都光滑。它们的定义不依赖基，因而在重叠图中相容，确实成为整体丛映射。
\end{proof}

\begin{definition}
	\label{g1:def:smooth-tensor-field}
	设 \(M\) 为光滑流形，可带边界，\(r,s\geq0\)。记
	\[
		T^r_sM=(\mathrm{T}M)^{\otimes r}\otimes(\mathrm{T}^*M)^{\otimes s}.
	\]
	它的光滑截面称为 \((r,s)\) 型\term[zhangliangchang]{张量场}[tensor field]。零次张量积约定为平凡实线丛，因此 \((0,0)\) 型张量场就是光滑实函数。\(\mathrm{T}M\) 的光滑截面称为\term[xiangliangchang]{向量场}[vector field]，其空间记为 \(\mathfrak X(M)\)；\(\Lambda^k\mathrm{T}^*M\) 的光滑截面称为光滑 \(k\) 次\term[weifenxingshi]{微分形式}[differential form]，其空间记为 \(\Omega^k(M)\)。\symidx{\(T^r_sM\)：张量丛}\symidx{\(\Omega^k(M)\)：光滑微分形式空间}
\end{definition}

设局部切标架为 \(e_1,\ldots,e_n\)，对偶标架为 \(\varepsilon^1,\ldots,\varepsilon^n\)。一个张量场写成
\[
	T=\sum T^{i_1\cdots i_r}_{j_1\cdots j_s}\,
	e_{i_1}\otimes\cdots\otimes e_{i_r}\otimes
	\varepsilon^{j_1}\otimes\cdots\otimes\varepsilon^{j_s}.
\]
它光滑当且仅当所有分量函数光滑，这是向量丛截面的局部判据。\cref{g1:prop:tensor-component-change}说明换标架后分量乘上相应过渡矩阵及其逆，因而光滑性不依赖选择。这些分量不是各自独立的几何不变量，只有连同换标架规律理解，才能恢复同一个 \(T\)。

\ProseParagraphBreak

张量场还可以通过它如何作用于其他场来识别。一个 \((r,s)\) 型张量在每一点输入 \(r\) 个协变向量和 \(s\) 个切向量；逐点求值得到光滑函数。这里的多线性不仅是对实数线性，而且是对光滑函数线性。例如对一个协变二阶张量，\(T(fX,Y)=fT(X,Y)\)，不应额外出现 \(f\) 的导数。

\begin{proposition}[张量场的多线性判据]
	\label{g1:prop:tensor-field-multilinearity}
	在 Hausdorff、第二可数光滑流形 \(M\)（可带边界）上，光滑 \((r,s)\) 型张量场与 \(C^\infty(M)\) 多线性映射
	\[
		\underbrace{\Omega^1(M)\times\cdots\times\Omega^1(M)}_{r\text{ 个}}
		\times\underbrace{\mathfrak X(M)\times\cdots\times\mathfrak X(M)}_{s\text{ 个}}
		\longrightarrow C^\infty(M)
	\]
	一一对应，对应为逐点求值。当 \(r=s=0\) 时，右端只表示指定一个光滑函数。
\end{proposition}

\begin{proof}
	张量场给出的映射具有所述多线性，坐标公式还说明输出光滑。反向设给定多线性映射 \(A\)，首先证明它在点 \(p\) 的值只看各输入场在 \(p\) 的取值。

	固定其他输入，若某个输入场 \(Z\) 在 \(p\) 的邻域内恒为零，取一个在 \(p\) 附近为一、支集包含于该邻域的光滑截断函数 \(\chi\)，其存在由\cref{g1:cor:smooth-cutoff-closed}保证。于是 \(\chi Z=0\)，由多线性，\(\chi A(\ldots,Z,\ldots)=0\)，故输出在 \(p\) 为零。因此局部相同的输入可互相替换。

	在 \(p\) 附近取一个局部标架，将它乘以支集包含于图域且在更小邻域为一的截断函数，再在图域外延零，就得到在较小邻域等于该标架的整体场；对偶标架同样处理。任一输入场 \(Z\) 在该邻域可展开为 \(\sum z^i e_i\)。再用一个支集包含于此邻域的截断函数把系数延成整体光滑函数，由刚才的局部性与 \(C^\infty(M)\) 多线性，在点 \(p\) 得
	\[
		A(\ldots,Z,\ldots)(p)=\sum_i z^i(p)A(\ldots,e_i,\ldots)(p),
	\]
	其中右端的 \(e_i\) 表示已经延拓的场。对每个输入重复，便证明输出只依赖于各输入在 \(p\) 的值，并给出一个有限维多线性函数，即张量 \(T_p\)。这里允许把任意点处向量作为输入，是因为局部常系数标架组合也能用同样方法延拓。

	在固定的较小图域内，\(T\) 的每个分量都是 \(A\) 作用于若干固定整体标架场后得到的光滑函数，故 \(T\) 光滑。它的逐点求值正是 \(A\)，且各点全部输入值都可由整体场实现，所以 \(T\) 唯一。
\end{proof}

这个判据也帮助排除看起来像张量的表达。在坐标区间上，令 \(B(f\partial_x,g\partial_x)=f g'\)。它对实数双线性，却有 \(B(X,hY)=hB(X,Y)+X(h)g\)，一般不对光滑函数线性，所以不是二阶张量。额外出现的导数表明它还在观察输入场在附近如何变化，而不只观察一个点的值。

\begin{definition}
	\label{g1:def:riemannian-metric-tensor}
	光滑流形 \(M\) 上的\term[limanduliang]{Riemann 度量}[Riemannian metric]是一个光滑截面 \(g\in\Gamma(\operatorname{Sym}^2\mathrm{T}^*M)\)，使每个 \(g_p\) 都是 \(\mathrm{T}_pM\) 上的正定内积。这里将对称张量与对称多线性形式作\secref{g1:sec:1202}中的自然识别。
\end{definition}

\cref{g1:thm:bundle-metric-existence}应用于 \(\mathrm{T}M\)，保证每个本书意义下的光滑流形都有 Riemann 度量。度量输入同一点的两个向量，不是直接输入两个流形点；由它构造曲线长度与距离，还需另作定义。在标架中记 \(G=(g(e_i,e_j))\)，换标架 \(e'=eA\) 时
\[
	G'=A^{\mathsf T}GA.
\]
因此一个正定矩阵函数只有按这个规律在交叠处相容，才表示整体度量。

\begin{proposition}[度量的对偶识别]
	\label{g1:prop:metric-duality-smooth}
	设 \(g\) 为光滑流形 \(M\)（可带边界）上的 Riemann 度量。映射
	\[
		g^\flat:\mathrm{T}M\longrightarrow \mathrm{T}^*M,\qquad
		v\longmapsto g(v,\cdot)
	\]
	是光滑丛同构，记其逆为 \(g^\sharp\)。对任意 \(f\in C^\infty(M)\)，向量场 \(\operatorname{grad}_g f=g^\sharp(\mathrm{d}f)\) 是唯一满足 \(g(\operatorname{grad}_g f,X)=\mathrm{d}f(X)\) 的光滑向量场。
\end{proposition}

\begin{proof}
	在局部标架中，\(g^\flat\) 的矩阵就是 \(G\)。正定性使 \(G\) 可逆，逆矩阵各项由余子式除以非零行列式表示，所以 \(G^{-1}\) 光滑。这给出光滑的逐纤维逆，且求值定义保证各图中的构造相容。令 \(\operatorname{grad}_g f=g^\sharp(\mathrm{d}f)\) 就得到等式；非退化性保证唯一。坐标表达为 \(\operatorname{grad}_g f=\sum_{i,j}g^{ij}\partial_jf\,\partial_i\)，其中 \((g^{ij})=(g_{ij})^{-1}\)。
\end{proof}

对微分形式，若 \(x^1,\ldots,x^n\) 是坐标，则每个 \(k\) 次形式唯一写成
\[
	\alpha=\sum_{i_1<\cdots<i_k}a_{i_1\cdots i_k},
	\mathrm{d}x^{i_1}\wedge\cdots\wedge\mathrm{d}x^{i_k},
\]
光滑性等价于这些系数光滑。\(k>n\) 时只有零形式，\(\Omega^0(M)=C^\infty(M)\)。逐点楔积定义 \(\Omega^p(M)\times\Omega^q(M)\to\Omega^{p+q}(M)\)，逐点插入向量场定义 \(\iota_X\alpha\in\Omega^{k-1}(M)\)。它们光滑，并满足\cref{g1:prop:wedge-algebra,g1:prop:interior-product-rule}的分次公式，因为每一点上都满足同样的代数恒等式。外微分将由\chapref{g1:ch:14}另行构造；当前的一形式 \(\mathrm{d}f\) 已有第十一章的定义。

\begin{example}[极坐标中的度量与面积]
	\label{g1:exm:polar-tensors-density}
	在 \(r>0\) 的一个角扇区中，\(x=r\cos\theta\)、\(y=r\sin\theta\)。Euclid 度量和标准二形式的坐标表达分别为
	\[
		g=\mathrm{d}r\otimes\mathrm{d}r+r^2\mathrm{d}\theta\otimes\mathrm{d}\theta,
		\mathrm{d}x\wedge\mathrm{d}y=r\,\mathrm{d}r\wedge\mathrm{d}\theta.
	\]
	第一式由 \(\mathrm{d}x=\cos\theta\,\mathrm{d}r-r\sin\theta\,\mathrm{d}\theta\)、\(\mathrm{d}y=\sin\theta\,\mathrm{d}r+r\cos\theta\,\mathrm{d}\theta\) 展开后交叉项抵消得到；第二式由楔积的反交换关系得到。\(g\) 的矩阵为 \(\operatorname{diag}(1,r^2)\)，所以体积密度为 \(r|\mathrm{d}r\wedge\mathrm{d}\theta|\)。若把角度坐标反向，二形式的相应系数改变符号，密度使用绝对行列式，因而仍为正。不同坐标中的系数虽不同，度量、二形式与密度各自仍是同一个整体对象。
\end{example}

\subsection{自然丛运算与拉回}
\label{g1:subsec:1204:02}

“在 \(N\) 上有一个对象，把它拉回到 \(M\)”有两种不同含义。只改变基空间，得到的是拉回丛 \(F^*E\) 上的数据；如果还要让它作用于 \(M\) 自己的切向量，就需要使用 \(\mathrm{d}F\)。前一步保留目标空间的纤维，后一步把源切向量作为输入，二者由微分联系。

\begin{proposition}[丛运算与基空间拉回]
	\label{g1:prop:bundle-operations-pullback}
	设 \(F:M\to N\) 为光滑映射，两流形可带边界，\(E,H\to N\) 为有限秩实光滑向量丛。存在自然光滑丛同构
	\[
		F^*(E\otimes H)\cong F^*E\otimes F^*H,\qquad
		F^*(E^*)\cong(F^*E)^*,\qquad
		F^*(\Lambda^kE)\cong\Lambda^k(F^*E),
	\]
	对直和、对称幂和同态丛也有相同性质。每个光滑截面 \(s\) 都给出拉回截面 \((F^*s)(p)=(p,s(F(p)))\)，且这些同构与各代数运算相容。
\end{proposition}

\begin{proof}
	在 \(p\) 处，拉回丛的纤维就是带基点标记的 \(E_{F(p)}\)。因此第一项同构把 \((p,e\otimes h)\) 送到 \((p,e)\otimes(p,h)\)，其余各项由同一个纤维的对偶、外幂等恒等识别给出。它们与线性换标架相容。在一张 \(E,H\) 的平凡化图拉回后的图中，上述映射在纤维分量上都是固定线性同构，所以光滑，逆也光滑。截面的局部系数从 \(s^a(y)\) 变为 \(s^a(F(x))\)，故光滑；逐点定义保证它与所列运算相容。
\end{proof}

例如，\(N\) 上的向量场 \(Y\) 总能给出 \(F^*\mathrm{T}N\) 的截面 \(p\mapsto(p,Y_{F(p)})\)，但这还是位于目标切空间的向量。为把 \(M\) 的向量送到那里，可以用覆盖恒等基映射的丛同态
\[
	\widehat{\mathrm{d}F}:\mathrm{T}M\longrightarrow F^*\mathrm{T}N,
	\qquad v_p\longmapsto(p,\mathrm{d}F_pv_p).
\]
\cref{g1:prop:differential-chain-rule}及第十一章的切丛坐标表明它光滑。其对象关系是
\[
	\begin{tikzcd}[column sep=large,row sep=large]
		\mathrm{T}M \arrow[r,"\widehat{\mathrm{d}F}"] \arrow[d] & F^*\mathrm{T}N \arrow[r] \arrow[d] & \mathrm{T}N \arrow[d] \\
		M \arrow[r,"\operatorname{id}_M"'] & M \arrow[r,"F"'] & N.
	\end{tikzcd}
\]
右上箭头忘记拉回纤维的基点标记，三条竖箭头均为丛投影。

\begin{proposition}[协变张量的拉回]
	\label{g1:prop:covariant-tensor-pullback}
	设 \(F:M\to N\)、\(G:N\to P\) 为光滑映射，各流形可带边界，\(T\) 是 \(N\) 上的光滑协变 \(k\) 阶张量。公式
	\[
		(F^*T)_p(v_1,\ldots,v_k)
		=T_{F(p)}(\mathrm{d}F_pv_1,\ldots,\mathrm{d}F_pv_k)
	\]
	定义 \(M\) 上的光滑协变 \(k\) 阶张量。拉回保持张量积、对称性、交替性及楔积，且对 \(P\) 上的同类张量满足 \((G\circ F)^*=F^*\circ G^*\)。当 \(T=g\) 为 Riemann 度量时，\(F^*g\) 半正定；它正定当且仅当 \(F\) 为浸入。
\end{proposition}

\begin{proof}
	先把 \(T\) 拉到 \(F^*((\mathrm{T}^*N)^{\otimes k})\)，再在每个输入中复合 \(\widehat{\mathrm{d}F}\)，便得到所给逐点公式。局部展开时，其系数由 \(T\) 的系数与 \(F\) 复合，再乘上 \(F\) 的一阶偏导数并有限求和，因此光滑。交换或交替输入不改变这种复合，故对称性和交替性保留。张量积与楔积的相容性由各自的多线性公式逐项得到；复合公式来自微分的链式法则。

	度量情形中，\((F^*g)_p(v,v)=g_{F(p)}(\mathrm{d}F_pv,\mathrm{d}F_pv)\geq0\)，且等于零当且仅当 \(\mathrm{d}F_pv=0\)。所以它在每点正定，恰好等价于每点微分单射。
\end{proof}

例如抛物线参数 \(F(t)=(t,t^2)\) 将 Euclid 度量拉回为 \((1+4t^2)\,\mathrm{d}t\otimes\mathrm{d}t\)，处处正定。投影 \(P(x,y)=x\) 却把直线度量拉回为 \(\mathrm{d}x\otimes\mathrm{d}x\)，在 \(\partial_y\) 方向为零。这两个计算分别体现浸入保留全部输入方向、淹没可以忘掉方向。

\ProseParagraphBreak

对向量场不能直接照抄同一公式。若要从 \(N\) 上的 \(Y\) 得到 \(M\) 上的 \(X\)，至少要解 \(\mathrm{d}F_pX_p=Y_{F(p)}\)。对投影 \(P:\mathbb R^2\to\mathbb R\) 与 \(Y=\partial_t\)，\(X=\partial_x+b(x,y)\partial_y\) 都是解，没有指定唯一选择；对包含 \(I(t)=(t,0)\) 与目标向量场 \(\partial_y\)，则没有解。相比之下，协变张量只需让目标观测复合已有的 \(\mathrm{d}F\)，所以总能拉回。

\begin{proposition}[可逆微分下的混合张量]
	\label{g1:prop:mixed-tensor-pullback}
	设 \(F:M\to N\) 为光滑映射，两流形可带边界，并假设每个 \(\mathrm{d}F_p\) 都是线性同构。则 \(N\) 上任意光滑 \((r,s)\) 型张量场都可按
	\[
		(F^*T)_p=
		\bigl((\mathrm{d}F_p)^{-1}\bigr)^{\otimes r}
		\otimes\bigl((\mathrm{d}F_p)^*\bigr)^{\otimes s}
		\bigl(T_{F(p)}\bigr)
	\]
	拉回为 \(M\) 上的同型张量场。该操作与张量积和收缩相容；两个满足条件的映射复合时，拉回仍满足反向复合公式。特别地，微分同胚满足这些条件。
\end{proposition}

\begin{proof}
	\(\mathrm{d}F\) 的局部矩阵光滑且可逆，故逆矩阵光滑，公式因而定义光滑张量场。收缩中一次切向量与一次协变向量的配对变成 \((\alpha\circ\mathrm{d}F)((\mathrm{d}F)^{-1}v)=\alpha(v)\)，所以收缩不变。张量积相容性由定义直接得到，复合公式则使用链式法则以及复合线性同构的逆按反向顺序复合。
\end{proof}

密度还需要单独辨析。密度丛本身作为向量丛，可以沿任意 \(F\) 拉回；但若要把目标的顶维密度变成源上的同维密度，公式中会出现 \(|\det \mathrm{D}F|\)。在局部微分同胚下，行列式处处非零，绝对值仍光滑，所以这一步没有问题。对一般光滑映射则可能失去光滑性：\(F(t)=t^2\) 会把 \(|\mathrm{d}y|\) 按该逐点规则变成 \(2|t|\,|\mathrm{d}t|\)，它在零点不光滑。\cref{g1:prop:covariant-tensor-pullback}适用于交替形式，不能因此不加条件地套到密度上。

\subsection{结构群约化的模型及其存在问题}
\label{g1:subsec:1204:03}

一个任意标架允许全部可逆线性换基。如果选定内积，就可以只使用正交单位标架，其间换基矩阵只能是正交矩阵。这个观察把“额外几何结构”与“允许怎样换标架”联系起来。这里先用局部标架给出定义；一般主丛的语言将在\chapref{g1:ch:18}系统展开。

\begin{definition}
	\label{g1:def:frame-structure-reduction}
	设 \(E\to M\) 为秩 \(r\geq1\) 的实光滑向量丛，\(H\subset\operatorname{GL}(r,\mathbb R)\) 为矩阵子群。一份 \(H\)-\term[jiegouqunyuehua]{结构群约化}[reduction of structure group]由覆盖上的局部光滑标架 \(e_i\) 给出，要求所有过渡矩阵 \(e_j=e_i g_{ij}\) 都取值于 \(H\)。在同一基点处把 \(e_i(p)\) 右乘全部 \(H\) 元素得到的标架子集，因过渡条件而与 \(i\) 无关；记为 \(P_{H,p}\)。定义同一约化的局部标架，允许在各自图中再右乘光滑的 \(H\) 值函数，并允许细化覆盖。
\end{definition}

局部标架之间的矩阵约定与\cref{g1:def:frame-bundle}一致：标架是同构 \(\mathbb R^r\to E_p\)，右作用为复合同构。在 \(P_{H,p}\) 上，\(H\) 的右作用仍自由且传递。几何结构由这组适配标架确定；限制允许的换基群，正是要求额外结构在换标架后保持不变。

\begin{proposition}[度量与正交标架]
	\label{g1:prop:metric-orthogonal-reduction}
	对 Hausdorff、第二可数光滑流形 \(M\)（可带边界）上的秩 \(r\) 实光滑向量丛 \(E\to M\)，给定丛度量等价于给定 \(\operatorname{O}(r)\) 结构群约化。这样的约化总存在。
\end{proposition}

\begin{proof}
	若给定度量，在任一局部标架 \(e_1,\ldots,e_r\) 上逐步定义
	\[
		u_k=e_k-\sum_{j<k}g(e_k,v_j)v_j,
		\qquad v_k=\frac{u_k}{\sqrt{g(u_k,u_k)}}.
	\]
	先前的 \(v_j\) 张成 \(e_1,\ldots,e_{k-1}\) 的空间，原标架线性无关，所以 \(u_k\neq0\)。分母因正定性处处为正，各步光滑，最终得到局部正交单位标架。两个这样的标架之间的矩阵 \(A\) 满足 \(A^{\mathsf T}A=I\)，所以属于 \(\operatorname{O}(r)\)。这些标架的 \(\operatorname{O}(r)\) 轨道就是全部正交单位标架。

	反向给定这样的约化，在每个所选标架中规定度量矩阵为 \(I\)。因为过渡矩阵正交，不同图中的规定一致，得到整体光滑正定度量。两种构造互逆。最后由\cref{g1:thm:bundle-metric-existence}，度量总存在，故约化也总存在。
\end{proof}

\begin{proposition}[定向与体积形式的标架]
	\label{g1:prop:orientation-volume-reductions}
	设 \(E\to M\) 为秩 \(r\geq1\) 的实光滑向量丛。连续选择每个纤维的定向，且在局部标架中选择一致，等价于 \(\operatorname{GL}^+(r,\mathbb R)\) 约化，其中上标 \(+\) 表示行列式为正。一个处处非零的光滑截面 \(\nu\in\Gamma(\det E^*)\) 等价于 \(\operatorname{SL}(r,\mathbb R)\) 约化。若底流形 Hausdorff、第二可数，可带边界，则这种 \(\nu\) 存在当且仅当 \(E\) 可定向。
\end{proposition}

\begin{proof}
	给定定向，局部取正向标架；两正向标架之间行列式为正。反向由正行列式的过渡保证各图的定向一致，所以第一项等价成立。

	给定非零 \(\nu\)，在一个局部标架上设 \(a=\nu(e_1,\ldots,e_r)\)，它是处处非零光滑函数。把 \(e_1\) 换成 \(e_1/a\)，其余不动，即得满足 \(\nu(e_1,\ldots,e_r)=1\) 的标架。这些标架之间的行列式为一，给出 \(\operatorname{SL}\) 约化。反向在约化标架中规定 \(\nu(e_1,\ldots,e_r)=1\)，过渡矩阵行列式为一使规定相容，故得到非零光滑顶次形式，且两种构造互逆。

	非零顶次形式将使其取值为正的标架定义为正向，因而给出定向。若已有定向，先由\cref{g1:thm:bundle-metric-existence}取度量，再取正向正交单位标架；把它们的顶次对偶楔积定义为 \(\nu\)。两这样的标架之间矩阵既正交又有正行列式，故行列式为一，顶次形式在交叠处相同。这就构造出所需 \(\nu\)。
\end{proof}

度量总给出正密度，而不必给出带符号的非零顶次形式。差别在于，密度的转换使用绝对行列式，反向标架不造成符号冲突。因而“有处处正的体积密度”不能用于证明可定向。

\begin{proposition}[非退化交替形式的标架]
	\label{g1:prop:symplectic-frame-reduction}
	设 \(E\to M\) 为秩 \(2m\) 的实光滑向量丛，\(m\geq1\)。给定光滑、逐纤维非退化的交替双线性形式 \(\omega\)，等价于过渡矩阵取值于
	\[
		\operatorname{Sp}(2m,\mathbb R)
		=\{A\in\operatorname{GL}(2m,\mathbb R):A^{\mathsf T}J_0A=J_0\},
		\qquad J_0=\begin{pmatrix}0&I_m\\-I_m&0\end{pmatrix}
	\]
	的结构群约化。这样的 \(\omega\) 总给出非零顶次形式 \(\omega^m/m!\)，因此使 \(E\) 可定向。
\end{proposition}

\begin{proof}
	需要证明适配标架可以光滑选择，而不只是在每个点分别选出一个基。在一点附近先取一个非零局部截面 \(u_1\)。非退化性保证在该点某个标架向量 \(w\) 满足 \(\omega(u_1,w)\neq0\)，缩小邻域后这个函数仍非零，令 \(v_1=w/\omega(u_1,w)\)，就有 \(\omega(u_1,v_1)=1\)。

	对任意局部向量 \(z\)，令
	\[
		z'=z+\omega(v_1,z)u_1-\omega(u_1,z)v_1.
	\]
	直接代入得 \(\omega(u_1,z')=\omega(v_1,z')=0\)。因此这个光滑投影的像为两者的共同交替正交空间 \(W\)，核为 \(\operatorname{span}\{u_1,v_1\}\)，秩恒为 \(2m-2\)。由\cref{g1:thm:constant-rank-bundle-morphism}，\(W\) 是光滑子丛。\(\omega|_W\) 非退化：若 \(z\in W\) 与 \(W\) 中所有向量配对为零，它又与 \(u_1,v_1\) 配对为零，因直和分解而与整个纤维配对为零，所以 \(z=0\)。在 \(W\) 上重复，得到光滑标架 \(u_1,\ldots,u_m,v_1,\ldots,v_m\)，其中 \(\omega(u_i,v_j)=\delta_{ij}\)，同类向量间配对为零。

	在这些标架中 \(\omega\) 的矩阵为 \(J_0\)，故过渡矩阵满足所列条件。反向，在这样的约化标架中规定矩阵为 \(J_0\)，便粘合成光滑非退化交替形式。最后在对偶标架中写 \(\omega=\sum_i u^i\wedge v^i\)。各个二次项两两交换，各自平方为零，故
	\[
		\frac{\omega^m}{m!}=u^1\wedge v^1\wedge\cdots\wedge u^m\wedge v^m\neq0.
	\]
	由\cref{g1:prop:orientation-volume-reductions}得到定向。
\end{proof}

另一个常见的逐点结构，是实向量丛上的光滑丛自同态 \(J\)，满足 \(J^2=-I\)。规定 \((a+ib)v=av+bJv\)，就把每个实纤维变成复向量空间。在切丛上，这样的 \(J\) 称为\term[jihufujiegou]{几乎复结构}[almost complex structure]；此处只要求代数恒等式与光滑性，尚没有要求特殊坐标的存在。

\begin{proposition}[复线性标架]
	\label{g1:prop:complex-frame-reduction}
	设 \(\pi:E\to M\) 为有限秩实光滑向量丛，底流形 Hausdorff、第二可数，可带边界。若有覆盖 \(\mathrm{id}_M\) 的光滑实线性丛自同态 \(J:E\to E\)，即 \(\pi\circ J=\pi\) 且每个 \(J_p:E_p\to E_p\) 实线性，并满足 \(J_p^2=-\mathrm{id}_{E_p}\)，则秩为偶数 \(2m\)，且存在局部光滑复标架；相应实标架的过渡属于 \(\operatorname{GL}(m,\mathbb C)\) 在 \(\operatorname{GL}(2m,\mathbb R)\) 中的标准实表示。反向，这样的约化定义唯一的 \(J\)。该结构总使 \(E\) 可定向。
\end{proposition}

\begin{proof}
	每个纤维通过上述复数乘法成为有限维复向量空间，实维数因此为偶数。为光滑地选基，先取任意丛度量 \(g\)，定义
	\[
		h(v,w)=g(v,w)+g(Jv,Jw).
	\]
	这是光滑正定度量，且 \(h(Jv,Jw)=h(v,w)\)。一个局部非零截面 \(u_1\) 与 \(Ju_1\) 实线性无关：若 \(Ju_1=a u_1\)，再作用 \(J\) 就得 \(-u_1=a^2u_1\)，矛盾。它们张成一个 \(J\) 不变的秩二子丛，其 \(h\) 正交补也是 \(J\) 不变的，因为 \(h(Jv,w)=-h(v,Jw)\)。由\cref{g1:prop:orthogonal-complement-bundle}正交补光滑，在其上重复选择，得到 \(u_1,Ju_1,\ldots,u_m,Ju_m\) 的局部光滑实标架，等价于复标架 \(u_1,\ldots,u_m\)。

	两个复标架之间是可逆复矩阵，视为实矩阵时与标准复乘法相交换。反向若过渡矩阵都复线性，在每张标架中规定 \(J\) 为乘以 \(i\)，交叠处便相容，给出所需 \(J\)。两项构造互逆。

	若复矩阵为 \(A+iB\)，按实部、虚部分组的实矩阵为 \(\begin{psmallmatrix}A&-B\\B&A\end{psmallmatrix}\)。在复数上改变基，把实部与虚部组合为 \(v+iw\)、\(v-iw\)，该矩阵变成对角块 \(A+iB\)、\(A-iB\)，所以它的实行列式为 \(|\det_{\mathbb C}(A+iB)|^2>0\)。所有复标架因而给出一致实定向。
\end{proof}

这里的存在性已经出现了整体限制。令 \(L\to S^1\) 为\cref{g1:exm:mobius-line-bundle}的 Möbius 实线丛，并取 \(E=L\oplus\underline{\mathbb R}\)。沿圆周绕一周，\(E\) 的标架过渡可表示为 \(\operatorname{diag}(-1,1)\)，所以其行列式线丛 \(\det E\) 与 \(L\) 相同。若 \(E\) 可定向，取度量并选正向单位体积向量，就会得到 \(\det E\) 的处处非零截面；但 \(L\) 没有这样的截面。因此 \(E\) 不可定向。它虽然秩为二，仍不能有上面两类结构中的任何一种。这说明偶数秩只是纤维上的必要条件，不能代替过渡函数的整体相容。

\ProseParagraphBreak

单位分解不能消除这些障碍。在二维向量空间上，\(\omega\) 与 \(-\omega\) 都可非退化，但等权平均为零；\(J\) 与 \(-J\) 都满足平方为 \(-I\)，平均后却不再满足。相反，正定双线性型的非负加权和，只要权重和为一，就仍正定。这正是丛度量总能粘合、而其他结构不能照搬同一证明的原因。即使某种逐点结构已存在，它是否满足进一步的微分条件，仍须另行研究。
% 纲要条目排为编号小节标题，并进入目录与书签；说明段仍保留于纲要注释中。


\begin{chaptersummary}
	向量丛把局部乘积通过光滑可逆矩阵粘合；三重交叠上的余循环条件保证识别相容。一个整体标架等价于整体平凡化，局部标架则总存在。丛同态的核与像要成为固定秩光滑子丛，需要相应的常秩条件；拉回改变底空间，商丛则逐纤维消去指定子丛，两者都必须同时构造局部平凡化。

	\ProseParagraphBreak

	张量积将多线性运算改写为线性运算，收缩来自向量与协变向量的求值。外幂的行列式配对确定了楔积的归一化，交换次数为 \(p,q\) 的形式会产生 \((-1)^{pq}\)。分量随标架改变，完整配对得到的数却不变。最高次形式记录符号，密度使用绝对行列式；度量总产生正密度，产生相应的体积形式还需要定向。

	\ProseParagraphBreak

	光滑单位分解由紧穷竭、局部有限的坐标细化、截断函数和归一化构成。局部有限性使各点附近的求和真正成为有限和；支集控制使局部对象能够光滑补零。它们共同给出闭集上的截面延拓和整体丛度量。正交补进一步分裂短正合丛列，但度量与分裂通常依赖选择，相对延拓也须保留明确的邻域相容条件。

	\ProseParagraphBreak

	张量场是在张量丛中的光滑截面，光滑分量与局部标架无关。任意光滑映射都能拉回向量丛，也能利用微分拉回协变张量；把目标向量变为源向量则一般缺少所需的逆。结构群约化用适配标架组织额外几何结构。正交约化总存在，定向、非退化交替形式和几乎复结构的存在却受整体相容性限制，不能仅凭各纤维有相同的线性模型就断定。
\end{chaptersummary}

\BookExercises

\begin{exercise}\label{g1:ex:mobius-sum-double-pullback}
	令 \(L\to S^1\) 为关系 \((t+2\pi,v)\sim(t,-v)\) 定义的 Möbius 线丛，令 \(f:S^1\to S^1\)，\(f(e^{\mathrm{i}t})=e^{2\mathrm{i}t}\)。
	\begin{enumerate}
		\item 构造 \(L\oplus L\) 的全局标架，并证明它是平凡秩二实向量丛。
		\item 构造 \(f^*L\) 的处处非零光滑截面，并解释这与 \(L\) 非平凡为何不矛盾。
	\end{enumerate}
\end{exercise}
\begin{solution}
	1. 记
	\[
		R(\theta)=\begin{pmatrix}\cos\theta&-\sin\theta\\\sin\theta&\cos\theta\end{pmatrix}.
	\]
	在角度参数 \(t\) 下，用 \(R(t/2)\) 的两列分别作为两个截面的纤维分量，即
	\[
		s_1(e^{\mathrm{i}t})=([t,\cos(t/2)],[t,\sin(t/2)]),\qquad
		s_2(e^{\mathrm{i}t})=([t,-\sin(t/2)],[t,\cos(t/2)]).
	\]
	把 \(t\) 换成 \(t+2\pi\) 后，四个分量都变号；同时两条线丛的粘合也都变号，所以这两组有序对不变。它们因此在圆周上良定义，并且在每个提升圆弧上由光滑函数表示。两截面的分量矩阵为 \(R(t/2)\)，行列式恒为一，故逐点线性无关。\cref{g1:prop:global-frame-triviality}给出平凡化
	\[
		S^1\times\mathbb R^2\longrightarrow L\oplus L,
		\qquad(p,(a,b))\longmapsto a s_1(p)+b s_2(p).
	\]
	其局部逆通过乘 \(R(-t/2)\) 求出，因而也光滑。这个构造用了两条截面，而不只是秩二丛的一条非零截面。

	2. 对 \(f^*L\)，定义
	\[
		s(e^{\mathrm{i}t})=(e^{\mathrm{i}t},[2t,1]).
	\]
	右边第二分量位于 \(e^{2\mathrm{i}t}=f(e^{\mathrm{i}t})\) 上，故确实属于拉回纤维。把 \(t\) 换成 \(t+2\pi\) 时，第二分量变为 \([2t+4\pi,1]=[2t,1]\)，因为绕目标圆周两次产生两次变号。于是截面良定义，局部光滑且处处非零，故给出拉回线丛的全局标架。\(L\oplus L\) 改变了纤维秩，\(f^*L\) 改变了沿底空间环路所经历的粘合；两个结论都没有在原线丛 \(L\) 上构造非零全局截面，因此与其非平凡性相容。
\end{solution}

\begin{exercise}\label{g1:ex:smooth-idempotent-subbundles}
	设 \(M\) 为连通的 Hausdorff、第二可数光滑流形，可带边界，\(E\to M\) 为有限秩实光滑向量丛。若同底光滑丛同态 \(P:E\to E\) 满足 \(P^2=P\)，证明 \(\ker P\) 与 \(\operatorname{im}P\) 是光滑子丛，且加法给出自然丛同构
	\[
		\ker P\oplus\operatorname{im}P\simeq E.
	\]
	再对 \(S^2\times\mathbb R^3\) 上的 \(P_p(v)=v-\langle v,p\rangle p\) 求核、像和相应商丛，并写出这个分解的两个投影。
\end{exercise}
\begin{solution}
	对任一纤维向量 \(v\)，有
	\[
		v=(v-Pv)+Pv,\qquad P(v-Pv)=0.
	\]
	而若 \(w\in\ker P\cap\operatorname{im}P\)，写成 \(w=Pu\) 后有 \(w=P^2u=Pw=0\)。所以逐纤维有直和分解 \(E_p=\ker P_p\oplus\operatorname{im}P_p\)。在适应该分解的一组线性基下，\(P_p\) 的矩阵是 \(\operatorname{diag}(0,I_k)\)，故 \(\operatorname{tr}P_p=\operatorname{rank}P_p\)。迹不依赖基，因为 \(\operatorname{tr}(AB)=\operatorname{tr}(BA)\)，所以局部矩阵的迹给出一个整体光滑函数。

	这个函数只取整数值，故局部常值：在一点附近把它与该点值的差控制在 \(1/2\) 以内，整数性就迫使两者相同。底空间连通，因而秩在整个 \(M\) 上恒定。由\cref{g1:thm:constant-rank-bundle-morphism}，核与像都是光滑子丛。加法映射逐纤维同构，且逆明确为
	\[
		E\longrightarrow\ker P\oplus\operatorname{im}P,\qquad
		v\longmapsto(v-Pv,Pv).
	\]
	这些映射在子丛局部标架中光滑，故确为光滑丛同构。这里投影 \(P\) 本身已经指定了补丛，所以不需要另选内积。

	对题给球面例，\(|p|=1\) 蕴含
	\[
		P_p^2=P_p,\qquad
		\ker P_p=\mathbb Rp,\qquad
		\operatorname{im}P_p=p^\perp.
	\]
	核丛由全局非零截面 \(p\mapsto p\) 张成，故是平凡线丛。像丛就是 \(\mathrm{T}S^2\) 在环境 \(\mathbb R^3\) 中的实现：对球面内曲线求导 \(\langle\gamma,\gamma\rangle=1\)，得到速度垂直于 \(p\)；反之，对 \(0\ne v\perp p\)，曲线
	\[
		\gamma(s)=\cos(s|v|)p+\frac{\sin(s|v|)}{|v|}v
	\]
	位于球面且满足 \(\gamma(0)=p\)、\(\gamma'(0)=v\)，零速度由常曲线实现。球面的局部图是一个坐标关于另两个坐标的光滑图像，所以其包含映射微分单射；上述速度计算遂识别了整个切空间。于是
	\[
		(S^2\times\mathbb R^3)/\ker P\simeq \mathrm{T}S^2,
		\qquad [v]\longmapsto v-\langle v,p\rangle p.
	\]
	所需两个投影分别为 \((I-P)_p(v)=\langle v,p\rangle p\) 与 \(P_p(v)=v-\langle v,p\rangle p\)，并给出 \((S^2\times\mathbb R)\oplus \mathrm{T}S^2\simeq S^2\times\mathbb R^3\)。这里切丛与一条平凡线丛的直和是平凡丛，并未据此断言切丛本身平凡。
\end{solution}

\begin{exercise}\label{g1:ex:latitude-pullback-bundle}
	对固定 \(\alpha\in[0,\pi]\)，定义
	\[
		f_\alpha:S^1\longrightarrow S^2,\qquad
		f_\alpha(e^{\mathrm{i}t})=(\sin\alpha\cos t,\sin\alpha\sin t,\cos\alpha).
	\]
	将球面切向量识别为环境中垂直于基点的向量。构造拉回丛 \(f_\alpha^*\mathrm{T}S^2\) 的全局标架；用它写出微分给出的同底丛同态 \(\mathrm{T}S^1\to f_\alpha^*\mathrm{T}S^2\)，并分别讨论 \(0<\alpha<\pi\) 与 \(\alpha=0,\pi\) 时的核、像和商丛。
\end{exercise}
\begin{solution}
	令 \(p(t)=f_\alpha(e^{\mathrm{i}t})\)，定义环境向量
	\[
		u_1(t)=(-\sin t,\cos t,0),\qquad
		u_2(t)=(-\cos\alpha\cos t,-\cos\alpha\sin t,\sin\alpha).
	\]
	直接计算得到 \(\langle p,u_1\rangle=\langle p,u_2\rangle=0\)、\(|u_1|=|u_2|=1\)、\(\langle u_1,u_2\rangle=0\)。因此两向量在 \(\mathrm{T}_{p(t)}S^2\) 中组成基。它们都是 \(2\pi\) 周期的光滑函数，所以
	\[
		e_a(e^{\mathrm{i}t})=(e^{\mathrm{i}t},u_a(t)),\qquad a=1,2,
	\]
	是拉回丛的全局光滑标架。这个公式在 \(\alpha=0,\pi\) 时仍成立：虽然 \(p(t)\) 此时为常点，所选基可以在这一固定平面内随 \(t\) 旋转。

	在源圆周上取全局切标架 \(\tau(t)=(-\sin t,\cos t)\)。微分作用由参数曲线求导得到
	\[
		\mathrm{d}(f_\alpha)_{e^{\mathrm{i}t}}(a\tau(t))
		=a\sin\alpha\,u_1(t).
	\]
	因此改写成到拉回丛的同底丛同态后，其在标架 \(\tau\) 与 \((e_1,e_2)\) 下的矩阵为列 \((\sin\alpha,0)^{\mathsf T}\)。若 \(0<\alpha<\pi\)，这个映射单射，核为零子丛，像为 \(e_1\) 张成的平凡线丛，而目标除以像的商丛由 \([e_2]\) 给出全局标架，亦为平凡线丛。具体商坐标为 \([a e_1+b e_2]\mapsto b\)。

	若 \(\alpha=0\) 或 \(\pi\)，微分为零，核是整个 \(\mathrm{T}S^1\)，像是零子丛；目标除以像仍为秩二的平凡拉回丛。拉回纤维的维数始终为二，微分的秩却可以为一或零。因而拉回丛不能被定义为微分的像，也不等同于源流形的切丛。
\end{solution}

\begin{exercise}\label{g1:ex:quotient-complement-shears}
	在 \(\mathbb R\) 上令 \(E=\mathbb R\times\mathbb R^2\)，\(F_t=\mathbb R(1,0)\)，并把商映射写成 \(q(t,a,b)=(t,b)\)。求出 \(q\) 的全部光滑线性右逆及其像补丛。再考虑丛自同构
	\[
		T_c(t,a,b)=(t,a+c(t)b,b),\qquad c\in C^\infty(\mathbb R).
	\]
	证明每个 \(T_c\) 在 \(F\) 上为恒等，且在 \(E/F\) 上诱导恒等；是否存在一个补丛在所有这些自同构下都不变？
\end{exercise}
\begin{solution}
	若 \(j:E/F\to E\) 是光滑丛同态且 \(qj=\operatorname{id}\)，其在每条纤维上的线性表达必为
	\[
		j_h(t,b)=(t,h(t)b,b).
	\]
	第一分量的系数 \(h(t)\) 由 \(j(t,1)\) 得到，所以光滑；反过来，任意光滑 \(h\) 都给出这样的右逆。因此全部右逆由光滑函数 \(h\) 参数化，其像为
	\[
		C_h|_t=\mathbb R(h(t),1).
	\]
	它由一个处处非零的光滑截面张成，且每个 \((a,b)\) 有唯一分解
	\[
		(a,b)=(a-h(t)b,0)+(h(t)b,b),
	\]
	所以 \(E=F\oplus C_h\)。反之，任一光滑补丛 \(C\) 上的 \(q|_C\) 逐纤维可逆，由\cref{g1:prop:bundle-morphism-matrices}有光滑逆，故也来自某个 \(h\)。

	矩阵 \(\begin{psmallmatrix}1&c(t)\\0&1\end{psmallmatrix}\) 的逆由把 \(c\) 换成 \(-c\) 得到，因此 \(T_c\) 是光滑丛自同构。若 \(b=0\)，它保持 \((a,0)\) 不变，故固定 \(F\) 的每个向量；第二纤维坐标 \(b\) 不变，故诱导商上的恒等。同时
	\[
		T_c\circ j_h=j_{h+c},\qquad T_c(C_h)=C_{h+c}.
	\]
	取 \(c\equiv1\) 就有 \(C_{h+1}\ne C_h\)：若两条直线相同，将生成向量的第二分量都归一为一，就会得到 \(h+1=h\)，矛盾。因此不存在在所有这些自同构下不变的补丛。这个例子中每个补丛都能整体光滑选出，但子丛与商映射并没有指定一个对全部保留这些数据的自同构不变的选择。
\end{solution}

\begin{exercise}
	\label{g1:exr:1202-tensor-rank}
	设 \(V,W\) 为有限维实向量空间，\(T\in W\otimes V^*\)，其对应的线性映射为 \(L=\Phi(T)\)。证明：把 \(T\) 写成纯张量之和所需的最少项数，等于 \(\operatorname{rank}L\)，约定零张量允许空和。由此判断在 \(V=W=\mathbb R^2\) 中，\(e_1\otimes e^1+e_2\otimes e^2\) 能否写成一个纯张量。再给出一个非零而收缩为零的 \(V\otimes V^*\) 中的元素，解释“迹为零”为什么比“张量为零”弱。
\end{exercise}
\begin{solution}
	若 \(T=\sum_{j=1}^k w_j\otimes\alpha_j\)，则
	\[
		L(v)=\sum_{j=1}^k\alpha_j(v)w_j,
	\]
	故 \(\operatorname{im}L\subseteq\operatorname{span}\{w_1,\ldots,w_k\}\)，得到 \(\operatorname{rank}L\le k\)。反过来，设 \(\operatorname{rank}L=r>0\)，取像空间的一组基 \(w_1,\ldots,w_r\)。对每个 \(v\)，唯一写 \(L(v)=\sum_{j=1}^r\alpha_j(v)w_j\)。由 \(L\) 线性及系数表示唯一，\(\alpha_j(v+v')=\alpha_j(v)+\alpha_j(v')\)、\(\alpha_j(cv)=c\alpha_j(v)\)，所以 \(\alpha_j\in V^*\)。于是 \(L=\Phi(\sum_jw_j\otimes\alpha_j)\)，由 \(\Phi\) 单射得到 \(T=\sum_jw_j\otimes\alpha_j\)。若 \(r=0\)，则 \(L=0\)，同样由同构知 \(T=0\)。最少项数因此恰为 \(r\)。

	给出的二维张量对应恒等映射，秩为 \(2\)，故不能写成一个纯张量。另一方面，\(S=e_1\otimes e^2\) 将 \(e_2\) 送到 \(e_1\)，所以非零；但其收缩是 \(e^2(e_1)=0\)。收缩只保留一个实数，是从整个张量空间到 \(\mathbb R\) 的线性观测，会消去许多非零元素；这一例中，它消去的是一个非零的非对角分量。
\end{solution}

\begin{exercise}
	\label{g1:exr:1202-quadratic-symmetric}
	设 \(V\) 为有限维实向量空间，\(B\in\operatorname{Sym}^2(V^*)\) 按本节约定视为对称双线性形式，令 \(q(v)=B(v,v)\)。证明 \(q\) 唯一决定 \(B\)，并写出恢复公式。反过来，在一组基的坐标中给定任意实齐次二次多项式，证明它来自唯一这样的 \(B\)。对
	\[
		q(xe_1+ye_2)=3x^2+4xy+2y^2
	\]
	写出 \(B\) 的矩阵及它在 \(\operatorname{Sym}^2(V^*)\) 中的表达，并计算 \(B(e_1+e_2,e_1-e_2)\)。
\end{exercise}
\begin{solution}
	双线性与对称性给出
	\[
		q(u+v)=B(u,u)+2B(u,v)+B(v,v),
	\]
	所以
	\[
		B(u,v)=\frac12\bigl(q(u+v)-q(u)-q(v)\bigr).
	\]
	这先证明唯一性。若给定齐次二次多项式
	\[
		q(x)=\sum_i a_i(x^i)^2+\sum_{i<j}b_{ij}x^ix^j,
	\]
	取矩阵 \(G\) 的对角元为 \(a_i\)，非对角元 \(G_{ij}=G_{ji}=b_{ij}/2\)，则 \(B(u,v)=[u]^{\mathsf T}G[v]\) 是对称双线性形式，并满足 \(B(v,v)=q(v)\)，证明存在。恢复公式只使用向量加法与实数值，因此换基并不改变它定义的对象。

	本题中
	\[
		G=\begin{pmatrix}3&2\\2&2\end{pmatrix},\qquad
		B=3e^1\odot e^1+4e^1\odot e^2+2e^2\odot e^2.
	\]
	中间项的系数为 \(4\)，因为 \((e^1\odot e^2)(u,v)=\tfrac12(e^1(u)e^2(v)+e^1(v)e^2(u))\)，对应矩阵的两个非对角元都是 \(1/2\)。所以系数 \(4\) 正好产生两个矩阵元 \(2\)。最后，
	\[
		B(e_1+e_2,e_1-e_2)
		=(1,1)\begin{pmatrix}3&2\\2&2\end{pmatrix}\binom{1}{-1}=1.
	\]
\end{solution}

\begin{exercise}
	\label{g1:exr:1202-wedge-insertion}
	设 \(V=\mathbb R^4\)，基与对偶基分别为 \(e_i,e^i\)，令
	\[
		\omega=e^1\wedge e^2+e^3\wedge e^4.
	\]
	计算 \(\omega\wedge\omega\) 以及 \(\iota_v\omega\)，其中 \(v=\sum_iv^ie_i\)。证明 \(v\mapsto\iota_v\omega\) 是 \(V\to V^*\) 的线性同构，且 \(\omega\) 不能写成两个一形式的楔积。最后计算 \(\iota_{e_1}(\omega\wedge\omega)\)，同时用插入的乘法公式核对符号与系数。
\end{exercise}
\begin{solution}
	两个平方项含重复的一形式，均为零。两个交叉项交换的是两个二次块，符号为 \((-1)^4=1\)，故
	\[
		\omega\wedge\omega=2e^1\wedge e^2\wedge e^3\wedge e^4.
	\]
	按插入固定第一位置的约定，
	\[
		\iota_v\omega=v^1e^2-v^2e^1+v^3e^4-v^4e^3.
	\]
	若此一形式为零，对偶基的独立性依次给出 \(v^2=v^1=v^4=v^3=0\)，所以该线性映射单射。两个空间都为四维，故为同构；也可直接把任意一形式 \(\sum_i a_ie^i\) 的原像写成 \(a_2e_1-a_1e_2+a_4e_3-a_3e_4\)。

	若 \(\omega=\alpha\wedge\beta\)，那么
	\[
		\omega\wedge\omega=\alpha\wedge\beta\wedge\alpha\wedge\beta=0,
	\]
	因为移项后出现两个相同的一形式。这与已计算出的非零最高次形式矛盾，所以 \(\omega\) 不是纯楔积。外幂由纯楔积张成，并不意味着每个元素都是纯楔积。

	由最高次形式的表达直接插入，得到
	\[
		\iota_{e_1}(\omega\wedge\omega)=2e^2\wedge e^3\wedge e^4.
	\]
	另一方面，\(\iota_{e_1}\omega=e^2\)，且 \(\deg\omega=2\)，故乘法公式给出
	\[
		\iota_{e_1}(\omega\wedge\omega)
		=e^2\wedge\omega+\omega\wedge e^2
		=2e^2\wedge\omega
		=2e^2\wedge e^3\wedge e^4.
	\]
	第二步交换一次与二次形式，符号为 \((-1)^{1\cdot2}=1\)，正好得到系数 \(2\)。
\end{solution}

\begin{exercise}
	\label{g1:exr:1202-density-orientation}
	在 \(V=\mathbb R^2\) 的标准基 \(e=(e_1,e_2)\) 中，取内积矩阵及新基
	\[
		G=\begin{pmatrix}2&1\\1&2\end{pmatrix},\qquad
		e'=eA,\qquad A=\begin{pmatrix}1&1\\0&-2\end{pmatrix}.
	\]
	证明 \(G\) 正定，计算新基中的内积矩阵，以及 \(\mu_g(e)\)、\(\mu_g(e')\)。取使标准基正向的定向，写出体积形式 \(\omega\) 和体积密度 \(\mu_g\) 在新旧两组对偶基中的表达。若把定向反过来，哪些对象发生改变？
\end{exercise}
\begin{solution}
	对非零 \((x,y)\)，有
	\[
		(x,y)G\binom{x}{y}=2x^2+2xy+2y^2=(x+y)^2+x^2+y^2>0,
	\]
	故 \(G\) 正定。直接相乘得
	\[
		G'=A^{\mathsf T}GA=\begin{pmatrix}2&0\\0&6\end{pmatrix}.
	\]
	因此 \(\mu_g(e)=\sqrt{\det G}=\sqrt3\)，而 \(\mu_g(e')=\sqrt{12}=2\sqrt3\)，也等于 \(|\det A|\mu_g(e)\)。

	记 \(\varepsilon=e^1\wedge e^2\)、\(\varepsilon'=e'^1\wedge e'^2\)。由于 \(\det A=-2\)，有 \(\varepsilon'=-\tfrac12\varepsilon\)。标准基虽不必是正交单位基，但正向体积形式在它上的值必须为正且绝对值等于 \(\sqrt3\)，所以
	\[
		\omega=\sqrt3\,\varepsilon=-2\sqrt3\,\varepsilon',\qquad
		\mu_g=\sqrt3\,|\varepsilon|=2\sqrt3\,|\varepsilon'|.
	\]
	新基为负向基，故体积形式的系数为负，而体积密度的系数仍正。反转定向会把 \(\omega\) 变为 \(-\omega\)，不改变 \(g\)、\(\mu_g\)、两组固定基或它们的变换矩阵。体积密度因此不能单独判断所选的是哪一个定向。
\end{solution}

\begin{exercise}
	\label{g1:exr:1202-polar-rotation}
	在平面一个 \(r>0\) 的极坐标扇区内，固定一点。令 \(g\) 为 Euclid 内积，\(J\) 为正向旋转 \(\pi/2\) 的线性映射，即 \(J\partial_x=\partial_y\)、\(J\partial_y=-\partial_x\)。求 \(J\) 在基 \((\partial_r,\partial_\theta)\) 中的矩阵与 \((1,1)\) 型张量表达。由这些表达验证 \(J^2=-I\)、\(\operatorname{tr}J=0\) 和 \(g(Ju,Jv)=g(u,v)\)。再求双线性形式 \(\beta(u,v)=g(Ju,v)\)，说明它为何是二次交替形式，并写成 \(\mathrm{d}r,\mathrm{d}\theta\) 的楔积。
\end{exercise}
\begin{solution}
	直角坐标中的矩阵为 \(J_0=\begin{psmallmatrix}0&-1\\1&0\end{psmallmatrix}\)。极坐标换基矩阵可写成
	\[
		A=\begin{pmatrix}\cos\theta&-\sin\theta\\\sin\theta&\cos\theta\end{pmatrix}
		\begin{pmatrix}1&0\\0&r\end{pmatrix}.
	\]
	旋转矩阵与 \(J_0\) 交换，所以
	\[
		J'=A^{-1}J_0A=\begin{pmatrix}0&-r\\1/r&0\end{pmatrix}.
	\]
	等价地，\(J\partial_r=r^{-1}\partial_\theta\)、\(J\partial_\theta=-r\partial_r\)，故
	\[
		J=\frac1r\,\partial_\theta\otimes\mathrm{d}r
		-r\,\partial_r\otimes\mathrm{d}\theta.
	\]
	矩阵平方为 \(-I\)，对角元和为零。极坐标内积矩阵是 \(G'=\operatorname{diag}(1,r^2)\)，相乘得 \((J')^{\mathsf T}G'J'=G'\)，从而验证内积保持不变。也可令 \(u=a\partial_r+b\partial_\theta\)、\(v=c\partial_r+d\partial_\theta\)，直接算出
	\[
		g(Ju,Jv)=(-rb)(-rd)+r^2(a/r)(c/r)=r^2bd+ac=g(u,v).
	\]

	同样，
	\[
		\beta(u,v)=g(-rb\partial_r+(a/r)\partial_\theta,\,
		c\partial_r+d\partial_\theta)=r(ad-bc).
	\]
	交换 \(u,v\) 时此值变号，重复输入时为零，因此它是交替二次形式，并且
	\[
		\beta=r\,\mathrm{d}r\wedge\mathrm{d}\theta
		=\mathrm{d}x\wedge\mathrm{d}y.
	\]
	这里 \(J\) 的矩阵含有 \(r\) 与 \(1/r\)，却仍表示固定的直角旋转；内积的矩阵同时发生相应变化，配对之后得到的有向面积形式与直角坐标中的对象一致。
\end{solution}

\begin{exercise}\label{g1:ex:smooth-partition-on-line}
	令 \(\eta\) 为\cref{g1:lem:smooth-bump-function}中的函数，\(b(t)=\eta(1-t^2)\)。对 \(k\in\mathbb Z\) 定义
	\[
		H(x)=\sum_{k\in\mathbb Z}b(x-k),\qquad
		\psi_k(x)=\frac{b(x-k)}{H(x)}.
	\]
	证明这些公式给出从属于开覆盖 \(O_k=(k-\tfrac32,k+\tfrac32)\) 的光滑单位分解，并证明
	\[
		H(x+1)=H(x),\qquad
		\psi_k(x)=\psi_0(x-k).
	\]
	求在整数点与半整数点处非零的权重，并证明对每个整数 \(j\ge1\)，有 \(\sum_k\psi_k^{(j)}(x)=0\)。
\end{exercise}
\begin{solution}
	函数 \(b\) 光滑，在 \((-1,1)\) 上为正，在外部为零，故 \(\operatorname{supp}b=[-1,1]\)。因此 \(\operatorname{supp}b(\,\cdot-k)=[k-1,k+1]\subset O_k\)。给定 \(x\)，邻域 \((x-\tfrac14,x+\tfrac14)\) 只能与指标落在有限区间 \([x-\tfrac54,x+\tfrac54]\) 内的这些闭区间相交，故支集族局部有限，\(H\) 光滑。每个实数与某个整数的距离不超过 \(1/2\)，相应项严格为正，所以 \(H>0\)。

	归一化后，\(\psi_k\) 光滑、非负、总和为 \(1\)，且其支集与对应的 \(b(x-k)\) 相同，全部从属条件成立。换指标 \(\ell=k-1\) 得
	\[
		H(x+1)=\sum_k b(x+1-k)=\sum_\ell b(x-\ell)=H(x).
	\]
	于是 \(H(x-k)=H(x)\)，给出 \(\psi_k(x)=b(x-k)/H(x-k)=\psi_0(x-k)\)。

	在整数 \(m\) 处，只有 \(b(m-m)=b(0)>0\)，相邻项的自变量为 \(\pm1\)，也为零，故 \(\psi_m(m)=1\)，其余项为零。在 \(m+\tfrac12\) 处，只有指标 \(m,m+1\) 的两项非零；由于 \(b\) 为偶函数，两项相等，权重各为 \(1/2\)。最后，在任意计算点的固定邻域内，恒等式 \(\sum_k\psi_k=1\) 都是有限和恒等式，逐项求 \(j\) 阶导数即得所求。若把覆盖误写为 \((k-1,k+1)\)，支集就不包含于对应开集，不能得到本书定义中的从属关系。
\end{solution}

\begin{exercise}\label{g1:ex:sections-on-closed-discrete-set}
	设 \(M\) 为有限维、Hausdorff 且第二可数的光滑流形，允许有边界，\(E\to M\) 为有限秩实光滑向量丛，\(D\subseteq M\) 是闭离散子集，即 \(D\) 闭，且每个 \(p\in D\) 有一个只与 \(D\) 交于 \(p\) 的开邻域。任意指定向量 \(v_p\in E_p\)。证明存在整体光滑截面 \(s\)，使 \(s(p)=v_p\) 对所有 \(p\in D\) 成立。进一步证明，可要求 \(s\) 紧支，当且仅当集合 \(\{p\in D:v_p\ne0\}\) 有限。
\end{exercise}
\begin{solution}
	在每个 \(p\in D\) 附近选择局部标架，并把邻域缩小到只含一个 \(D\) 中的点。用 \(v_p\) 在该点的分量作常系数线性组合，得到局部光滑截面 \(s_p\)，满足 \(s_p(p)=v_p\)。因为这个邻域与 \(D\) 只交于 \(p\)，它已延拓该邻域内的全部指定值。故 \(p\mapsto v_p\) 正是\cref{g1:def:smooth-section-closed-set}意义下闭集 \(D\) 上的光滑截面，\cref{g1:thm:section-extension}给出整体延拓。

	若整体截面 \(s\) 紧支，令 \(K=\operatorname{supp}s\)。由于 \(D\) 闭，\(D\cap K\) 紧；它又是离散空间，而紧离散空间只能有限，因为由所有单点组成的开覆盖必须有有限子覆盖。因此非零数据的指标集包含于有限集 \(D\cap K\)，也有限。

	反之，设只有 \(p_1,\ldots,p_N\) 的指定值非零。为每个 \(p_j\) 取前面的局部截面 \(s_j\)，其定义域 \(U_j\) 只与 \(D\) 交于 \(p_j\)。对紧闭集 \(\{p_j\}\subset U_j\) 选紧支光滑截断函数 \(\chi_j\)，在 \(p_j\) 附近为 \(1\)，支集包含于 \(U_j\)。将 \(\chi_j s_j\) 补零到 \(M\)，然后取有限和
	\[
		s=\sum_{j=1}^N\chi_j s_j.
	\]
	该截面在 \(p_j\) 取值 \(v_{p_j}\)，在 \(D\) 的其余点为零，支集包含于有限个紧集之并，所以紧。没有非零数据时取零截面即可。
\end{solution}

\begin{exercise}\label{g1:ex:relative-lift-of-section}
	设 \(M\) 为有限维、Hausdorff 且第二可数的光滑流形，允许有边界，\(q:E\to F\) 是有限秩实光滑向量丛之间覆盖 \(\operatorname{id}_M\) 的逐纤维满射同态。证明每个 \(t\in\Gamma(F)\) 都有整体光滑提升 \(s\in\Gamma(E)\)，即 \(q\circ s=t\)。若 \(A\subseteq M\) 闭，且已在 \(A\) 的开邻域 \(U\) 上给定提升 \(s_U\)，证明可选整体提升与 \(s_U\) 在 \(A\) 附近相同。说明两条整体提升之差属于哪个截面空间。
\end{exercise}
\begin{solution}
	由\cref{g1:thm:constant-rank-bundle-morphism}，\(K=\ker q\) 为光滑子丛，且
	\[
		0\longrightarrow K\longrightarrow E\xrightarrow{\ q\ }F\longrightarrow0
	\]
	为短正合丛列。\cref{g1:thm:exact-bundle-splitting}给出光滑右逆 \(j:F\to E\)。于是 \(s_0=j\circ t\) 光滑，且 \(q\circ s_0=t\)。

	在 \(U\) 上，\(d=s_U-s_0|_U\) 满足 \(q(d)=0\)，故逐点落在 \(K\) 中，并且是 \(K\) 的光滑截面。最后这一点可在适配 \(K\) 的局部标架中看出：后面的分量恒为零，前面的分量仍光滑。用\cref{g1:thm:section-extension}把 \(d\) 延拓为 \(K\) 的整体截面 \(\widetilde d\)，保持 \(A\) 附近的值，令
	\[
		s=s_0+\widetilde d.
	\]
	则 \(q(s)=q(s_0)=t\)，且在 \(A\) 附近 \(s=s_U\)。若 \(A\) 紧，还可使改变量 \(\widetilde d\) 紧支。

	若 \(s_1,s_2\) 都提升 \(t\)，则 \(q(s_1-s_2)=0\)，其差属于 \(\Gamma(K)\)；反之，任意给一个提升加上 \(\Gamma(K)\) 中的截面，仍为提升。因此一旦选定一个提升，所有提升恰好由它加上核子丛的任意整体截面得到。
\end{solution}

\begin{exercise}\label{g1:ex:metric-dependent-splitting}
	设 \(M\) 为光滑流形，\(a\in C^\infty(M)\)。在平凡丛 \(E=M\times\mathbb R^2\) 上定义
	\[
		h_{a,p}\bigl((u,v),(u',v')\bigr)
		=(u+a(p)v)(u'+a(p)v')+vv'.
	\]
	证明 \(h_a\) 是丛度量，求第一坐标子丛 \(L=M\times(\mathbb R\times\{0\})\) 的正交补、到 \(L\) 的正交投影，以及由该正交补给出的投影 \(q:E\to M\times\mathbb R\)、\(q(p,u,v)=(p,v)\) 的右逆。进一步证明，\(q\) 的每个光滑丛右逆都可由某个这样的 \(h_a\) 得到。
\end{exercise}
\begin{solution}
	定义显然给出对称双线性型，其矩阵
	\[
		H_a=\begin{pmatrix}1&a\\a&1+a^2\end{pmatrix}
	\]
	光滑。对非零 \((u,v)\)，平方和 \((u+av)^2+v^2\) 严格为正，所以 \(h_a\) 为丛度量。

	向量 \((u,v)\) 与 \((1,0)\) 正交，当且仅当 \(u+av=0\)。因此 \(L^\perp\) 有整体标架 \((-a,1)\)。分解
	\[
		(u,v)=(u+av,0)+v(-a,1)
	\]
	给出正交投影 \(P_a(p,u,v)=(p,u+a(p)v,0)\)。由于 \(q(p,-a(p)t,t)=(p,t)\)，对应的丛右逆为
	\[
		j_a(p,t)=(p,-a(p)t,t).
	\]

	任意逐纤维线性的右逆 \(j\) 都由其在 \(t=1\) 处的值决定；右逆条件使第二坐标必须为 \(1\)，所以 \(j(p,t)=(p,c(p)t,t)\)。丛同态的光滑性等价于 \(c\) 光滑。取 \(a=-c\)，便有 \(j=j_a\)。可见同一条短正合丛列的不同分裂都能在本例中通过改变度量实现；丛列本身并未规定其中哪一个。
\end{solution}

\begin{exercise}
	\label{g1:ex:tensor-field-changing-frame}
	在 \(\mathbb R^2\) 上取 \(X=\partial_x+y\partial_y\)、\(\alpha=\mathrm{d}x+\mathrm{d}y\)，令 \(L=X\otimes\alpha\)，把它视为同态场 \(L(w)=\alpha(w)X\)。换用全局坐标 \(u=x\)、\(v=e^xy\)，求 \(X,\alpha,L\) 的新分量，并在两套坐标中核对 \(\operatorname{tr}L=\alpha(X)\)。
\end{exercise}
\begin{solution}
	坐标变化的逆为 \(x=u\)、\(y=ve^{-u}\)，两个方向都光滑。由 \(X(u)=1\)、\(X(v)=e^xy+ye^x=2v\)，得
	\[
		X=\partial_u+2v\partial_v.
	\]
	又由 \(\mathrm{d}x=\mathrm{d}u\)、\(\mathrm{d}y=e^{-u}\mathrm{d}v-ve^{-u}\mathrm{d}u\)，得
	\[
		\alpha=(1-ve^{-u})\mathrm{d}u+e^{-u}\mathrm{d}v.
	\]
	向量与协变向量按不同方向变换，合起来给出同态的矩阵。在旧、新坐标基中分别为
	\[
		[L]_{x,y}=\begin{pmatrix}1&1\\y&y\end{pmatrix},\qquad
		[L]_{u,v}=\begin{pmatrix}1-ve^{-u}&e^{-u}\\2v(1-ve^{-u})&2ve^{-u}\end{pmatrix}.
	\]
	所以旧坐标中迹为 \(1+y\)，新坐标中为 \(1+ve^{-u}\)，正是同一函数。直接求值也有 \(\alpha(X)=1+y=(1-ve^{-u})+2ve^{-u}\)。矩阵分量随标架改变，收缩时换基因子相消，因而迹保持不变。
\end{solution}

\begin{exercise}
	\label{g1:ex:exponential-polar-pullbacks}
	设 \(F:\mathbb R^2\to\mathbb R^2\setminus\{0\}\)，\(F(u,v)=(e^u\cos v,e^u\sin v)\)。在目标上取 Euclid 度量 \(g=\mathrm{d}x\otimes\mathrm{d}x+\mathrm{d}y\otimes\mathrm{d}y\)、一形式 \(\alpha=x\,\mathrm{d}y-y\,\mathrm{d}x\)、二形式 \(\omega=\mathrm{d}x\wedge\mathrm{d}y\)。求三者的拉回，并求 \(F^*g\) 下函数 \(u\) 的梯度。说明即使 \(F\) 不是整体单射，这些构造仍然有意义。
\end{exercise}
\begin{solution}
	有
	\[
		F^*(\mathrm{d}x)=e^u\cos v\,\mathrm{d}u-e^u\sin v\,\mathrm{d}v,
		\qquad
		F^*(\mathrm{d}y)=e^u\sin v\,\mathrm{d}u+e^u\cos v\,\mathrm{d}v.
	\]
	平方的对称交叉项相消，故
	\[
		F^*g=e^{2u}(\mathrm{d}u\otimes\mathrm{d}u+\mathrm{d}v\otimes\mathrm{d}v).
	\]
	把函数 \(x,y\) 同时与 \(F\) 复合，再代入上式中的一形式，得到 \(F^*\alpha=e^{2u}\mathrm{d}v\)。按楔积展开，则 \(F^*\omega=e^{2u}\mathrm{d}u\wedge\mathrm{d}v\)。该度量的逆矩阵为 \(e^{-2u}I\)，所以 \(\operatorname{grad}_{F^*g}u=e^{-2u}\partial_u\)。

	\(F(u,v+2\pi)=F(u,v)\)，因此 \(F\) 不单射；但 \(\det \mathrm{D}F=e^{2u}>0\)，它在每点附近为局部微分同胚。协变张量的拉回本来就不要求单射，拉回度量的正定性只要求微分单射，故本题没有矛盾。若在目标取 \(h(x,y)=\log\sqrt{x^2+y^2}\)，则 \(h\circ F=u\)，并有
	\[
		\mathrm{d}F(\operatorname{grad}_{F^*g}u)
		=e^{-2u}(x,y)=\operatorname{grad}_g h\big|_{F(u,v)},
	\]
	也可直接核对这个局部等距识别下的梯度关系。
\end{solution}

\begin{exercise}
	\label{g1:ex:oriented-plane-complex-structure}
	设 \(E\to M\) 是可定向的秩二实光滑向量丛，固定定向和丛度量 \(g\)。在每张正向正交单位标架 \((e_1,e_2)\) 中规定 \(Je_1=e_2\)、\(Je_2=-e_1\)。证明这定义整体光滑 \(J\)，且 \(J^2=-I\)。若 \(\nu\) 是由定向和 \(g\) 决定的单位体积形式，证明 \(\nu(v,w)=g(Jv,w)\)。最后将度量改为 \(\widetilde g=e^{2f}g\)，其中 \(f\) 光滑，比较两者对应的 \(J\) 与体积形式。
\end{exercise}
\begin{solution}
	两张正向正交单位标架之间的变换矩阵在 \(\operatorname{SO}(2)\) 中。任一这样的矩阵都可写为 \(\begin{psmallmatrix}a&-b\\b&a\end{psmallmatrix}\)，其中 \(a^2+b^2=1\)，它与 \(J_0=\begin{psmallmatrix}0&-1\\1&0\end{psmallmatrix}\) 交换。因此各图中规定的 \(J\) 在交叠处一致，且局部矩阵恒为 \(J_0\)，所以 \(J\) 整体光滑并满足 \(J^2=-I\)。

	对 \(v=ae_1+be_2\)、\(w=ce_1+de_2\)，有 \(g(Jv,w)=ad-bc\)。同一正向单位标架中，\(\nu=e^1\wedge e^2\)，故 \(\nu(v,w)=ad-bc\)，得到所需恒等式。这也说明 \(J\) 由度量与定向唯一确定：对每个 \(v\)，非退化性使满足 \(g(Jv,\cdot)=\nu(v,\cdot)\) 的向量唯一。

	对新度量，\((e^{-f}e_1,e^{-f}e_2)\) 为正向正交单位标架。在它上面规定的四分之一转动仍是原来的 \(J\)，因为两个基向量作了相同的正缩放。对应的对偶标架各乘 \(e^f\)，所以 \(\widetilde\nu=e^{2f}\nu\)。因此正向共形缩放不改变 \(J\)，但改变面积的尺度。
\end{solution}

\begin{exercise}
	\label{g1:ex:prescribed-positive-density}
	设 \(M\) 为 \(n\geq1\) 维 Hausdorff、第二可数光滑流形，可带边界，\(\mu\) 是一个光滑正密度。
	\begin{enumerate}
		\item 证明存在 Riemann 度量 \(g\) 的体积密度恰为 \(\mu\)，不要求 \(M\) 可定向。
		\item 若两个度量 \(g_0,g_1\) 的体积密度均为 \(\mu\)，构造连接它们且始终保持体积密度 \(\mu\) 的光滑度量族。说明普通凸组合为何未必保持体积密度。
	\end{enumerate}
\end{exercise}
\begin{solution}
	1. 先由\cref{g1:thm:bundle-metric-existence}给 \(\mathrm{T}M\) 取一个度量 \(h\)，记其正体积密度为 \(\mu_h\)。两个正密度的比 \(\rho=\mu/\mu_h\) 是定义在 \(M\) 上的正光滑函数：在任一局部标架中是两个正光滑系数的比，而二者换标架时都乘以同一绝对行列式，所以比值与标架无关。若令 \(g=\rho^{2/n}h\)，其局部矩阵行列式为 \(\rho^2\det H\)，故 \(\mu_g=\rho\mu_h=\mu\)。构造只使用密度的正性与比值，不需要选择定向。

	2. 先令 \(k_t=(1-t)g_0+tg_1\)、\(0\leq t\leq1\)。它对每个非零向量均取正值，故是一族光滑度量，但体积密度可能变化。再按第一问逐个归一化：
	\[
		g_t=\left(\frac{\mu}{\mu_{k_t}}\right)^{2/n}k_t.
	\]
	\(\mu_{k_t}\) 的系数是正行列式的平方根，联合光滑依赖于 \((t,p)\)；因此上式给出联合光滑度量族，且每个 \(g_t\) 的体积密度均为 \(\mu\)。端点处比值为一，故确实连接给定 \(g_0,g_1\)。

	普通凸组合不保体积的例子已可取在平凡二维丛中：\(G_0=\operatorname{diag}(2,1/2)\)、\(G_1=\operatorname{diag}(1/2,2)\) 的行列式都为一，但平均矩阵为 \((5/4)I\)，其体积密度系数为 \(5/4\)，不再是一。正定性对平均稳定，固定行列式却不是线性条件，因而需要额外归一化。
\end{solution}
